% ===========================================================================
%  Chapter 1 -- Introduction
%  Disease-Adaptive Heterogeneous Graph Learning with Anatomically Aligned
%  Multi-Sequence MRI Representations
%
%  PROSPECTIVE THESIS-FRAMING VERSION 2.0
%  Written before completion of the full experimental programme.
%  The research questions, hypotheses, scope boundaries and novelty position are
%  intended to remain stable. Fields marked [TO CONFIRM], [TO RECORD] or
%  [LATER INTEGRATION] must be completed transparently when the relevant facts or
%  results exist; they must not be used to rewrite hypotheses after results are seen.
%
%  Compiles standalone:   pdflatex chapter1 ; biber chapter1 ; pdflatex chapter1
%
%  To fold into the thesis, delete everything from \documentclass down to
%  \begin{document}, and the \printbibliography/\end{document} at the foot, then
%  % ===========================================================================
%  Chapter 1 -- Introduction
%  Disease-Adaptive Heterogeneous Graph Learning with Anatomically Aligned
%  Multi-Sequence MRI Representations
%
%  PROSPECTIVE THESIS-FRAMING VERSION 2.0
%  Written before completion of the full experimental programme.
%  The research questions, hypotheses, scope boundaries and novelty position are
%  intended to remain stable. Fields marked [TO CONFIRM], [TO RECORD] or
%  [LATER INTEGRATION] must be completed transparently when the relevant facts or
%  results exist; they must not be used to rewrite hypotheses after results are seen.
%
%  Compiles standalone:   pdflatex chapter1 ; biber chapter1 ; pdflatex chapter1
%
%  To fold into the thesis, delete everything from \documentclass down to
%  \begin{document}, and the \printbibliography/\end{document} at the foot, then
%  % ===========================================================================
%  Chapter 1 -- Introduction
%  Disease-Adaptive Heterogeneous Graph Learning with Anatomically Aligned
%  Multi-Sequence MRI Representations
%
%  PROSPECTIVE THESIS-FRAMING VERSION 2.0
%  Written before completion of the full experimental programme.
%  The research questions, hypotheses, scope boundaries and novelty position are
%  intended to remain stable. Fields marked [TO CONFIRM], [TO RECORD] or
%  [LATER INTEGRATION] must be completed transparently when the relevant facts or
%  results exist; they must not be used to rewrite hypotheses after results are seen.
%
%  Compiles standalone:   pdflatex chapter1 ; biber chapter1 ; pdflatex chapter1
%
%  To fold into the thesis, delete everything from \documentclass down to
%  \begin{document}, and the \printbibliography/\end{document} at the foot, then
%  % ===========================================================================
%  Chapter 1 -- Introduction
%  Disease-Adaptive Heterogeneous Graph Learning with Anatomically Aligned
%  Multi-Sequence MRI Representations
%
%  PROSPECTIVE THESIS-FRAMING VERSION 2.0
%  Written before completion of the full experimental programme.
%  The research questions, hypotheses, scope boundaries and novelty position are
%  intended to remain stable. Fields marked [TO CONFIRM], [TO RECORD] or
%  [LATER INTEGRATION] must be completed transparently when the relevant facts or
%  results exist; they must not be used to rewrite hypotheses after results are seen.
%
%  Compiles standalone:   pdflatex chapter1 ; biber chapter1 ; pdflatex chapter1
%
%  To fold into the thesis, delete everything from \documentclass down to
%  \begin{document}, and the \printbibliography/\end{document} at the foot, then
%  \input{chapter1} from the main thesis file.
% ===========================================================================

\documentclass[12pt,a4paper]{report}

\usepackage{lmodern}
\usepackage[T1]{fontenc}
\usepackage[utf8]{inputenc}
\usepackage[english]{babel}
\usepackage{csquotes}
\usepackage{microtype}
\usepackage[margin=2.5cm]{geometry}
\usepackage{setspace}
\onehalfspacing

\usepackage{amsmath,amssymb}
\usepackage{booktabs}
\usepackage{longtable}
\usepackage{array}
\usepackage{graphicx}
\usepackage{thesisstyle}   % shared TikZ figure styles; see thesisstyle.sty
\usepackage[hidelinks]{hyperref}

\usepackage[
  backend=biber,
  style=numeric-comp,
  sorting=none,
  maxbibnames=6,
  minbibnames=3,
  giveninits=true,
  doi=true,
  url=false,
  isbn=false
]{biblatex}

% Production bibliography cleaned for thesis compilation. Citation keys are
% unchanged from the working literature inventory.
\addbibresource{lumbar_spine_mri_ai_references_clean.bib}

\newcommand{\toconfirm}[1]{\textbf{[TO CONFIRM: #1]}}
\newcommand{\torecord}[1]{\textbf{[TO RECORD AFTER IMPLEMENTATION: #1]}}
\newcommand{\laterinsert}[1]{\textbf{[LATER INTEGRATION: #1]}}

\begin{document}

\chapter{Introduction}
\label{ch:introduction}

% ===========================================================================
\section{Introduction and Thesis Position}
\label{sec:intro-position}

Lumbar degenerative disease presents an unusually demanding problem for medical
image analysis. The anatomical region is compact, but the diagnostic task is not:
one examination contains several intervertebral levels, several pathological
compartments, left--right distinctions, multiple magnetic resonance imaging (MRI)
sequences acquired in different planes, and severity labels whose boundaries are
partly subjective. A clinically meaningful automated system must therefore do more
than recognise that degeneration is present. It must identify \emph{what} is
abnormal, \emph{where} it is abnormal, \emph{how severe} the abnormality is, and
\emph{which imaging evidence} supports that judgement. It must also remain
interpretable as a measurement of radiological grading rather than be mistaken for
an autonomous clinical diagnosis.

The need for such systems arises from both disease burden and reporting workload.
Low back pain remains one of the largest contributors to disability worldwide, and
lumbar degenerative change is a frequent structural correlate examined by MRI
\cite{compte2023are,wang2024machine}. Routine interpretation is repetitive because
the reader must inspect multiple levels across sagittal and axial acquisitions,
while the severity of central canal, foraminal and subarticular narrowing is judged
against morphological criteria that are not perfectly reproducible between human
readers \cite{lurie2008reliability,schizas2010qualitative,lee2011a,lee2010a}.
Automation is attractive in this setting not because a neural network is assumed to
possess a superior biological truth, but because a reproducible system may help
standardise a high-volume structured task whose human reference process itself has
known variability.

The technical field has advanced rapidly. Early systems demonstrated that lumbar
anatomy could be localised and individual discs graded with convolutional networks;
more recent pipelines combine explicit localisation, 2.5D regions of interest,
multiview fusion, attention and large public benchmarks
\cite{lu2018deepspine,pan2021automatically,batra2025mscan,richards2026the}. The
central problem has consequently shifted. The thesis no longer asks whether a
convolutional neural network or transformer can classify a lumbar MRI in isolation.
It asks whether the \emph{structure of the grading problem itself} has been
represented faithfully enough.

This thesis takes the position that lumbar MRI grading is a structured prediction
problem with four coupled sources of information. First, disease occurs at ordered
anatomical levels rather than at unrelated image locations. Second, several
conditions coexist within one motion segment and lateral compartments occur as
paired left--right structures. Third, the diagnostic value of an MRI sequence is
target-dependent: sagittal T1, sagittal T2/STIR and axial T2 provide complementary
rather than interchangeable evidence \cite{hallinan2021deep}. Fourth, a model that
performs well on one benchmark may degrade when scanners, protocols, population and
reference-standard conventions change across institutions
\cite{mcsweeney2023external,zhang2023deep,guan2022domain}.

The research therefore evaluates three methodological contributions and one
translational validation programme. The integrated framework evaluating these components
is designated \textbf{AMOG-Net} (\textbf{A}natomy-aware \textbf{M}ulti-sequence \textbf{O}rdinal \textbf{G}raph Network). The first contribution is anatomically aligned cross-sequence
self-supervision, in which DICOM patient-space correspondence defines which views
should agree during representation learning. The second is disease-conditioned
adaptive sequence routing, in which the model learns target- and patient-dependent
weights over available MRI sequences and is explicitly tested under both missing
and degraded inputs. The third is a heterogeneous disease--anatomy graph whose
nodes represent level--condition--laterality targets and whose typed edges encode
adjacent-level, same-level cross-condition and bilateral relations. These are then
evaluated under a deliberately separate transfer programme: development occurs on
a multinational public benchmark; the selected model is frozen; zero-shot
performance is measured on an independent Rizgary Hospital cohort; and controlled
few-shot adaptation is performed only afterwards.

Several related techniques are necessary but are not claimed as doctoral novelty.
Anatomical localisation, segmentation, 2.5D region construction, conventional CNN
or transformer encoders, class-imbalance handling, ordinal loss functions,
calibration and parameter-efficient fine-tuning are established tools. They are
used to construct a fair experiment around the thesis contributions rather than
presented as original merely because they appear in the same architecture. This
novelty boundary is important in a field where anatomy-guided contextual grading,
cross-sequence attention and graph-based disc modelling have already been reported
\cite{chai2026anatomy,nguyen2026crossspine,baur2024automated}. The question is not
whether any one of those mechanisms exists elsewhere; it is whether the particular
representation tested here---typed relationships among the benchmark's disease
outputs, anatomical cross-sequence supervision, target-conditioned evidence
allocation and explicit institutional transfer---adds measurable knowledge beyond
strong matched baselines.

% Written after the confirmatory experiments were locked. Reports the outcome
% only; the problem statement, research questions and hypotheses are unchanged.
The principal empirical outcome is that the complete system improves agreement
with the reference standard over a single-sequence baseline by $+0.0177$
quadratic weighted kappa, on every one of seven training seeds and after
correction for multiple comparisons, but that the anatomical mechanisms proposed
to achieve this do not account for the improvement: attribution analysis
indicates the convolutional encoder already localises the graded structure
without them.

% ===========================================================================
\section{Background and Motivation}
\label{sec:intro-background}

\subsection{Clinical Motivation: A Repetitive but Structured Imaging Task}
\label{sec:intro-clinical-motivation}

Lumbar degeneration is not a single lesion. It is a family of related structural
changes involving the intervertebral disc, facet joints, ligamentum flavum,
vertebral endplates and the neural spaces that these structures surround. Disc
dehydration and height loss alter load distribution; annular bulging or herniation
may encroach on the canal and foramina; posterior element hypertrophy may further
reduce the space available to neural tissue. These changes are assessed at several
levels, most commonly from L1--L2 through L5--S1, and their significance depends on
which anatomical compartment is narrowed.

Three stenotic compartments are especially relevant to the benchmark task used in
this thesis. Central canal stenosis concerns the thecal sac and cauda equina,
foraminal narrowing concerns the exiting nerve roots, and subarticular or lateral
recess stenosis concerns the traversing roots. Their visual criteria are related but
not identical, and neither are the acquisitions from which they are best judged.
Central canal and lateral recess morphology depend heavily on axial T2 information,
whereas neural foraminal assessment benefits particularly from sagittal or
parasagittal T1 because perineural fat provides the relevant contrast
\cite{hallinan2021deep,schizas2010qualitative,lee2010a}. A model that treats all
sequences as equivalent therefore ignores a clinical property of the task before
learning has even begun.

The labels themselves are another source of complexity. Morphological grading
systems convert continuous anatomical variation into ordinal categories. Human
agreement is strongest for some compartments and weaker for others: the reliability
study of Lurie et al. found substantially better inter-reader agreement for central
canal stenosis than for foraminal and subarticular findings
\cite{lurie2008reliability}. This matters methodologically. If the reference standard
contains boundary uncertainty, model evaluation must distinguish gross errors from
one-grade disagreements and must interpret performance relative to the reliability
of the grading instrument. Apparent model error and reference-standard disagreement
are not always separable.

This combination---high examination volume, repeated level-wise assessment,
complementary sequence evidence and imperfect categorical boundaries---makes lumbar
MRI a suitable domain for artificial intelligence, but only if the automation is
framed correctly. The objective is not to replace full clinical interpretation.
Symptoms, neurological findings, prior imaging and management context remain outside
the information available to an image-only grading system. The narrower objective is
to make a defined radiological grading task more consistent, measurable and
portable. Reader-assistance studies support this framing: deep-learning assistance
has been shown to reduce reporting time and improve or preserve inter-reader
agreement, with the greatest gains often occurring among less experienced readers
\cite{lim2022improved}. This thesis therefore treats the system as a potential
component of decision support, not as an autonomous treatment decision-maker.

\subsection{Technical Motivation: From Classification to Structured Prediction}
\label{sec:intro-technical-motivation}

The evolution of lumbar MRI automation has progressively narrowed the unit of
prediction. Early work often asked whether a lesion or abnormality was present in an
image. Subsequent systems localised discs, graded individual levels, and shared
encoders across several pathologies. Public benchmark design has now made the task
more explicit. The RSNA Lumbar Degenerative Imaging Spine Classification resource,
published as LumbarDISC, contains 2,697 patients and 8,593 MRI series collected
across eight institutions in six countries \cite{richards2026the}. The associated
challenge requires five condition/laterality targets at five lumbar levels, producing
25 ordinal outputs per study \cite{rsna2024rsna,richards2026the}. A correct model
therefore does not simply output ``mild'', ``moderate'' or ``severe'' for an image;
it attaches one of those grades to a particular condition at a particular level and,
for lateral compartments, a particular side.

This target structure exposes a limitation in the prevailing formulation. Most
systems share image features but ultimately emit independent output heads. Shared
features allow several tasks to benefit from a common representation, but a
prediction at L4--L5 does not explicitly inform an uncertain prediction at L3--L4;
a confident left foraminal prediction does not interact structurally with its right
counterpart; and the central, foraminal and subarticular targets at one motion
segment do not exchange typed messages simply because they coexist anatomically.
The representation is therefore multi-task but not fully relational.

Graph learning provides a natural language for testing whether those relationships
carry useful information. Graph neural networks are not novel by themselves, and a
lumbar application has already reconstructed an individual disc as a point cloud and
graded it with a graph network \cite{baur2024automated}. That precedent is important
because it prevents an unjustified claim that ``using a graph in lumbar MRI'' is new.
The unresolved question is narrower: whether the \emph{prediction targets} can be
represented as nodes connected by clinically motivated relation types, and whether
those relations improve grading beyond independent heads, an ordered level
transformer and a homogeneous graph. The distinction between a graph over points
within one disc and a graph over disease--anatomy targets across the spine is the
novelty boundary this thesis tests.

A second technical motivation concerns multimodal evidence. Recent methods have
already shown that cross-view and cross-sequence attention can improve lumbar
assessment \cite{batra2025mscan,nguyen2026crossspine}, while anatomy-guided
inter-level context has also become an active direction \cite{chai2026anatomy}.
These developments mean that novelty cannot rest on ``using multiple sequences'' or
``using attention''. The remaining question is whether evidence allocation should be
\emph{conditioned on the target and the patient}, and whether one trained model can
remain operational when an expected sequence is unavailable or present but of poor
quality. A fixed fusion rule assumes the same contribution of sagittal T1, sagittal
T2/STIR and axial T2 for every target and every study. Clinical sequence--pathology
correspondence makes that assumption doubtful.

A third motivation is representation learning. Generic ImageNet pretraining is
convenient but does not exploit the particular correspondence available in MRI: the
same anatomical level can be observed under different contrasts and planes. Earlier
spinal work demonstrated that self-supervision derived from patient-level or
longitudinal information can reduce dependence on manual labels
\cite{jamaludin2017selfsupervised}, and disc-centric contrastive work continues that
line \cite{acharya2026disccentric}. DICOM geometry makes a stronger, directly
anatomical pairing possible. If sagittal and axial observations can be matched in
patient space, then the correspondence itself can serve as supervision without
pretending that the images should look visually alike.

Finally, there is a translational motivation. Systematic reviews repeatedly find
that internally impressive medical-imaging models underperform when re-evaluated
outside their development distribution \cite{compte2023are,wang2024machine}.
Lumbar-specific external validations show the same pattern, with fine-grained
grading generally less portable than localisation or detection
\cite{mcsweeney2023external,zhang2023deep}. A model intended for eventual use in
Iraq or the wider region therefore cannot be justified solely by performance on a
multinational benchmark. The model must be transported to a genuinely independent
hospital cohort, its zero-shot degradation measured before local tuning, and the
amount of local annotation required to recover performance quantified explicitly.
That annotation--performance curve is more actionable than a single statement that
``domain adaptation improved accuracy''.

\subsection{Why the Problem Should Be Addressed Now}
\label{sec:intro-why-now}

Three developments make the present study timely. First, LumbarDISC supplies a
large, multi-institutional and anatomically structured benchmark whose target schema
is rich enough to test relational modelling rather than generic image
classification \cite{richards2026the}. Second, public anatomical resources such as
SPIDER reduce the need to claim localisation or segmentation as the central
contribution \cite{vandergraaf2024lumbar}. Third, the current methodological
literature is sufficiently mature to provide strong controls: multi-view
cross-attention, anatomy-guided contextual models, contemporary CNN and transformer
backbones, self-supervised pretraining, calibration and parameter-efficient transfer
all exist as comparators or supporting tools. The research can therefore ask a more
precise question than whether a new architecture attains a high score.

The resulting opportunity is methodological rather than merely computational. The
field now has enough data and enough baseline competence to test whether explicit
clinical structure provides value over capacity alone. If a typed target graph,
anatomical self-supervision or disease-conditioned router fails to outperform its
control, that negative result is still informative. Conversely, if a gain disappears
when topology is shuffled, when a simpler transformer is matched for capacity, or
when the model is transported to another hospital, the interpretation must change.
This insistence on falsifiable components is part of the motivation for the study,
not an afterthought added to evaluation.

% ===========================================================================
\section{Problem Statement}
\label{sec:intro-problem}

Despite rapid progress in automated lumbar MRI analysis, four connected
methodological problems remain.

\subsection{Problem 1: The Output Space Is Structured but Commonly Treated as Independent}
\label{sec:intro-problem-relational}

The benchmark output space contains ordered lumbar levels, multiple conditions at
each level and bilateral counterparts. These targets arise from one connected
anatomical system, yet the dominant design predicts them independently after a
shared encoder. This separation is convenient for implementation but does not test
whether adjacent-level context, within-level co-occurrence or bilateral relations
carry predictive information. Existing graph work in lumbar MRI does not resolve
this question because its graph is constructed within a disc rather than across the
disease targets \cite{baur2024automated}. CrossSpine introduces anatomical priors
and level-specific behaviour but still does not formulate the full 25 target outputs
as a typed communication graph \cite{nguyen2026crossspine}. The first problem is
therefore not the absence of contextual modelling in general; it is the absence of a
controlled test of explicit disease--anatomy relations among the prediction targets.

A relational model also carries a specific risk: useful context can become
``information contagion''. A severe lesion at one level should not cause the model
to elevate an adjacent normal level simply because degeneration often co-occurs.
Any graph proposed as a solution must therefore preserve local visual evidence,
include non-relational and homogeneous controls, and be stress-tested on naturally
isolated lesions. Without those controls, improved aggregate performance would not
establish that the anatomical topology is beneficial.

\subsection{Problem 2: Sequence Fusion Is Usually Fixed While Evidence Is Target- and Quality-Dependent}
\label{sec:intro-problem-routing}

Sagittal T1, sagittal T2/STIR and axial T2 are complementary acquisitions. The
foramina, central canal and lateral recess are not equally well represented in each
view, and routine data may be incomplete or technically degraded. Strong recent
systems address multi-view fusion through cross-attention
\cite{batra2025mscan,nguyen2026crossspine}, but a remaining deployment question is
whether one model can dynamically allocate evidence by condition, level, patient and
sequence quality. A model trained only on complete, high-quality benchmark studies
may fail when one sequence is missing, truncated or motion-corrupted. Missingness
and corruption are also different problems: an absent sequence can be masked,
whereas a degraded sequence remains present and can attract misleading attention.
The second problem is therefore to learn adaptive sequence allocation without
turning the router into an unexamined source of instability.

\subsection{Problem 3: Cross-Sequence Correspondence Is Available but Underused as Supervision}
\label{sec:intro-problem-ssl}

Self-supervised learning can reduce dependence on costly expert labels, but common
contrastive objectives construct positive pairs through augmentations of the same
image. Multi-sequence MRI contains a more domain-specific relationship: different
acquisitions can depict the same anatomical level even when their intensity pattern,
orientation and diagnostic role differ. DICOM patient-space geometry provides the
physical correspondence required to identify those pairs. Prior spinal
self-supervision demonstrates the value of metadata-derived or longitudinal
relationships \cite{jamaludin2017selfsupervised}, but anatomical cross-sequence
correspondence has not been established as the supervision signal for the target
formulation studied here.

The problem is not solved simply by forcing all sequence representations to become
identical. T1 and T2 contrasts contain sequence-specific diagnostic information, so
excessive invariance could improve embedding alignment while damaging grading. The
study must therefore separate the self-supervised projection space from downstream
features and test explicitly whether sequence-specific information survives
pretraining.

\subsection{Problem 4: Internal Accuracy Does Not Establish Transportability or Clinical Error Alignment}
\label{sec:intro-problem-transfer}

The strongest evidence against uncritical internal benchmarking is the repeated
performance decline under external validation
\cite{compte2023are,mcsweeney2023external,zhang2023deep}. Scanner manufacturer,
field strength, resolution, axial prescription, preprocessing and patient mix can
all change. More subtly, the reference standard may change: a public benchmark may
use structured expert annotation while a hospital's routine reports reflect local
reporting habits and may omit findings that were not considered clinically salient.
An external performance drop can therefore combine imaging-domain shift,
localisation failure and reference-standard shift.

A defensible transfer study must measure zero-shot performance before local tuning,
keep the local test set fixed, separate it from the adaptation pool, and harmonise
the local reference standard as far as resources allow. The present work therefore
uses the Rizgary Hospital cohort for external evaluation rather than for iterative
architecture development. The principal local comparison is restricted to central
canal stenosis, for which a defensible overlap with the public benchmark can be
created; local disc bulge, protrusion and extrusion remain a separate morphology
task rather than being relabelled as stenosis.

The error objective also requires care. Severity grades are ordinal, and distant
under-grading is clinically different from adjacent disagreement. Yet the literature
does not establish that ordinal losses automatically improve classification:
Niemeyer et al. found several ordinal formulations unable to outperform standard
cross-entropy for Pfirrmann grading \cite{niemeyer2021a}. A recent preprint has begun to evaluate distance-aware ordinal assessment on
LumbarDISC \cite{trinc2026beyond}; however, the combination of asymmetric clinical
cost with target-level relational modelling appears unaddressed in the literature
reviewed for this thesis.
At the same time, reducing
Severe$\rightarrow$Normal/Mild errors by simply shifting the model toward predicting
Severe would create an unacceptable false-positive burden. The fourth problem is
therefore to evaluate error direction, calibration and selective prediction without
converting radiological grading into an unsupported claim of treatment utility.

\subsection{Integrated Problem Statement}
\label{sec:intro-integrated-problem}

Taken together, the research problem can be stated as follows:

\begin{quote}
Current automated lumbar MRI grading systems can localise anatomy and predict
per-level disease severity with increasingly strong internal performance, but the
prevailing formulation does not fully exploit the relational structure of the
25-target output space, does not explicitly learn target- and quality-dependent
sequence allocation, underuses DICOM-defined cross-sequence anatomical
correspondence as a self-supervised signal, and provides limited evidence on how
these design choices behave under true institutional transfer. A controlled
methodology is therefore required to determine whether explicitly modelling these
properties improves grading, robustness, label efficiency and portability beyond
strong matched baselines, while preserving local evidence and respecting the
uncertainty of the clinical reference standard.
\end{quote}

\laterinsert{If implementation reveals a previously unrecognised but reproducible
problem that materially changes the framing---for example a systematic DICOM
geometry failure mode across a substantial cohort---add a short factual paragraph
here and document the protocol amendment. Do not add a new problem merely because a
proposed contribution underperformed.}

% ===========================================================================
\section{Research Gap and Novelty Boundary}
\label{sec:intro-gap}

Chapter~\ref{ch:litreview} develops the evidence in detail; the purpose here is to
state the gap concisely enough to define the thesis.

\subsection{Gap A: Explicit Relational Modelling of Disease--Anatomy Targets}
\label{sec:intro-gap-graph}

The reviewed literature contains multi-task systems, inter-level attention,
anatomical priors and one lumbar graph model, but no reviewed study constructs the
benchmark's level--condition--laterality targets as a heterogeneous graph with typed
relations for adjacent levels, within-level condition coexistence and bilateral
symmetry. Baur et al. use a graph within one reconstructed disc
\cite{baur2024automated}; Nguyen et al. introduce level-specific priors and
cross-sequence attention without target-to-target graph communication
\cite{nguyen2026crossspine}; Chai et al. demonstrate anatomy-guided contextual
multi-level modelling without the disease-target graph formulation proposed here
\cite{chai2026anatomy}. The intended contribution is therefore not ``graph neural
networks for lumbar MRI'' but a controlled test of an explicit typed graph over the
prediction targets themselves.

\subsection{Gap B: Anatomically Aligned Self-Supervision and Disease-Conditioned Evidence Routing}
\label{sec:intro-gap-sequence}

The literature establishes both the complementarity of MRI sequences and the value
of self-supervised pretraining, but these strands remain only partially connected in
lumbar grading. Anatomical correspondence across sequences is usually treated as a
preprocessing requirement for fusion rather than as supervision. At the same time,
multisequence systems generally learn a common fusion mechanism rather than a
condition- and quality-dependent allocation that is deliberately challenged under
missing and corrupted modalities. The thesis therefore tests two linked ideas: that
DICOM geometry can define cross-sequence positive relationships, and that a single
model can use the resulting features through a target-conditioned router whose
failure modes are directly measurable.

\subsection{Gap C: Quantified Transfer and Adaptation Under a Harmonised External Task}
\label{sec:intro-gap-transfer}

Existing external validations establish that transport is difficult, but they
usually report one transported model at one point in the adaptation process---often
with no adaptation at all. The operational question for a receiving hospital is
instead how performance changes before adaptation and how quickly it recovers as
local labels are supplied. The literature reviewed for this thesis contains no
Middle Eastern evaluation of the same scope and no quantified parameter-efficient
annotation curve for this lumbar target formulation. The gap is therefore not
simply ``lack of external validation''; it is lack of a staged zero-shot-to-few-shot
transfer experiment with fixed local testing, controlled annotation budgets and
explicit reference-standard harmonisation.

\subsection{Supporting Gap: Directional Error and Calibrated Uncertainty}
\label{sec:intro-gap-error}

Ordinal severity and asymmetric error are clinically meaningful, but ordinal
objectives have not consistently outperformed categorical baselines
\cite{niemeyer2021a}. Calibration is also rarely reported in lumbar MRI despite the
known overconfidence of modern neural networks \cite{guo2017calibration}. These are
therefore treated as supporting questions rather than headline novelty. A
cost-sensitive objective is retained only if it reduces clinically consequential
under-grading without creating indiscriminate Severe false positives, and selective
prediction is evaluated as a model behaviour rather than proof of a validated
clinical workflow.

The novelty claim throughout the thesis will remain appropriately bounded. The
phrasing ``to our knowledge'' or ``in the reviewed literature'' is used where
necessary, and the final claim will be updated against the literature available at
the time of submission. A new paper that adopts one element of the proposed system
does not invalidate the PhD if the controlled experimental question remains
distinct; it does, however, change how originality must be described.

\laterinsert{Immediately before thesis submission, update this novelty boundary
against publications appearing after the current literature-search cut-off. Record
which claims remain original and which should be reframed as independent validation,
extension or comparative evidence.}

% ===========================================================================
\section{Research Aim}
\label{sec:intro-aim}

The overall aim of this PhD is:

\begin{quote}
\emph{to determine whether explicitly modelling anatomical relationships among
lumbar disease targets, anatomically aligned evidence across MRI sequences, and
target-dependent sequence availability improves the accuracy, robustness, label
efficiency and cross-institutional portability of automated lumbar spine MRI
severity grading, relative to strong non-relational and fixed-fusion baselines.}
\end{quote}

The aim is intentionally phrased as a determination rather than a promise of
superiority. The thesis succeeds scientifically if it establishes under controlled
conditions where these mechanisms help, where they fail, and what their limits are.

% ===========================================================================
\section{Research Objectives}
\label{sec:intro-objectives}

The aim is operationalised through the following objectives.

\begin{enumerate}
  \item To construct a reproducible patient-level lumbar MRI analysis pipeline that
        preserves DICOM patient-space geometry, performs reliable level
        localisation, and supplies identical anatomically defined regions of
        interest to all grading baselines and proposed models.

  \item To establish competitive independent and contextual grading baselines on
        the public LumbarDISC development data, including fixed fusion,
        cross-attention and ordered inter-level context where reproducibility
        permits, so that improvements cannot be attributed to a weak comparator.

  \item To develop and evaluate anatomically aligned cross-sequence
        self-supervision in which positive relationships are defined by the same
        patient and lumbar level across MRI sequences using DICOM geometry.

  \item To quantify whether this anatomical self-supervision improves supervised
        grading under 10\%, 25\%, 50\% and 100\% label availability and whether it
        preserves sequence-specific diagnostic information.

  \item To develop and evaluate disease-conditioned adaptive sequence routing that
        learns patient- and target-dependent weights over available MRI evidence
        while explicitly representing sequence availability and acquisition quality.

  \item To test the router under both missing-sequence and corrupted-sequence
        conditions, and to compare its behaviour with concatenation, averaging and
        ordinary cross-attention under matched inputs and training budgets.

  \item To construct a heterogeneous disease--anatomy graph over the 25
        level--condition--laterality targets, with typed adjacent-level,
        same-level cross-condition and bilateral relations, and to compare it with
        independent heads, an ordered level transformer, a homogeneous graph,
        ungated message passing and shuffled topology.

  \item To test whether relational reasoning introduces adjacent-level false
        positives in naturally isolated lesions and to retain the graph only where
        any performance or robustness benefit is not purchased by clinically
        incoherent information propagation.

  \item To compare categorical, ordinal and clinically asymmetric cost-sensitive
        grading objectives, together with probability calibration and uncertainty
        analysis, while explicitly reporting both severe under-grading and severe
        over-grading.

  \item To quantify zero-shot transfer from the multinational public benchmark to
        the independent Rizgary Hospital cohort on the shared central-canal grading
        task, separating localisation, image-domain and reference-standard effects
        as far as the retrospective data allow.

  \item To estimate an annotation--performance recovery curve using controlled
        local adaptation sizes $N\in\{10,25,50,100\}$ and to compare where a limited
        trainable-parameter budget is most effectively placed.

  \item To report all principal comparisons with patient-level uncertainty,
        repeated training where feasible, contribution-specific ablations, negative
        results and explicit methodological limitations.
\end{enumerate}

These objectives define a staged experiment rather than a monolithic architecture.
A component is not retained because it appears in the proposal; it is retained only
if the corresponding experiment demonstrates a reproducible contribution or a
meaningful robustness, calibration or interpretability advantage.

% ===========================================================================
\section{Research Questions}
\label{sec:intro-rqs}

The research questions are fixed prospectively and correspond directly to the
methodological tests in Chapter~\ref{ch:methodology}.

\begin{description}
  \item[RQ1 --- Relational disease modelling.] Does representing lumbar disease
  targets as a typed heterogeneous graph improve grading over independent target
  heads, an ordered five-level transformer and a homogeneous graph, and are any
  gains concentrated where single-target evidence is weakest?

  \item[RQ2 --- Anatomically aligned self-supervision.] Does cross-sequence
  pretraining based on DICOM-defined anatomical correspondence improve grading,
  label efficiency and transfer compared with ImageNet initialisation, generic
  medical pretraining and ordinary augmentation-based contrastive learning?

  \item[RQ3 --- Disease-conditioned modality routing.] Can one model learn
  target- and patient-dependent sequence weights while remaining robust to both
  missing and quality-degraded MRI sequences, without requiring a separately
  trained network for every combination of sagittal T1, sagittal T2/STIR and axial
  T2?

  \item[RQ4 --- Cross-institutional transfer.] What is the zero-shot degradation
  when a benchmark-trained model is applied to a previously unseen Middle Eastern
  clinical cohort, how much of that degradation is attributable to localisation
  versus grading/reference-standard shift, and how does performance recover as the
  amount of local annotation available for adaptation increases?

  \item[RQ5 --- Clinically asymmetric error and uncertainty.] Do ordinal and
  cost-sensitive objectives reduce clinically consequential distant errors,
  particularly Severe$\rightarrow$Normal/Mild, and can calibrated uncertainty
  support selective prediction without overstating clinical safety?
\end{description}

These questions are intentionally answerable by negative as well as positive
results. If a typed graph performs no better than a homogeneous graph, RQ1 is still
answered. If anatomical self-supervision improves label efficiency but not external
transfer, RQ2 yields a more limited but still informative result. This design avoids
defining doctoral success as the requirement that every proposed component win.

% ===========================================================================
\section{Research Hypotheses}
\label{sec:intro-hypotheses}

The study uses directional hypotheses where the literature justifies a prediction,
but avoids arbitrary target accuracies that would have no defensible basis before
implementation.

\begin{description}
  \item[H1 --- Anatomical self-supervision.] Anatomical cross-sequence
  self-supervision will improve label-efficient grading and cross-institutional
  robustness relative to generic pretraining, particularly when only a fraction of
  benchmark labels is used.

  \item[H2 --- Adaptive routing.] Disease-conditioned routing with explicit
  availability masking, modality dropout and quality-aware gating will degrade less
  under missing-sequence and corruption experiments than fixed concatenation or
  ordinary cross-attention trained only on clean complete studies.

  \item[H3 --- Heterogeneous relational structure.] A typed heterogeneous graph
  will outperform independent target heads and a homogeneous graph if the proposed
  anatomical edge types carry useful information; performance under a shuffled-edge
  control should deteriorate if the topology itself matters.

  \item[H4 --- External degradation.] Zero-shot performance on the Rizgary Hospital
  cohort will be lower than internal benchmark performance because of scanner,
  protocol, population, localisation and reference-standard shift; the magnitude
  and pattern of that degradation are treated as findings rather than as failure of
  the research programme.

  \item[H5 --- Annotation-dependent recovery.] Local adaptation will improve
  performance as annotation increases, but no predetermined percentage of recovery
  is required. The shape of the recovery curve and the location in the network at
  which limited adaptation is most effective are empirical outcomes.

  \item[H6 --- Asymmetric grading error.] Cost-sensitive and ordinal objectives may
  reduce distant Severe$\rightarrow$Normal/Mild errors even if they do not improve
  aggregate macro-F1; any reduction must be evaluated against false-positive Severe
  predictions and against the cross-entropy baseline
  \cite{niemeyer2021a}.
\end{description}

The hypotheses will not be rewritten after inspection of the held-out public or
Rizgary test results. If an implementation decision changes the operational test of
a hypothesis, that change will be entered into the protocol changelog with its date
and rationale.

% ===========================================================================
\section{Conceptual Framework}
\label{sec:intro-conceptual}

The thesis can be understood as a progression from \emph{local evidence} to
\emph{structured evidence} and finally to \emph{transported evidence}.

At the first level, a model receives anatomically localised image evidence for a
single level and target. This is the baseline condition: if local evidence alone is
sufficient, more complex mechanisms should not be retained. At the second level,
the model is permitted to combine complementary MRI sequences, but the combination
must respect anatomical correspondence and sequence availability. Anatomical
self-supervision asks whether the same structure under different acquisitions can
produce a more label-efficient representation, while disease-conditioned routing
asks how strongly each sequence should contribute to a given target.

At the third level, a target is no longer treated as an isolated endpoint. Its
representation becomes a node in a typed graph whose relations correspond to three
known structures of the lumbar grading problem: the ordered chain of levels,
coexisting compartments at a level, and bilateral pairing. Message passing is soft
rather than deterministic; no rule states that disease at one level forces disease
at its neighbour. The graph is therefore a hypothesis about useful context, not a
hard-coded disease progression model.

At the fourth level, predictions are evaluated as ordinal and probabilistic outputs.
The model is judged not only by aggregate discrimination but by error distance,
Severe-class performance, calibration and the false-positive cost of any
cost-sensitive objective. This stage asks whether a statistically improved model is
also better aligned with the shape of clinically consequential grading error.

The final level is institutional transport. The representation, router, graph and
output head developed on public data are frozen before external evaluation. The
Rizgary Hospital cohort then asks a different question from internal validation:
what remains true when the acquisition and reporting environment changes? Only
after zero-shot performance has been measured is local adaptation permitted. This
ordering prevents the local test cohort from becoming an informal development set.

Conceptually, the thesis therefore tests the proposition

\begin{equation}
\text{robust grading}
\;=\;
f(\text{local anatomy},\;\text{cross-sequence evidence},\;
\text{target relations},\;\text{uncertainty},\;\text{domain}),
\label{eq:intro-conceptual}
\end{equation}

where the equation is explanatory rather than a literal computational model. Each
term corresponds to a controlled experiment, and no term is assumed necessary until
its ablation supports that conclusion.

% ===========================================================================
\section{Significance of the Study}
\label{sec:intro-significance}

\subsection{Scientific Significance}
\label{sec:intro-significance-scientific}

The principal scientific significance lies in changing the unit of reasoning from
an isolated image label to a structured set of disease targets. Many improvements in
medical imaging can be produced by additional parameters, larger backbones or more
aggressive tuning. Such gains are difficult to interpret mechanistically. By
contrast, this thesis attaches each proposed contribution to a specific hypothesis
and a destructive control: anatomical edges are compared with homogeneous and
shuffled topology; adaptive routing is compared with fixed fusion and challenged by
missing and corrupted modalities; anatomical self-supervision is compared with
generic pretraining and tested for over-invariance. The resulting evidence can show
not merely whether the full model scored higher, but \emph{why} it did or did not.

This matters because recent lumbar literature already demonstrates sophisticated
contextual models \cite{chai2026anatomy,batra2025mscan,nguyen2026crossspine}. The
research contribution cannot therefore be the generic use of attention, graphs or
multiple sequences. Its value lies in isolating the anatomical assumptions and
measuring them against alternatives under one controlled pipeline.

\subsection{Methodological Significance}
\label{sec:intro-significance-methodological}

The study places unusually strong emphasis on validity controls. Patient identity,
not slice or target, defines the data split; self-supervised learning is prohibited
from seeing held-out patients; localisation is evaluated separately from grading;
zero-shot transfer is measured before adaptation; the external reference standard is
kept distinct from routine report extraction; and every claimed model contribution
requires an ablation. These choices respond directly to recurring weaknesses in the
medical-imaging literature, where leakage, non-comparable splits and weak external
validation can produce optimistic performance estimates
\cite{compte2023are,wang2024machine}.

The transfer design is methodologically significant for a second reason. The local
hospital cohort is not used simply as another test set. The experiment attempts to
disentangle localisation failure, image-domain shift and reference-standard shift,
then measures recovery as local labels are added. This converts ``generalisation''
from a binary claim into a measurable process.

\subsection{Translational and Regional Significance}
\label{sec:intro-significance-regional}

The regional significance follows from the use of an independent Iraqi hospital
cohort that was not curated to mirror the public benchmark. Models developed on
large international datasets are often assumed to transfer into settings whose
scanner mix, acquisition practice, disease prevalence and reporting conventions were
not represented in development. Testing that assumption is particularly important
for institutions that do not have the resources to reproduce a large annotation
campaign locally. If useful performance can be recovered from tens rather than
hundreds of local annotations, that is a practical finding; if not, the negative
result is equally important because it quantifies the annotation burden required for
responsible transfer.

The study does not claim that one external hospital represents all of Iraq or the
Middle East. Rizgary Hospital provides evidence of transport to one genuinely
independent setting. Regional generalisation would require additional hospitals and
is therefore a future study rather than an inference drawn from one centre.

\subsection{Clinical Significance and its Boundary}
\label{sec:intro-significance-clinical}

The intended clinical significance is improved technical support for structured MRI
grading: consistency, reproducibility, calibrated confidence and explicit handling
of incomplete imaging. The thesis does not claim that this will improve treatment
outcomes, referral decisions or radiologist productivity. Those claims require a
prospective reader-assistance or impact study, as demonstrated by work that measures
reader performance directly \cite{lim2022improved}. The present study therefore
uses language such as ``decision-support potential'' and ``technical portability'',
not ``clinical benefit''.

% Written after Chapter 4. States only what the executed experiments support.
Of the benefits anticipated above, the executed experiments support
\emph{discrimination} and, within it, the reduction of clinically distant
errors. Quadratic weighted kappa improves by $+0.0177$ over a single-sequence
baseline, Severe$\rightarrow$Normal confusions fall from $5.7\%$ to $3.8\%$ and
recall on Severe targets rises from $61.1\%$ to $63.1\%$, the last two
attributable to the ordinal and cost-sensitive objective rather than to any
proposed structural mechanism.

Robustness to missing sequences improves, but the improvement is attributable to
modality dropout rather than to disease-conditioned routing, which changes
sequence reliance by an amount indistinguishable from zero. Label efficiency and
cross-institutional transfer were not evaluated and no claim is made for either.
Calibration is reported but was not the object of a confirmatory comparison.

No inference is drawn from these figures about patient benefit. The study
measures agreement with a reference standard on a retrospective benchmark; it
does not measure any clinical outcome.

% ===========================================================================
\section{Proposed Contributions}
\label{sec:intro-contributions}

At the protocol stage, the contributions below collectively form the proposed \textbf{AMOG-Net} framework and are \emph{proposed contributions}.
They become demonstrated thesis contributions only if the experiments support them.

\subsection{Contribution 1: Anatomically Aligned Cross-Sequence Self-Supervision}
\label{sec:intro-contrib-ssl}

The first proposed contribution uses DICOM-defined patient-space correspondence to
construct positive relationships between different MRI sequences depicting the same
patient and lumbar level. The objective is not to invent a new generic contrastive
loss; established contrastive or non-contrastive objectives may be used. The
contribution under test is the anatomical definition of correspondence and its value
for label efficiency, grading and transport. A dedicated preservation experiment
checks whether alignment erases sequence-specific diagnostic information.

\subsection{Contribution 2: Disease-Conditioned and Quality-Aware Sequence Routing}
\label{sec:intro-contrib-routing}

The second proposed contribution allows one study to be fused differently for
different disease targets. Routing depends on target identity, patient features,
sequence availability and quality evidence, and is evaluated under genuine or
simulated missingness and corruption. The contribution is not the existence of a
softmax gate or mixture-of-experts mechanism; it is the target-conditioned allocation
of clinically complementary evidence and the explicit robustness evaluation that
accompanies it.

\subsection{Contribution 3: Heterogeneous Disease--Anatomy Graph Reasoning}
\label{sec:intro-contrib-graph}

The third proposed contribution represents the 25 benchmark outputs as nodes in a
heterogeneous graph. Typed edges encode adjacent-level, within-level cross-condition
and bilateral relations. The graph includes a mechanism for preserving local
evidence when neighbour messages are unhelpful and is evaluated against independent,
ordered-transformer, homogeneous, ungated and shuffled-topology controls. A focal
isolated-lesion test determines whether relational reasoning improves uncertain
predictions without artificially spreading disease to adjacent normal levels.

\subsection{Contribution 4: A Structured Cross-Institutional Transfer Evaluation}
\label{sec:intro-contrib-transfer}

The translational contribution is an experimental design rather than a new network
layer. The selected public-data model is evaluated zero-shot on a fixed Rizgary
Hospital test set before any local adaptation. A separate local adaptation pool is
then sampled at controlled annotation sizes and used to compare constrained update
strategies. The analysis is strengthened by a verified-localisation control and,
where resources permit, harmonised independent reader grading of central canal
stenosis. This design estimates both the cost of transport and the cost of recovery.

\subsection{Supporting Contributions}
\label{sec:intro-contrib-supporting}

Ordinal and cost-sensitive objectives, calibration, uncertainty estimation and
parameter-efficient adaptation support the principal contributions but are not
claimed as original methods. Their value lies in testing whether the proposed
structured model behaves sensibly under clinically relevant error definitions and
limited local annotation. They are retained only when their contribution is
measurable.

% --------------------------------------------------------------------------
% Written at final submission, replacing the [LATER INTEGRATION] placeholder.
% Only components supported by the locked experiments are claimed. Components
% that failed are reported below as negative findings and remain in the
% research history and in the ablation ladder.
% --------------------------------------------------------------------------
\subsection{Demonstrated Contributions}
\label{sec:intro-contrib-demonstrated}

The experiments reported in Chapter~\ref{ch:results} support three
contributions. They are not the three proposed above, and the difference is
itself part of what the thesis contributes.

\paragraph{D1 --- A verified anatomy-aware multi-sequence grading system.}
AMOG-Net improves quadratic weighted kappa over a single-sequence baseline by
$+0.0177$ (95\% CI $[+0.0097, +0.0260]$), on every one of seven training seeds
and surviving false discovery rate correction across the pre-specified
comparison family. Two further comparisons also survive: typed heterogeneous
edges and the ordinal, cost-sensitive objective. The system, its frozen data
partition and its behavioural test suite are released.

\paragraph{D2 --- A controlled ablation protocol for anatomical priors.} Each
claimed mechanism is tested against a control matched in capacity rather than
against its own absence: a degree-preserving edge shuffle for graph topology, a
router-free fixed-fusion arm for sequence routing, an architecture-matched
untrained network as an attribution floor, and a label-shuffled model as a
pipeline-level negative control. Section~\ref{sec:results-threats} documents
three analyses that produced incorrect positive conclusions before the
corresponding control was introduced, in each case in the direction of the
hypothesis under test.

\paragraph{D3 --- Evidence that anatomical structural priors do not improve
grading at this data scale, with a mechanism.} Neither anatomically aligned
cross-sequence self-supervision nor disease-conditioned routing produces a
detectable improvement, and anatomically correct graph topology does not
outperform a degree-preserving random shuffle. These are bounds rather than
absences: the paired intervals constrain the three mechanisms to at most
$1.32\%$, $1.55\%$ and $1.05\%$ of baseline performance respectively. Gradient-weighted class activation
mapping shows why: the convolutional encoder already concentrates $2.4$ times the
chance share of its attention on the annotated target before any prior is added.
The localisation these priors were designed to supply is largely accomplished
upstream of them.

\subsection{Proposed Contributions Not Supported}
\label{sec:intro-contrib-negative}

Contribution~1 (anatomical cross-sequence self-supervision) is rejected on
accuracy, on cross-sequence robustness under controlled input ablation, and on
attention concentration. The pretext task is learnable and was learned ---
validation InfoNCE $1.0771$ against a chance value of $3.4657$ --- but the
representation confers no measurable downstream benefit.

Contribution~2 (disease-conditioned routing) divides. The router does learn
clinically sensible target-dependent sequence weights, reproducing the expected
allocation in all fifteen runs in which a gate is present. It does not improve
grading ($+0.0053$ QWK, five seeds of seven, bounded at $1.55\%$ of baseline)
and it does not improve robustness to missing sequences.
Its allocation tracks the annotation structure of the benchmark rather than
creating a dependence of its own.

Contribution~3 (heterogeneous disease--anatomy graph) is partly supported. Typed
relational structure improves grading over a homogeneous graph ($+0.0093$ QWK on
every one of seven seeds, surviving correction at $p_{\mathrm{FDR}} = 0.009$);
the anatomical identity of the relations does not. The resulting
claim is narrower than proposed and more specific than the original: relational
typing carries information, anatomical adjacency does not.

Contribution~4 (cross-institutional transfer evaluation) was not executed. Cohort
preparation is complete and is reported in Section~\ref{sec:results-rq4};
adjudication of conflicting reference labels and a correct de-identification
procedure remain outstanding. It is retained here as unexecuted rather than
withdrawn.

% ===========================================================================
\section{Scope and Delimitations}
\label{sec:intro-scope}

Clear boundaries are necessary because lumbar MRI supports a much broader clinical
and computational agenda than one PhD can test.

\subsection{Clinical Scope}
\label{sec:intro-scope-clinical}

The primary public task is degenerative lumbar disease severity grading at five
levels from L1--L2 through L5--S1. On LumbarDISC, the full target space contains
central canal stenosis, left and right neural foraminal narrowing, and left and
right subarticular stenosis, each graded as Normal/Mild, Moderate or Severe
\cite{richards2026the}. The study does not attempt to diagnose every cause of low
back pain, infer symptoms, determine surgical candidacy or predict long-term patient
outcomes.

The primary local transfer task is narrower. Rizgary Hospital reports support a
defensible central-canal comparison more readily than the complete benchmark schema.
The confirmatory external analysis is therefore restricted to central canal stenosis
at the five lumbar levels unless new expert grading expands the local reference
standard. Disc bulge, protrusion and extrusion are retained as a separate local
morphology task and are not heuristically converted into stenosis grades.

\subsection{Imaging Scope}
\label{sec:intro-scope-imaging}

The principal imaging inputs are sagittal T1, sagittal T2/STIR and axial T2 MRI.
Other sequences may appear in clinical studies but do not define the primary model
unless a justified later protocol amendment is made. DICOM geometry is preserved
because cross-plane correspondence is part of the scientific design. PNG or other
rendered formats may be used for visual inspection but do not replace the geometric
metadata required for alignment.

\subsection{Methodological Scope}
\label{sec:intro-scope-method}

Anatomical localisation and segmentation are enabling technologies. They are
measured and controlled because a poor crop can invalidate grading, but a new
localisation architecture is not the headline doctoral contribution. Likewise,
backbone selection is controlled rather than treated as a novelty competition. The
research may use ResNet-, EfficientNet-, Swin- or other contemporary encoders, but
the primary backbone is frozen before the main contribution-specific ablations.

The graph topology is defined anatomically rather than discovered freely from the
test data. Adjacent levels, same-level conditions and bilateral counterparts form the
pre-specified relation families. Learning the globally optimal topology is a
different research problem and is excluded from the primary scope.

\subsection{Data and Institutional Scope}
\label{sec:intro-scope-data}

Model development uses the public LumbarDISC benchmark and, where permitted, SPIDER
for anatomical localisation or segmentation support. The locally held labelled
competition subset contains approximately 1,975 studies at the protocol stage;
\torecord{final public development, validation and held-out counts after quality
control}. The local project inventory currently identifies 299 narrative reports
and 294 candidate matched multi-sequence cases, while a wider raw DICOM inventory
contains additional studies. These numbers are not presented as the final cohort
until one-to-one reconciliation is complete:
\torecord{final Rizgary case-flow counts and reasons for exclusion}.

Before analytical use of the local cohort, the Rizgary data-preparation phase will
follow the approved institutional governance protocol. A controlled research copy will
be de-identified before access by the model-development team, and the permitted storage
and compute environment will be verified before local processing begins. These safeguards
precede, rather than form part of, model development.

Rizgary Hospital is a single external institution. The study can therefore support a
claim of transport to that independent setting, not universal Middle Eastern
performance. Additional regional hospitals would be required for a multicentre
regional generalisation claim.

\subsection{Ethical and Governance Scope}
\label{sec:intro-scope-ethics}

The local component is retrospective and does not alter patient care. Formal use of
the local data remains conditional on the relevant institutional permissions:
\toconfirm{hospital/ethics approval reference, date and responsible institution};
\toconfirm{retrospective consent or waiver basis}; and
\toconfirm{permitted compute environment, including whether external cloud/GPU
processing is authorised}. Identifiable raw DICOM is not supplied to students or
external compute; a controlled de-identified research copy is required before local
analysis.

\subsection{Evaluation Scope}
\label{sec:intro-scope-evaluation}

The thesis evaluates technical grading performance, robustness, calibration,
computational cost and cross-institutional portability. The unit of statistical
independence is the patient. It does not claim prospective clinical impact,
radiologist replacement or treatment benefit. Reader-assistance evaluation is a
logical future study if the technical system proves sufficiently reliable.

% ===========================================================================
\section{Assumptions and Operational Definitions}
\label{sec:intro-definitions}

The following terms are used consistently throughout the thesis.

\begin{longtable}{p{3.6cm}p{10.0cm}}
\toprule
\textbf{Term} & \textbf{Operational meaning in this thesis}\\
\midrule
\endfirsthead
\toprule
\textbf{Term} & \textbf{Operational meaning in this thesis}\\
\midrule
\endhead
Lumbar level & One of L1--L2, L2--L3, L3--L4, L4--L5 or L5--S1.\\
Target & A level--condition--laterality output on LumbarDISC; 25 targets per study at full benchmark scope.\\
Severity grade & Normal/Mild, Moderate or Severe, encoded as an ordered three-class outcome.\\
Local evidence & Image-derived representation for one target before graph message passing.\\
Anatomical correspondence & Physical correspondence between acquisitions defined through DICOM patient-space geometry and lumbar localisation, not through image appearance alone.\\
Sequence routing & Learned allocation of weight over available MRI sequence features for a specific patient and target.\\
Missing modality & An expected MRI sequence is absent and explicitly unavailable to the model.\\
Degraded modality & A sequence is present but technically compromised, for example by motion, truncation, noise, slice loss or poor coverage.\\
Heterogeneous graph & A graph whose nodes are disease--anatomy targets and whose edges belong to clinically defined relation types.\\
Zero-shot transfer & Evaluation of the frozen public-data model on the fixed local test set before local fine-tuning, threshold selection or local calibration.\\
Few-shot adaptation & Model adaptation using only the separate local adaptation pool under pre-specified annotation budgets.\\
Reference-standard shift & Difference between the process used to generate labels in the benchmark and the process used to establish the local external reference.\\
Selective prediction & Model abstention or referral based on uncertainty; this is a technical behaviour, not evidence of a validated clinical workflow.\\
\bottomrule
\end{longtable}

These definitions are intentionally narrower than ordinary-language usage. In
particular, ``external validation'' does not mean that a model has been proven safe
for deployment, and ``uncertainty'' does not mean that the network has measured a
patient's clinical risk.

% ===========================================================================
\section{Overview of the Research Design}
\label{sec:intro-design-overview}

The research is a retrospective, multi-dataset medical-imaging methods study. Its
logic is cumulative, but the conclusions remain component-specific.

\begin{enumerate}
  \item \textbf{Data and anatomy foundation.} Public data are inventoried,
        partitioned by patient, reconstructed in DICOM patient space, anatomically
        localised and converted into disease-specific 2.5D regions of interest.

  \item \textbf{Strong baseline establishment.} An independent ROI classifier and
        competitive multiview/context baselines establish what can be achieved
        without the proposed relational formulation.

  \item \textbf{Representation experiment.} Anatomically aligned cross-sequence
        self-supervision is evaluated against generic and medical pretraining under
        several label fractions and sequence-specific preservation probes.

  \item \textbf{Fusion experiment.} Disease-conditioned sequence routing is tested
        against fixed and attention-based fusion under complete, missing and
        corrupted sequence conditions.

  \item \textbf{Relational experiment.} Homogeneous and heterogeneous target graphs
        are compared with independent and ordered-transformer controls, including
        shuffled topology and isolated-lesion stress testing.

  \item \textbf{Output and uncertainty experiment.} Categorical, ordinal and
        cost-sensitive objectives are evaluated together with calibration and
        uncertainty, with both under-grading and over-grading reported.

  \item \textbf{External transport experiment.} The selected public-data model is
        frozen and evaluated zero-shot on the fixed Rizgary Hospital test set. A
        verified-ROI control and harmonised reader reference are used where feasible
        to interpret the source of degradation.

  \item \textbf{Adaptation experiment.} Controlled local annotation budgets are
        sampled from a separate adaptation pool to estimate how quickly performance
        recovers and where a limited trainable-parameter budget is most effective.
\end{enumerate}

This order deliberately prevents the final architecture from becoming an indivisible
engineering object. Each stage introduces one assumption that can be rejected
without invalidating the remainder of the thesis. The final model is therefore the
simplest configuration supported by the evidence, not necessarily the configuration
containing every proposed module.

% ===========================================================================
\section{Expected Outcomes and Criteria for Scientific Success}
\label{sec:intro-success}

The protocol does not define success as reaching a predetermined accuracy. Such a
threshold would be difficult to justify because performance depends on target,
severity distribution, reference-standard reliability and the final patient split.
Scientific success is instead defined by whether the research questions can be
answered under controlled conditions.

For RQ1, success means establishing whether typed anatomical relations add value
beyond capacity-matched controls and whether they avoid adjacent-level
contamination. A null result is informative if the comparison is sufficiently
powered and the topology controls are valid. For RQ2, success means establishing
whether anatomical cross-sequence pairing improves label efficiency, transfer or
neither, while documenting whether sequence-specific information is preserved. For
RQ3, success means quantifying robustness under missing and degraded modalities and
determining whether adaptive routing adds value over simpler fusion. For RQ4,
success means producing an interpretable zero-shot transfer estimate and a controlled
adaptation curve, even if the external drop is large. For RQ5, success means
characterising the trade-off between severe under-grading, severe over-grading and
probability reliability rather than selecting a loss solely because it raises one
aggregate metric.

At final submission, the thesis will distinguish three categories explicitly:
\emph{supported contributions}, \emph{negative findings}, and \emph{unresolved
questions}. This is preferable to rewriting the original proposal so that only
successful components appear to have been planned.

% Written after Chapter 4. High-level answers only; metrics live in Chapter 4.
\begin{description}
  \item[RQ1 --- Relational disease modelling.] \textbf{Partly supported.} A typed
  heterogeneous graph improves grading over independent heads and over a
  homogeneous graph, on every seed. The anatomical identity of the relations does
  not contribute: an anatomically correct topology does not outperform a
  degree-preserving random shuffle.

  \item[RQ2 --- Anatomically aligned self-supervision.] \textbf{Not supported.}
  The pretext task is learnable and was learned, but the resulting representation
  yields no detectable improvement in grading, in robustness to withheld
  sequences, or in where the model attends.

  \item[RQ3 --- Disease-conditioned modality routing.] \textbf{Mechanism
  supported, benefit not.} The router reliably learns the clinically expected
  sequence weights, but a router-free model shows the same sequence dependence
  under intervention, and routing improves neither accuracy nor missing-sequence
  robustness.

  \item[RQ4 --- Cross-institutional transfer.] \textbf{Not executed.} Cohort
  preparation is complete; reference-label adjudication and de-identification
  remain outstanding.

  \item[RQ5 --- Clinically asymmetric error and uncertainty.] \textbf{Supported
  for error asymmetry.} Ordinal and cost-sensitive objectives improve both
  aggregate agreement and recall on Severe targets, on every seed, and form the
  most reliable single contributor to final performance. Selective prediction
  under calibrated uncertainty was not evaluated confirmatorily.
\end{description}

% ===========================================================================
\section{Thesis Organisation}
\label{sec:intro-organisation}

At the current protocol stage, the thesis is organised around the progression from
problem definition to evidence.

\textbf{Chapter 1 --- Introduction} defines the clinical and methodological
motivation, research problem, gaps, aim, objectives, research questions, hypotheses,
proposed contributions, significance and scope of the thesis.

\textbf{Chapter 2 --- Literature Review} examines the clinical grading systems that
define the reference labels, MRI sequence--pathology correspondence, localisation
and segmentation, backbone and contextual architectures, class imbalance and
ordinal grading, public datasets, external validation, interpretability and the
specific gaps that motivate the present study.

\textbf{Chapter 3 --- Methodology} translates those gaps into a protocol-grade
experimental design. It defines the public and local cohorts, leakage controls,
DICOM reconstruction, localisation and ROI construction, self-supervised
pretraining, adaptive routing, heterogeneous graph reasoning, grading objectives,
calibration, ablations, external validation, domain adaptation, statistical analysis
and reproducibility safeguards.

\textbf{Chapter 4 --- Results and Analysis} will report the staged E0--E8
experiments, including component-specific ablations, robustness analyses,
calibration, zero-shot external transfer, few-shot adaptation and failure analysis.
% Written after Chapter 4. Retained experiment groups; research questions unchanged.
The retained experiment groups are the eight-rung ablation ladder E0--E7 with two
matched controls on the graph configuration, a controlled input ablation testing
sequence dependence by intervention rather than by learned allocation, and a
Grad-CAM attribution analysis against an architecture-matched untrained floor.
The cross-institutional transfer group was prepared but not executed and is
reported as such. The research questions are unchanged from those fixed in
Section~\ref{sec:intro-rqs}.

\textbf{Chapter 5 --- Discussion and Conclusions} will interpret the findings
against the literature, distinguish demonstrated contributions from negative
results, state limitations, discuss implications for transport to other institutions
and identify the next studies required before any clinical-impact claim.

\toconfirm{final institutional thesis structure if the university requires results
and discussion as separate chapters, publication-based chapters, or an additional
conclusion chapter}. The scientific ordering above should remain recognisable even
if the administrative chapter numbering changes.

% ===========================================================================
\section{Rules for Completing Prospective Placeholders}
\label{sec:intro-placeholders}

Because this chapter is written before completion of implementation, several fields
are intentionally prospective. The following rules govern their later completion.

\begin{enumerate}
  \item \textbf{[TO CONFIRM]} is used only for institutional or governance facts
        that are currently unknown, such as an approval reference or permitted
        compute environment. These fields are replaced only with documented facts.

  \item \textbf{[TO RECORD AFTER IMPLEMENTATION]} is used for executed technical or
        cohort quantities that cannot responsibly be specified before the pipeline
        runs, such as final quality-controlled sample counts.

  \item \textbf{[LATER INTEGRATION]} is used for a small number of places where the
        final thesis may state the principal findings or update the novelty boundary.
        These additions may report evidence; they must not alter the prospective
        research questions or make an unsuccessful hypothesis disappear.

  \item Any substantive protocol change made after results become visible is
        documented with date, reason and affected analysis. Confirmatory and
        exploratory analyses remain distinguishable.

  \item The introduction will not be retrofitted to imply that the final retained
        architecture was inevitable. The final thesis should preserve the fact that
        ACSSL, adaptive routing, heterogeneous graph reasoning and supporting losses
        were hypotheses subjected to testing.
\end{enumerate}

These rules turn placeholders into a transparency mechanism rather than an invitation
to hindsight bias.

% ===========================================================================
\section{Chapter Summary}
\label{sec:intro-summary}

This chapter has defined the thesis as a study of structured and transferable lumbar
MRI grading rather than as a competition between image backbones. The clinical task
is inherently multi-level, multi-condition, bilateral and multi-sequence, while its
reference labels are ordinal and imperfectly reproducible. The technical literature
has established that localisation should precede grading, that sophisticated
multiview systems can perform strongly on public data, and that external performance
frequently degrades. Those achievements narrow rather than remove the research gap.

The thesis therefore asks whether three explicit structural mechanisms add value:
anatomically aligned cross-sequence self-supervision, disease-conditioned and
quality-aware sequence routing, and a typed heterogeneous graph over
level--condition--laterality targets. These mechanisms are evaluated component by
component against strong controls rather than bundled into a single architecture and
assumed beneficial. Supporting ordinal, cost-sensitive and uncertainty methods are
used to test the clinical shape of error without being advertised as standalone
novelty.

The translational component separates model development from external evaluation.
LumbarDISC supplies the public development benchmark; SPIDER may support anatomical
localisation; and the independent Rizgary Hospital cohort is reserved for zero-shot
central-canal transfer and controlled local adaptation. The local reference process
is treated explicitly as a potential source of shift rather than assumed equivalent
to the public benchmark. Clinical impact remains outside the scope until a separate
reader-assistance or prospective study is conducted.

The aim, objectives, RQ1--RQ5 and H1--H6 stated here form the prospective frame for
Chapters~\ref{ch:litreview} and~\ref{ch:methodology}. The remainder of the thesis is
therefore not required to prove that every proposed mechanism succeeds. It is
required to determine, with reproducible evidence, which assumptions survive
controlled testing and institutional transfer.

% Written at final submission.
In summary, the complete system improves on a single-sequence baseline
reproducibly and by a margin that survives correction for multiple comparisons
while remaining far too small to alter clinical practice,
and the objective function designed around clinically asymmetric error is the
most reliable single contributor to that improvement. Of the three proposed
structural mechanisms, relational typing is supported, anatomical topology and
disease-conditioned routing are not, and anatomically aligned self-supervision is
rejected on every axis on which it was tested; attribution analysis indicates
that the convolutional encoder already performs the localisation these mechanisms
were intended to supply. The principal limitations are that all results are
internal to one benchmark, that cross-institutional transfer was not executed,
and that the study is underpowered for single-component effects of the magnitude
observed.

% ===========================================================================
\printbibliography[heading=bibintoc,title={References}]

\end{document}
 from the main thesis file.
% ===========================================================================

\documentclass[12pt,a4paper]{report}

\usepackage{lmodern}
\usepackage[T1]{fontenc}
\usepackage[utf8]{inputenc}
\usepackage[english]{babel}
\usepackage{csquotes}
\usepackage{microtype}
\usepackage[margin=2.5cm]{geometry}
\usepackage{setspace}
\onehalfspacing

\usepackage{amsmath,amssymb}
\usepackage{booktabs}
\usepackage{longtable}
\usepackage{array}
\usepackage{graphicx}
\usepackage{thesisstyle}   % shared TikZ figure styles; see thesisstyle.sty
\usepackage[hidelinks]{hyperref}

\usepackage[
  backend=biber,
  style=numeric-comp,
  sorting=none,
  maxbibnames=6,
  minbibnames=3,
  giveninits=true,
  doi=true,
  url=false,
  isbn=false
]{biblatex}

% Production bibliography cleaned for thesis compilation. Citation keys are
% unchanged from the working literature inventory.
\addbibresource{lumbar_spine_mri_ai_references_clean.bib}

\newcommand{\toconfirm}[1]{\textbf{[TO CONFIRM: #1]}}
\newcommand{\torecord}[1]{\textbf{[TO RECORD AFTER IMPLEMENTATION: #1]}}
\newcommand{\laterinsert}[1]{\textbf{[LATER INTEGRATION: #1]}}

\begin{document}

\chapter{Introduction}
\label{ch:introduction}

% ===========================================================================
\section{Introduction and Thesis Position}
\label{sec:intro-position}

Lumbar degenerative disease presents an unusually demanding problem for medical
image analysis. The anatomical region is compact, but the diagnostic task is not:
one examination contains several intervertebral levels, several pathological
compartments, left--right distinctions, multiple magnetic resonance imaging (MRI)
sequences acquired in different planes, and severity labels whose boundaries are
partly subjective. A clinically meaningful automated system must therefore do more
than recognise that degeneration is present. It must identify \emph{what} is
abnormal, \emph{where} it is abnormal, \emph{how severe} the abnormality is, and
\emph{which imaging evidence} supports that judgement. It must also remain
interpretable as a measurement of radiological grading rather than be mistaken for
an autonomous clinical diagnosis.

The need for such systems arises from both disease burden and reporting workload.
Low back pain remains one of the largest contributors to disability worldwide, and
lumbar degenerative change is a frequent structural correlate examined by MRI
\cite{compte2023are,wang2024machine}. Routine interpretation is repetitive because
the reader must inspect multiple levels across sagittal and axial acquisitions,
while the severity of central canal, foraminal and subarticular narrowing is judged
against morphological criteria that are not perfectly reproducible between human
readers \cite{lurie2008reliability,schizas2010qualitative,lee2011a,lee2010a}.
Automation is attractive in this setting not because a neural network is assumed to
possess a superior biological truth, but because a reproducible system may help
standardise a high-volume structured task whose human reference process itself has
known variability.

The technical field has advanced rapidly. Early systems demonstrated that lumbar
anatomy could be localised and individual discs graded with convolutional networks;
more recent pipelines combine explicit localisation, 2.5D regions of interest,
multiview fusion, attention and large public benchmarks
\cite{lu2018deepspine,pan2021automatically,batra2025mscan,richards2026the}. The
central problem has consequently shifted. The thesis no longer asks whether a
convolutional neural network or transformer can classify a lumbar MRI in isolation.
It asks whether the \emph{structure of the grading problem itself} has been
represented faithfully enough.

This thesis takes the position that lumbar MRI grading is a structured prediction
problem with four coupled sources of information. First, disease occurs at ordered
anatomical levels rather than at unrelated image locations. Second, several
conditions coexist within one motion segment and lateral compartments occur as
paired left--right structures. Third, the diagnostic value of an MRI sequence is
target-dependent: sagittal T1, sagittal T2/STIR and axial T2 provide complementary
rather than interchangeable evidence \cite{hallinan2021deep}. Fourth, a model that
performs well on one benchmark may degrade when scanners, protocols, population and
reference-standard conventions change across institutions
\cite{mcsweeney2023external,zhang2023deep,guan2022domain}.

The research therefore evaluates three methodological contributions and one
translational validation programme. The integrated framework evaluating these components
is designated \textbf{AMOG-Net} (\textbf{A}natomy-aware \textbf{M}ulti-sequence \textbf{O}rdinal \textbf{G}raph Network). The first contribution is anatomically aligned cross-sequence
self-supervision, in which DICOM patient-space correspondence defines which views
should agree during representation learning. The second is disease-conditioned
adaptive sequence routing, in which the model learns target- and patient-dependent
weights over available MRI sequences and is explicitly tested under both missing
and degraded inputs. The third is a heterogeneous disease--anatomy graph whose
nodes represent level--condition--laterality targets and whose typed edges encode
adjacent-level, same-level cross-condition and bilateral relations. These are then
evaluated under a deliberately separate transfer programme: development occurs on
a multinational public benchmark; the selected model is frozen; zero-shot
performance is measured on an independent Rizgary Hospital cohort; and controlled
few-shot adaptation is performed only afterwards.

Several related techniques are necessary but are not claimed as doctoral novelty.
Anatomical localisation, segmentation, 2.5D region construction, conventional CNN
or transformer encoders, class-imbalance handling, ordinal loss functions,
calibration and parameter-efficient fine-tuning are established tools. They are
used to construct a fair experiment around the thesis contributions rather than
presented as original merely because they appear in the same architecture. This
novelty boundary is important in a field where anatomy-guided contextual grading,
cross-sequence attention and graph-based disc modelling have already been reported
\cite{chai2026anatomy,nguyen2026crossspine,baur2024automated}. The question is not
whether any one of those mechanisms exists elsewhere; it is whether the particular
representation tested here---typed relationships among the benchmark's disease
outputs, anatomical cross-sequence supervision, target-conditioned evidence
allocation and explicit institutional transfer---adds measurable knowledge beyond
strong matched baselines.

% Written after the confirmatory experiments were locked. Reports the outcome
% only; the problem statement, research questions and hypotheses are unchanged.
The principal empirical outcome is that the complete system improves agreement
with the reference standard over a single-sequence baseline by $+0.0177$
quadratic weighted kappa, on every one of seven training seeds and after
correction for multiple comparisons, but that the anatomical mechanisms proposed
to achieve this do not account for the improvement: attribution analysis
indicates the convolutional encoder already localises the graded structure
without them.

% ===========================================================================
\section{Background and Motivation}
\label{sec:intro-background}

\subsection{Clinical Motivation: A Repetitive but Structured Imaging Task}
\label{sec:intro-clinical-motivation}

Lumbar degeneration is not a single lesion. It is a family of related structural
changes involving the intervertebral disc, facet joints, ligamentum flavum,
vertebral endplates and the neural spaces that these structures surround. Disc
dehydration and height loss alter load distribution; annular bulging or herniation
may encroach on the canal and foramina; posterior element hypertrophy may further
reduce the space available to neural tissue. These changes are assessed at several
levels, most commonly from L1--L2 through L5--S1, and their significance depends on
which anatomical compartment is narrowed.

Three stenotic compartments are especially relevant to the benchmark task used in
this thesis. Central canal stenosis concerns the thecal sac and cauda equina,
foraminal narrowing concerns the exiting nerve roots, and subarticular or lateral
recess stenosis concerns the traversing roots. Their visual criteria are related but
not identical, and neither are the acquisitions from which they are best judged.
Central canal and lateral recess morphology depend heavily on axial T2 information,
whereas neural foraminal assessment benefits particularly from sagittal or
parasagittal T1 because perineural fat provides the relevant contrast
\cite{hallinan2021deep,schizas2010qualitative,lee2010a}. A model that treats all
sequences as equivalent therefore ignores a clinical property of the task before
learning has even begun.

The labels themselves are another source of complexity. Morphological grading
systems convert continuous anatomical variation into ordinal categories. Human
agreement is strongest for some compartments and weaker for others: the reliability
study of Lurie et al. found substantially better inter-reader agreement for central
canal stenosis than for foraminal and subarticular findings
\cite{lurie2008reliability}. This matters methodologically. If the reference standard
contains boundary uncertainty, model evaluation must distinguish gross errors from
one-grade disagreements and must interpret performance relative to the reliability
of the grading instrument. Apparent model error and reference-standard disagreement
are not always separable.

This combination---high examination volume, repeated level-wise assessment,
complementary sequence evidence and imperfect categorical boundaries---makes lumbar
MRI a suitable domain for artificial intelligence, but only if the automation is
framed correctly. The objective is not to replace full clinical interpretation.
Symptoms, neurological findings, prior imaging and management context remain outside
the information available to an image-only grading system. The narrower objective is
to make a defined radiological grading task more consistent, measurable and
portable. Reader-assistance studies support this framing: deep-learning assistance
has been shown to reduce reporting time and improve or preserve inter-reader
agreement, with the greatest gains often occurring among less experienced readers
\cite{lim2022improved}. This thesis therefore treats the system as a potential
component of decision support, not as an autonomous treatment decision-maker.

\subsection{Technical Motivation: From Classification to Structured Prediction}
\label{sec:intro-technical-motivation}

The evolution of lumbar MRI automation has progressively narrowed the unit of
prediction. Early work often asked whether a lesion or abnormality was present in an
image. Subsequent systems localised discs, graded individual levels, and shared
encoders across several pathologies. Public benchmark design has now made the task
more explicit. The RSNA Lumbar Degenerative Imaging Spine Classification resource,
published as LumbarDISC, contains 2,697 patients and 8,593 MRI series collected
across eight institutions in six countries \cite{richards2026the}. The associated
challenge requires five condition/laterality targets at five lumbar levels, producing
25 ordinal outputs per study \cite{rsna2024rsna,richards2026the}. A correct model
therefore does not simply output ``mild'', ``moderate'' or ``severe'' for an image;
it attaches one of those grades to a particular condition at a particular level and,
for lateral compartments, a particular side.

This target structure exposes a limitation in the prevailing formulation. Most
systems share image features but ultimately emit independent output heads. Shared
features allow several tasks to benefit from a common representation, but a
prediction at L4--L5 does not explicitly inform an uncertain prediction at L3--L4;
a confident left foraminal prediction does not interact structurally with its right
counterpart; and the central, foraminal and subarticular targets at one motion
segment do not exchange typed messages simply because they coexist anatomically.
The representation is therefore multi-task but not fully relational.

Graph learning provides a natural language for testing whether those relationships
carry useful information. Graph neural networks are not novel by themselves, and a
lumbar application has already reconstructed an individual disc as a point cloud and
graded it with a graph network \cite{baur2024automated}. That precedent is important
because it prevents an unjustified claim that ``using a graph in lumbar MRI'' is new.
The unresolved question is narrower: whether the \emph{prediction targets} can be
represented as nodes connected by clinically motivated relation types, and whether
those relations improve grading beyond independent heads, an ordered level
transformer and a homogeneous graph. The distinction between a graph over points
within one disc and a graph over disease--anatomy targets across the spine is the
novelty boundary this thesis tests.

A second technical motivation concerns multimodal evidence. Recent methods have
already shown that cross-view and cross-sequence attention can improve lumbar
assessment \cite{batra2025mscan,nguyen2026crossspine}, while anatomy-guided
inter-level context has also become an active direction \cite{chai2026anatomy}.
These developments mean that novelty cannot rest on ``using multiple sequences'' or
``using attention''. The remaining question is whether evidence allocation should be
\emph{conditioned on the target and the patient}, and whether one trained model can
remain operational when an expected sequence is unavailable or present but of poor
quality. A fixed fusion rule assumes the same contribution of sagittal T1, sagittal
T2/STIR and axial T2 for every target and every study. Clinical sequence--pathology
correspondence makes that assumption doubtful.

A third motivation is representation learning. Generic ImageNet pretraining is
convenient but does not exploit the particular correspondence available in MRI: the
same anatomical level can be observed under different contrasts and planes. Earlier
spinal work demonstrated that self-supervision derived from patient-level or
longitudinal information can reduce dependence on manual labels
\cite{jamaludin2017selfsupervised}, and disc-centric contrastive work continues that
line \cite{acharya2026disccentric}. DICOM geometry makes a stronger, directly
anatomical pairing possible. If sagittal and axial observations can be matched in
patient space, then the correspondence itself can serve as supervision without
pretending that the images should look visually alike.

Finally, there is a translational motivation. Systematic reviews repeatedly find
that internally impressive medical-imaging models underperform when re-evaluated
outside their development distribution \cite{compte2023are,wang2024machine}.
Lumbar-specific external validations show the same pattern, with fine-grained
grading generally less portable than localisation or detection
\cite{mcsweeney2023external,zhang2023deep}. A model intended for eventual use in
Iraq or the wider region therefore cannot be justified solely by performance on a
multinational benchmark. The model must be transported to a genuinely independent
hospital cohort, its zero-shot degradation measured before local tuning, and the
amount of local annotation required to recover performance quantified explicitly.
That annotation--performance curve is more actionable than a single statement that
``domain adaptation improved accuracy''.

\subsection{Why the Problem Should Be Addressed Now}
\label{sec:intro-why-now}

Three developments make the present study timely. First, LumbarDISC supplies a
large, multi-institutional and anatomically structured benchmark whose target schema
is rich enough to test relational modelling rather than generic image
classification \cite{richards2026the}. Second, public anatomical resources such as
SPIDER reduce the need to claim localisation or segmentation as the central
contribution \cite{vandergraaf2024lumbar}. Third, the current methodological
literature is sufficiently mature to provide strong controls: multi-view
cross-attention, anatomy-guided contextual models, contemporary CNN and transformer
backbones, self-supervised pretraining, calibration and parameter-efficient transfer
all exist as comparators or supporting tools. The research can therefore ask a more
precise question than whether a new architecture attains a high score.

The resulting opportunity is methodological rather than merely computational. The
field now has enough data and enough baseline competence to test whether explicit
clinical structure provides value over capacity alone. If a typed target graph,
anatomical self-supervision or disease-conditioned router fails to outperform its
control, that negative result is still informative. Conversely, if a gain disappears
when topology is shuffled, when a simpler transformer is matched for capacity, or
when the model is transported to another hospital, the interpretation must change.
This insistence on falsifiable components is part of the motivation for the study,
not an afterthought added to evaluation.

% ===========================================================================
\section{Problem Statement}
\label{sec:intro-problem}

Despite rapid progress in automated lumbar MRI analysis, four connected
methodological problems remain.

\subsection{Problem 1: The Output Space Is Structured but Commonly Treated as Independent}
\label{sec:intro-problem-relational}

The benchmark output space contains ordered lumbar levels, multiple conditions at
each level and bilateral counterparts. These targets arise from one connected
anatomical system, yet the dominant design predicts them independently after a
shared encoder. This separation is convenient for implementation but does not test
whether adjacent-level context, within-level co-occurrence or bilateral relations
carry predictive information. Existing graph work in lumbar MRI does not resolve
this question because its graph is constructed within a disc rather than across the
disease targets \cite{baur2024automated}. CrossSpine introduces anatomical priors
and level-specific behaviour but still does not formulate the full 25 target outputs
as a typed communication graph \cite{nguyen2026crossspine}. The first problem is
therefore not the absence of contextual modelling in general; it is the absence of a
controlled test of explicit disease--anatomy relations among the prediction targets.

A relational model also carries a specific risk: useful context can become
``information contagion''. A severe lesion at one level should not cause the model
to elevate an adjacent normal level simply because degeneration often co-occurs.
Any graph proposed as a solution must therefore preserve local visual evidence,
include non-relational and homogeneous controls, and be stress-tested on naturally
isolated lesions. Without those controls, improved aggregate performance would not
establish that the anatomical topology is beneficial.

\subsection{Problem 2: Sequence Fusion Is Usually Fixed While Evidence Is Target- and Quality-Dependent}
\label{sec:intro-problem-routing}

Sagittal T1, sagittal T2/STIR and axial T2 are complementary acquisitions. The
foramina, central canal and lateral recess are not equally well represented in each
view, and routine data may be incomplete or technically degraded. Strong recent
systems address multi-view fusion through cross-attention
\cite{batra2025mscan,nguyen2026crossspine}, but a remaining deployment question is
whether one model can dynamically allocate evidence by condition, level, patient and
sequence quality. A model trained only on complete, high-quality benchmark studies
may fail when one sequence is missing, truncated or motion-corrupted. Missingness
and corruption are also different problems: an absent sequence can be masked,
whereas a degraded sequence remains present and can attract misleading attention.
The second problem is therefore to learn adaptive sequence allocation without
turning the router into an unexamined source of instability.

\subsection{Problem 3: Cross-Sequence Correspondence Is Available but Underused as Supervision}
\label{sec:intro-problem-ssl}

Self-supervised learning can reduce dependence on costly expert labels, but common
contrastive objectives construct positive pairs through augmentations of the same
image. Multi-sequence MRI contains a more domain-specific relationship: different
acquisitions can depict the same anatomical level even when their intensity pattern,
orientation and diagnostic role differ. DICOM patient-space geometry provides the
physical correspondence required to identify those pairs. Prior spinal
self-supervision demonstrates the value of metadata-derived or longitudinal
relationships \cite{jamaludin2017selfsupervised}, but anatomical cross-sequence
correspondence has not been established as the supervision signal for the target
formulation studied here.

The problem is not solved simply by forcing all sequence representations to become
identical. T1 and T2 contrasts contain sequence-specific diagnostic information, so
excessive invariance could improve embedding alignment while damaging grading. The
study must therefore separate the self-supervised projection space from downstream
features and test explicitly whether sequence-specific information survives
pretraining.

\subsection{Problem 4: Internal Accuracy Does Not Establish Transportability or Clinical Error Alignment}
\label{sec:intro-problem-transfer}

The strongest evidence against uncritical internal benchmarking is the repeated
performance decline under external validation
\cite{compte2023are,mcsweeney2023external,zhang2023deep}. Scanner manufacturer,
field strength, resolution, axial prescription, preprocessing and patient mix can
all change. More subtly, the reference standard may change: a public benchmark may
use structured expert annotation while a hospital's routine reports reflect local
reporting habits and may omit findings that were not considered clinically salient.
An external performance drop can therefore combine imaging-domain shift,
localisation failure and reference-standard shift.

A defensible transfer study must measure zero-shot performance before local tuning,
keep the local test set fixed, separate it from the adaptation pool, and harmonise
the local reference standard as far as resources allow. The present work therefore
uses the Rizgary Hospital cohort for external evaluation rather than for iterative
architecture development. The principal local comparison is restricted to central
canal stenosis, for which a defensible overlap with the public benchmark can be
created; local disc bulge, protrusion and extrusion remain a separate morphology
task rather than being relabelled as stenosis.

The error objective also requires care. Severity grades are ordinal, and distant
under-grading is clinically different from adjacent disagreement. Yet the literature
does not establish that ordinal losses automatically improve classification:
Niemeyer et al. found several ordinal formulations unable to outperform standard
cross-entropy for Pfirrmann grading \cite{niemeyer2021a}. A recent preprint has begun to evaluate distance-aware ordinal assessment on
LumbarDISC \cite{trinc2026beyond}; however, the combination of asymmetric clinical
cost with target-level relational modelling appears unaddressed in the literature
reviewed for this thesis.
At the same time, reducing
Severe$\rightarrow$Normal/Mild errors by simply shifting the model toward predicting
Severe would create an unacceptable false-positive burden. The fourth problem is
therefore to evaluate error direction, calibration and selective prediction without
converting radiological grading into an unsupported claim of treatment utility.

\subsection{Integrated Problem Statement}
\label{sec:intro-integrated-problem}

Taken together, the research problem can be stated as follows:

\begin{quote}
Current automated lumbar MRI grading systems can localise anatomy and predict
per-level disease severity with increasingly strong internal performance, but the
prevailing formulation does not fully exploit the relational structure of the
25-target output space, does not explicitly learn target- and quality-dependent
sequence allocation, underuses DICOM-defined cross-sequence anatomical
correspondence as a self-supervised signal, and provides limited evidence on how
these design choices behave under true institutional transfer. A controlled
methodology is therefore required to determine whether explicitly modelling these
properties improves grading, robustness, label efficiency and portability beyond
strong matched baselines, while preserving local evidence and respecting the
uncertainty of the clinical reference standard.
\end{quote}

\laterinsert{If implementation reveals a previously unrecognised but reproducible
problem that materially changes the framing---for example a systematic DICOM
geometry failure mode across a substantial cohort---add a short factual paragraph
here and document the protocol amendment. Do not add a new problem merely because a
proposed contribution underperformed.}

% ===========================================================================
\section{Research Gap and Novelty Boundary}
\label{sec:intro-gap}

Chapter~\ref{ch:litreview} develops the evidence in detail; the purpose here is to
state the gap concisely enough to define the thesis.

\subsection{Gap A: Explicit Relational Modelling of Disease--Anatomy Targets}
\label{sec:intro-gap-graph}

The reviewed literature contains multi-task systems, inter-level attention,
anatomical priors and one lumbar graph model, but no reviewed study constructs the
benchmark's level--condition--laterality targets as a heterogeneous graph with typed
relations for adjacent levels, within-level condition coexistence and bilateral
symmetry. Baur et al. use a graph within one reconstructed disc
\cite{baur2024automated}; Nguyen et al. introduce level-specific priors and
cross-sequence attention without target-to-target graph communication
\cite{nguyen2026crossspine}; Chai et al. demonstrate anatomy-guided contextual
multi-level modelling without the disease-target graph formulation proposed here
\cite{chai2026anatomy}. The intended contribution is therefore not ``graph neural
networks for lumbar MRI'' but a controlled test of an explicit typed graph over the
prediction targets themselves.

\subsection{Gap B: Anatomically Aligned Self-Supervision and Disease-Conditioned Evidence Routing}
\label{sec:intro-gap-sequence}

The literature establishes both the complementarity of MRI sequences and the value
of self-supervised pretraining, but these strands remain only partially connected in
lumbar grading. Anatomical correspondence across sequences is usually treated as a
preprocessing requirement for fusion rather than as supervision. At the same time,
multisequence systems generally learn a common fusion mechanism rather than a
condition- and quality-dependent allocation that is deliberately challenged under
missing and corrupted modalities. The thesis therefore tests two linked ideas: that
DICOM geometry can define cross-sequence positive relationships, and that a single
model can use the resulting features through a target-conditioned router whose
failure modes are directly measurable.

\subsection{Gap C: Quantified Transfer and Adaptation Under a Harmonised External Task}
\label{sec:intro-gap-transfer}

Existing external validations establish that transport is difficult, but they
usually report one transported model at one point in the adaptation process---often
with no adaptation at all. The operational question for a receiving hospital is
instead how performance changes before adaptation and how quickly it recovers as
local labels are supplied. The literature reviewed for this thesis contains no
Middle Eastern evaluation of the same scope and no quantified parameter-efficient
annotation curve for this lumbar target formulation. The gap is therefore not
simply ``lack of external validation''; it is lack of a staged zero-shot-to-few-shot
transfer experiment with fixed local testing, controlled annotation budgets and
explicit reference-standard harmonisation.

\subsection{Supporting Gap: Directional Error and Calibrated Uncertainty}
\label{sec:intro-gap-error}

Ordinal severity and asymmetric error are clinically meaningful, but ordinal
objectives have not consistently outperformed categorical baselines
\cite{niemeyer2021a}. Calibration is also rarely reported in lumbar MRI despite the
known overconfidence of modern neural networks \cite{guo2017calibration}. These are
therefore treated as supporting questions rather than headline novelty. A
cost-sensitive objective is retained only if it reduces clinically consequential
under-grading without creating indiscriminate Severe false positives, and selective
prediction is evaluated as a model behaviour rather than proof of a validated
clinical workflow.

The novelty claim throughout the thesis will remain appropriately bounded. The
phrasing ``to our knowledge'' or ``in the reviewed literature'' is used where
necessary, and the final claim will be updated against the literature available at
the time of submission. A new paper that adopts one element of the proposed system
does not invalidate the PhD if the controlled experimental question remains
distinct; it does, however, change how originality must be described.

\laterinsert{Immediately before thesis submission, update this novelty boundary
against publications appearing after the current literature-search cut-off. Record
which claims remain original and which should be reframed as independent validation,
extension or comparative evidence.}

% ===========================================================================
\section{Research Aim}
\label{sec:intro-aim}

The overall aim of this PhD is:

\begin{quote}
\emph{to determine whether explicitly modelling anatomical relationships among
lumbar disease targets, anatomically aligned evidence across MRI sequences, and
target-dependent sequence availability improves the accuracy, robustness, label
efficiency and cross-institutional portability of automated lumbar spine MRI
severity grading, relative to strong non-relational and fixed-fusion baselines.}
\end{quote}

The aim is intentionally phrased as a determination rather than a promise of
superiority. The thesis succeeds scientifically if it establishes under controlled
conditions where these mechanisms help, where they fail, and what their limits are.

% ===========================================================================
\section{Research Objectives}
\label{sec:intro-objectives}

The aim is operationalised through the following objectives.

\begin{enumerate}
  \item To construct a reproducible patient-level lumbar MRI analysis pipeline that
        preserves DICOM patient-space geometry, performs reliable level
        localisation, and supplies identical anatomically defined regions of
        interest to all grading baselines and proposed models.

  \item To establish competitive independent and contextual grading baselines on
        the public LumbarDISC development data, including fixed fusion,
        cross-attention and ordered inter-level context where reproducibility
        permits, so that improvements cannot be attributed to a weak comparator.

  \item To develop and evaluate anatomically aligned cross-sequence
        self-supervision in which positive relationships are defined by the same
        patient and lumbar level across MRI sequences using DICOM geometry.

  \item To quantify whether this anatomical self-supervision improves supervised
        grading under 10\%, 25\%, 50\% and 100\% label availability and whether it
        preserves sequence-specific diagnostic information.

  \item To develop and evaluate disease-conditioned adaptive sequence routing that
        learns patient- and target-dependent weights over available MRI evidence
        while explicitly representing sequence availability and acquisition quality.

  \item To test the router under both missing-sequence and corrupted-sequence
        conditions, and to compare its behaviour with concatenation, averaging and
        ordinary cross-attention under matched inputs and training budgets.

  \item To construct a heterogeneous disease--anatomy graph over the 25
        level--condition--laterality targets, with typed adjacent-level,
        same-level cross-condition and bilateral relations, and to compare it with
        independent heads, an ordered level transformer, a homogeneous graph,
        ungated message passing and shuffled topology.

  \item To test whether relational reasoning introduces adjacent-level false
        positives in naturally isolated lesions and to retain the graph only where
        any performance or robustness benefit is not purchased by clinically
        incoherent information propagation.

  \item To compare categorical, ordinal and clinically asymmetric cost-sensitive
        grading objectives, together with probability calibration and uncertainty
        analysis, while explicitly reporting both severe under-grading and severe
        over-grading.

  \item To quantify zero-shot transfer from the multinational public benchmark to
        the independent Rizgary Hospital cohort on the shared central-canal grading
        task, separating localisation, image-domain and reference-standard effects
        as far as the retrospective data allow.

  \item To estimate an annotation--performance recovery curve using controlled
        local adaptation sizes $N\in\{10,25,50,100\}$ and to compare where a limited
        trainable-parameter budget is most effectively placed.

  \item To report all principal comparisons with patient-level uncertainty,
        repeated training where feasible, contribution-specific ablations, negative
        results and explicit methodological limitations.
\end{enumerate}

These objectives define a staged experiment rather than a monolithic architecture.
A component is not retained because it appears in the proposal; it is retained only
if the corresponding experiment demonstrates a reproducible contribution or a
meaningful robustness, calibration or interpretability advantage.

% ===========================================================================
\section{Research Questions}
\label{sec:intro-rqs}

The research questions are fixed prospectively and correspond directly to the
methodological tests in Chapter~\ref{ch:methodology}.

\begin{description}
  \item[RQ1 --- Relational disease modelling.] Does representing lumbar disease
  targets as a typed heterogeneous graph improve grading over independent target
  heads, an ordered five-level transformer and a homogeneous graph, and are any
  gains concentrated where single-target evidence is weakest?

  \item[RQ2 --- Anatomically aligned self-supervision.] Does cross-sequence
  pretraining based on DICOM-defined anatomical correspondence improve grading,
  label efficiency and transfer compared with ImageNet initialisation, generic
  medical pretraining and ordinary augmentation-based contrastive learning?

  \item[RQ3 --- Disease-conditioned modality routing.] Can one model learn
  target- and patient-dependent sequence weights while remaining robust to both
  missing and quality-degraded MRI sequences, without requiring a separately
  trained network for every combination of sagittal T1, sagittal T2/STIR and axial
  T2?

  \item[RQ4 --- Cross-institutional transfer.] What is the zero-shot degradation
  when a benchmark-trained model is applied to a previously unseen Middle Eastern
  clinical cohort, how much of that degradation is attributable to localisation
  versus grading/reference-standard shift, and how does performance recover as the
  amount of local annotation available for adaptation increases?

  \item[RQ5 --- Clinically asymmetric error and uncertainty.] Do ordinal and
  cost-sensitive objectives reduce clinically consequential distant errors,
  particularly Severe$\rightarrow$Normal/Mild, and can calibrated uncertainty
  support selective prediction without overstating clinical safety?
\end{description}

These questions are intentionally answerable by negative as well as positive
results. If a typed graph performs no better than a homogeneous graph, RQ1 is still
answered. If anatomical self-supervision improves label efficiency but not external
transfer, RQ2 yields a more limited but still informative result. This design avoids
defining doctoral success as the requirement that every proposed component win.

% ===========================================================================
\section{Research Hypotheses}
\label{sec:intro-hypotheses}

The study uses directional hypotheses where the literature justifies a prediction,
but avoids arbitrary target accuracies that would have no defensible basis before
implementation.

\begin{description}
  \item[H1 --- Anatomical self-supervision.] Anatomical cross-sequence
  self-supervision will improve label-efficient grading and cross-institutional
  robustness relative to generic pretraining, particularly when only a fraction of
  benchmark labels is used.

  \item[H2 --- Adaptive routing.] Disease-conditioned routing with explicit
  availability masking, modality dropout and quality-aware gating will degrade less
  under missing-sequence and corruption experiments than fixed concatenation or
  ordinary cross-attention trained only on clean complete studies.

  \item[H3 --- Heterogeneous relational structure.] A typed heterogeneous graph
  will outperform independent target heads and a homogeneous graph if the proposed
  anatomical edge types carry useful information; performance under a shuffled-edge
  control should deteriorate if the topology itself matters.

  \item[H4 --- External degradation.] Zero-shot performance on the Rizgary Hospital
  cohort will be lower than internal benchmark performance because of scanner,
  protocol, population, localisation and reference-standard shift; the magnitude
  and pattern of that degradation are treated as findings rather than as failure of
  the research programme.

  \item[H5 --- Annotation-dependent recovery.] Local adaptation will improve
  performance as annotation increases, but no predetermined percentage of recovery
  is required. The shape of the recovery curve and the location in the network at
  which limited adaptation is most effective are empirical outcomes.

  \item[H6 --- Asymmetric grading error.] Cost-sensitive and ordinal objectives may
  reduce distant Severe$\rightarrow$Normal/Mild errors even if they do not improve
  aggregate macro-F1; any reduction must be evaluated against false-positive Severe
  predictions and against the cross-entropy baseline
  \cite{niemeyer2021a}.
\end{description}

The hypotheses will not be rewritten after inspection of the held-out public or
Rizgary test results. If an implementation decision changes the operational test of
a hypothesis, that change will be entered into the protocol changelog with its date
and rationale.

% ===========================================================================
\section{Conceptual Framework}
\label{sec:intro-conceptual}

The thesis can be understood as a progression from \emph{local evidence} to
\emph{structured evidence} and finally to \emph{transported evidence}.

At the first level, a model receives anatomically localised image evidence for a
single level and target. This is the baseline condition: if local evidence alone is
sufficient, more complex mechanisms should not be retained. At the second level,
the model is permitted to combine complementary MRI sequences, but the combination
must respect anatomical correspondence and sequence availability. Anatomical
self-supervision asks whether the same structure under different acquisitions can
produce a more label-efficient representation, while disease-conditioned routing
asks how strongly each sequence should contribute to a given target.

At the third level, a target is no longer treated as an isolated endpoint. Its
representation becomes a node in a typed graph whose relations correspond to three
known structures of the lumbar grading problem: the ordered chain of levels,
coexisting compartments at a level, and bilateral pairing. Message passing is soft
rather than deterministic; no rule states that disease at one level forces disease
at its neighbour. The graph is therefore a hypothesis about useful context, not a
hard-coded disease progression model.

At the fourth level, predictions are evaluated as ordinal and probabilistic outputs.
The model is judged not only by aggregate discrimination but by error distance,
Severe-class performance, calibration and the false-positive cost of any
cost-sensitive objective. This stage asks whether a statistically improved model is
also better aligned with the shape of clinically consequential grading error.

The final level is institutional transport. The representation, router, graph and
output head developed on public data are frozen before external evaluation. The
Rizgary Hospital cohort then asks a different question from internal validation:
what remains true when the acquisition and reporting environment changes? Only
after zero-shot performance has been measured is local adaptation permitted. This
ordering prevents the local test cohort from becoming an informal development set.

Conceptually, the thesis therefore tests the proposition

\begin{equation}
\text{robust grading}
\;=\;
f(\text{local anatomy},\;\text{cross-sequence evidence},\;
\text{target relations},\;\text{uncertainty},\;\text{domain}),
\label{eq:intro-conceptual}
\end{equation}

where the equation is explanatory rather than a literal computational model. Each
term corresponds to a controlled experiment, and no term is assumed necessary until
its ablation supports that conclusion.

% ===========================================================================
\section{Significance of the Study}
\label{sec:intro-significance}

\subsection{Scientific Significance}
\label{sec:intro-significance-scientific}

The principal scientific significance lies in changing the unit of reasoning from
an isolated image label to a structured set of disease targets. Many improvements in
medical imaging can be produced by additional parameters, larger backbones or more
aggressive tuning. Such gains are difficult to interpret mechanistically. By
contrast, this thesis attaches each proposed contribution to a specific hypothesis
and a destructive control: anatomical edges are compared with homogeneous and
shuffled topology; adaptive routing is compared with fixed fusion and challenged by
missing and corrupted modalities; anatomical self-supervision is compared with
generic pretraining and tested for over-invariance. The resulting evidence can show
not merely whether the full model scored higher, but \emph{why} it did or did not.

This matters because recent lumbar literature already demonstrates sophisticated
contextual models \cite{chai2026anatomy,batra2025mscan,nguyen2026crossspine}. The
research contribution cannot therefore be the generic use of attention, graphs or
multiple sequences. Its value lies in isolating the anatomical assumptions and
measuring them against alternatives under one controlled pipeline.

\subsection{Methodological Significance}
\label{sec:intro-significance-methodological}

The study places unusually strong emphasis on validity controls. Patient identity,
not slice or target, defines the data split; self-supervised learning is prohibited
from seeing held-out patients; localisation is evaluated separately from grading;
zero-shot transfer is measured before adaptation; the external reference standard is
kept distinct from routine report extraction; and every claimed model contribution
requires an ablation. These choices respond directly to recurring weaknesses in the
medical-imaging literature, where leakage, non-comparable splits and weak external
validation can produce optimistic performance estimates
\cite{compte2023are,wang2024machine}.

The transfer design is methodologically significant for a second reason. The local
hospital cohort is not used simply as another test set. The experiment attempts to
disentangle localisation failure, image-domain shift and reference-standard shift,
then measures recovery as local labels are added. This converts ``generalisation''
from a binary claim into a measurable process.

\subsection{Translational and Regional Significance}
\label{sec:intro-significance-regional}

The regional significance follows from the use of an independent Iraqi hospital
cohort that was not curated to mirror the public benchmark. Models developed on
large international datasets are often assumed to transfer into settings whose
scanner mix, acquisition practice, disease prevalence and reporting conventions were
not represented in development. Testing that assumption is particularly important
for institutions that do not have the resources to reproduce a large annotation
campaign locally. If useful performance can be recovered from tens rather than
hundreds of local annotations, that is a practical finding; if not, the negative
result is equally important because it quantifies the annotation burden required for
responsible transfer.

The study does not claim that one external hospital represents all of Iraq or the
Middle East. Rizgary Hospital provides evidence of transport to one genuinely
independent setting. Regional generalisation would require additional hospitals and
is therefore a future study rather than an inference drawn from one centre.

\subsection{Clinical Significance and its Boundary}
\label{sec:intro-significance-clinical}

The intended clinical significance is improved technical support for structured MRI
grading: consistency, reproducibility, calibrated confidence and explicit handling
of incomplete imaging. The thesis does not claim that this will improve treatment
outcomes, referral decisions or radiologist productivity. Those claims require a
prospective reader-assistance or impact study, as demonstrated by work that measures
reader performance directly \cite{lim2022improved}. The present study therefore
uses language such as ``decision-support potential'' and ``technical portability'',
not ``clinical benefit''.

% Written after Chapter 4. States only what the executed experiments support.
Of the benefits anticipated above, the executed experiments support
\emph{discrimination} and, within it, the reduction of clinically distant
errors. Quadratic weighted kappa improves by $+0.0177$ over a single-sequence
baseline, Severe$\rightarrow$Normal confusions fall from $5.7\%$ to $3.8\%$ and
recall on Severe targets rises from $61.1\%$ to $63.1\%$, the last two
attributable to the ordinal and cost-sensitive objective rather than to any
proposed structural mechanism.

Robustness to missing sequences improves, but the improvement is attributable to
modality dropout rather than to disease-conditioned routing, which changes
sequence reliance by an amount indistinguishable from zero. Label efficiency and
cross-institutional transfer were not evaluated and no claim is made for either.
Calibration is reported but was not the object of a confirmatory comparison.

No inference is drawn from these figures about patient benefit. The study
measures agreement with a reference standard on a retrospective benchmark; it
does not measure any clinical outcome.

% ===========================================================================
\section{Proposed Contributions}
\label{sec:intro-contributions}

At the protocol stage, the contributions below collectively form the proposed \textbf{AMOG-Net} framework and are \emph{proposed contributions}.
They become demonstrated thesis contributions only if the experiments support them.

\subsection{Contribution 1: Anatomically Aligned Cross-Sequence Self-Supervision}
\label{sec:intro-contrib-ssl}

The first proposed contribution uses DICOM-defined patient-space correspondence to
construct positive relationships between different MRI sequences depicting the same
patient and lumbar level. The objective is not to invent a new generic contrastive
loss; established contrastive or non-contrastive objectives may be used. The
contribution under test is the anatomical definition of correspondence and its value
for label efficiency, grading and transport. A dedicated preservation experiment
checks whether alignment erases sequence-specific diagnostic information.

\subsection{Contribution 2: Disease-Conditioned and Quality-Aware Sequence Routing}
\label{sec:intro-contrib-routing}

The second proposed contribution allows one study to be fused differently for
different disease targets. Routing depends on target identity, patient features,
sequence availability and quality evidence, and is evaluated under genuine or
simulated missingness and corruption. The contribution is not the existence of a
softmax gate or mixture-of-experts mechanism; it is the target-conditioned allocation
of clinically complementary evidence and the explicit robustness evaluation that
accompanies it.

\subsection{Contribution 3: Heterogeneous Disease--Anatomy Graph Reasoning}
\label{sec:intro-contrib-graph}

The third proposed contribution represents the 25 benchmark outputs as nodes in a
heterogeneous graph. Typed edges encode adjacent-level, within-level cross-condition
and bilateral relations. The graph includes a mechanism for preserving local
evidence when neighbour messages are unhelpful and is evaluated against independent,
ordered-transformer, homogeneous, ungated and shuffled-topology controls. A focal
isolated-lesion test determines whether relational reasoning improves uncertain
predictions without artificially spreading disease to adjacent normal levels.

\subsection{Contribution 4: A Structured Cross-Institutional Transfer Evaluation}
\label{sec:intro-contrib-transfer}

The translational contribution is an experimental design rather than a new network
layer. The selected public-data model is evaluated zero-shot on a fixed Rizgary
Hospital test set before any local adaptation. A separate local adaptation pool is
then sampled at controlled annotation sizes and used to compare constrained update
strategies. The analysis is strengthened by a verified-localisation control and,
where resources permit, harmonised independent reader grading of central canal
stenosis. This design estimates both the cost of transport and the cost of recovery.

\subsection{Supporting Contributions}
\label{sec:intro-contrib-supporting}

Ordinal and cost-sensitive objectives, calibration, uncertainty estimation and
parameter-efficient adaptation support the principal contributions but are not
claimed as original methods. Their value lies in testing whether the proposed
structured model behaves sensibly under clinically relevant error definitions and
limited local annotation. They are retained only when their contribution is
measurable.

% --------------------------------------------------------------------------
% Written at final submission, replacing the [LATER INTEGRATION] placeholder.
% Only components supported by the locked experiments are claimed. Components
% that failed are reported below as negative findings and remain in the
% research history and in the ablation ladder.
% --------------------------------------------------------------------------
\subsection{Demonstrated Contributions}
\label{sec:intro-contrib-demonstrated}

The experiments reported in Chapter~\ref{ch:results} support three
contributions. They are not the three proposed above, and the difference is
itself part of what the thesis contributes.

\paragraph{D1 --- A verified anatomy-aware multi-sequence grading system.}
AMOG-Net improves quadratic weighted kappa over a single-sequence baseline by
$+0.0177$ (95\% CI $[+0.0097, +0.0260]$), on every one of seven training seeds
and surviving false discovery rate correction across the pre-specified
comparison family. Two further comparisons also survive: typed heterogeneous
edges and the ordinal, cost-sensitive objective. The system, its frozen data
partition and its behavioural test suite are released.

\paragraph{D2 --- A controlled ablation protocol for anatomical priors.} Each
claimed mechanism is tested against a control matched in capacity rather than
against its own absence: a degree-preserving edge shuffle for graph topology, a
router-free fixed-fusion arm for sequence routing, an architecture-matched
untrained network as an attribution floor, and a label-shuffled model as a
pipeline-level negative control. Section~\ref{sec:results-threats} documents
three analyses that produced incorrect positive conclusions before the
corresponding control was introduced, in each case in the direction of the
hypothesis under test.

\paragraph{D3 --- Evidence that anatomical structural priors do not improve
grading at this data scale, with a mechanism.} Neither anatomically aligned
cross-sequence self-supervision nor disease-conditioned routing produces a
detectable improvement, and anatomically correct graph topology does not
outperform a degree-preserving random shuffle. These are bounds rather than
absences: the paired intervals constrain the three mechanisms to at most
$1.32\%$, $1.55\%$ and $1.05\%$ of baseline performance respectively. Gradient-weighted class activation
mapping shows why: the convolutional encoder already concentrates $2.4$ times the
chance share of its attention on the annotated target before any prior is added.
The localisation these priors were designed to supply is largely accomplished
upstream of them.

\subsection{Proposed Contributions Not Supported}
\label{sec:intro-contrib-negative}

Contribution~1 (anatomical cross-sequence self-supervision) is rejected on
accuracy, on cross-sequence robustness under controlled input ablation, and on
attention concentration. The pretext task is learnable and was learned ---
validation InfoNCE $1.0771$ against a chance value of $3.4657$ --- but the
representation confers no measurable downstream benefit.

Contribution~2 (disease-conditioned routing) divides. The router does learn
clinically sensible target-dependent sequence weights, reproducing the expected
allocation in all fifteen runs in which a gate is present. It does not improve
grading ($+0.0053$ QWK, five seeds of seven, bounded at $1.55\%$ of baseline)
and it does not improve robustness to missing sequences.
Its allocation tracks the annotation structure of the benchmark rather than
creating a dependence of its own.

Contribution~3 (heterogeneous disease--anatomy graph) is partly supported. Typed
relational structure improves grading over a homogeneous graph ($+0.0093$ QWK on
every one of seven seeds, surviving correction at $p_{\mathrm{FDR}} = 0.009$);
the anatomical identity of the relations does not. The resulting
claim is narrower than proposed and more specific than the original: relational
typing carries information, anatomical adjacency does not.

Contribution~4 (cross-institutional transfer evaluation) was not executed. Cohort
preparation is complete and is reported in Section~\ref{sec:results-rq4};
adjudication of conflicting reference labels and a correct de-identification
procedure remain outstanding. It is retained here as unexecuted rather than
withdrawn.

% ===========================================================================
\section{Scope and Delimitations}
\label{sec:intro-scope}

Clear boundaries are necessary because lumbar MRI supports a much broader clinical
and computational agenda than one PhD can test.

\subsection{Clinical Scope}
\label{sec:intro-scope-clinical}

The primary public task is degenerative lumbar disease severity grading at five
levels from L1--L2 through L5--S1. On LumbarDISC, the full target space contains
central canal stenosis, left and right neural foraminal narrowing, and left and
right subarticular stenosis, each graded as Normal/Mild, Moderate or Severe
\cite{richards2026the}. The study does not attempt to diagnose every cause of low
back pain, infer symptoms, determine surgical candidacy or predict long-term patient
outcomes.

The primary local transfer task is narrower. Rizgary Hospital reports support a
defensible central-canal comparison more readily than the complete benchmark schema.
The confirmatory external analysis is therefore restricted to central canal stenosis
at the five lumbar levels unless new expert grading expands the local reference
standard. Disc bulge, protrusion and extrusion are retained as a separate local
morphology task and are not heuristically converted into stenosis grades.

\subsection{Imaging Scope}
\label{sec:intro-scope-imaging}

The principal imaging inputs are sagittal T1, sagittal T2/STIR and axial T2 MRI.
Other sequences may appear in clinical studies but do not define the primary model
unless a justified later protocol amendment is made. DICOM geometry is preserved
because cross-plane correspondence is part of the scientific design. PNG or other
rendered formats may be used for visual inspection but do not replace the geometric
metadata required for alignment.

\subsection{Methodological Scope}
\label{sec:intro-scope-method}

Anatomical localisation and segmentation are enabling technologies. They are
measured and controlled because a poor crop can invalidate grading, but a new
localisation architecture is not the headline doctoral contribution. Likewise,
backbone selection is controlled rather than treated as a novelty competition. The
research may use ResNet-, EfficientNet-, Swin- or other contemporary encoders, but
the primary backbone is frozen before the main contribution-specific ablations.

The graph topology is defined anatomically rather than discovered freely from the
test data. Adjacent levels, same-level conditions and bilateral counterparts form the
pre-specified relation families. Learning the globally optimal topology is a
different research problem and is excluded from the primary scope.

\subsection{Data and Institutional Scope}
\label{sec:intro-scope-data}

Model development uses the public LumbarDISC benchmark and, where permitted, SPIDER
for anatomical localisation or segmentation support. The locally held labelled
competition subset contains approximately 1,975 studies at the protocol stage;
\torecord{final public development, validation and held-out counts after quality
control}. The local project inventory currently identifies 299 narrative reports
and 294 candidate matched multi-sequence cases, while a wider raw DICOM inventory
contains additional studies. These numbers are not presented as the final cohort
until one-to-one reconciliation is complete:
\torecord{final Rizgary case-flow counts and reasons for exclusion}.

Before analytical use of the local cohort, the Rizgary data-preparation phase will
follow the approved institutional governance protocol. A controlled research copy will
be de-identified before access by the model-development team, and the permitted storage
and compute environment will be verified before local processing begins. These safeguards
precede, rather than form part of, model development.

Rizgary Hospital is a single external institution. The study can therefore support a
claim of transport to that independent setting, not universal Middle Eastern
performance. Additional regional hospitals would be required for a multicentre
regional generalisation claim.

\subsection{Ethical and Governance Scope}
\label{sec:intro-scope-ethics}

The local component is retrospective and does not alter patient care. Formal use of
the local data remains conditional on the relevant institutional permissions:
\toconfirm{hospital/ethics approval reference, date and responsible institution};
\toconfirm{retrospective consent or waiver basis}; and
\toconfirm{permitted compute environment, including whether external cloud/GPU
processing is authorised}. Identifiable raw DICOM is not supplied to students or
external compute; a controlled de-identified research copy is required before local
analysis.

\subsection{Evaluation Scope}
\label{sec:intro-scope-evaluation}

The thesis evaluates technical grading performance, robustness, calibration,
computational cost and cross-institutional portability. The unit of statistical
independence is the patient. It does not claim prospective clinical impact,
radiologist replacement or treatment benefit. Reader-assistance evaluation is a
logical future study if the technical system proves sufficiently reliable.

% ===========================================================================
\section{Assumptions and Operational Definitions}
\label{sec:intro-definitions}

The following terms are used consistently throughout the thesis.

\begin{longtable}{p{3.6cm}p{10.0cm}}
\toprule
\textbf{Term} & \textbf{Operational meaning in this thesis}\\
\midrule
\endfirsthead
\toprule
\textbf{Term} & \textbf{Operational meaning in this thesis}\\
\midrule
\endhead
Lumbar level & One of L1--L2, L2--L3, L3--L4, L4--L5 or L5--S1.\\
Target & A level--condition--laterality output on LumbarDISC; 25 targets per study at full benchmark scope.\\
Severity grade & Normal/Mild, Moderate or Severe, encoded as an ordered three-class outcome.\\
Local evidence & Image-derived representation for one target before graph message passing.\\
Anatomical correspondence & Physical correspondence between acquisitions defined through DICOM patient-space geometry and lumbar localisation, not through image appearance alone.\\
Sequence routing & Learned allocation of weight over available MRI sequence features for a specific patient and target.\\
Missing modality & An expected MRI sequence is absent and explicitly unavailable to the model.\\
Degraded modality & A sequence is present but technically compromised, for example by motion, truncation, noise, slice loss or poor coverage.\\
Heterogeneous graph & A graph whose nodes are disease--anatomy targets and whose edges belong to clinically defined relation types.\\
Zero-shot transfer & Evaluation of the frozen public-data model on the fixed local test set before local fine-tuning, threshold selection or local calibration.\\
Few-shot adaptation & Model adaptation using only the separate local adaptation pool under pre-specified annotation budgets.\\
Reference-standard shift & Difference between the process used to generate labels in the benchmark and the process used to establish the local external reference.\\
Selective prediction & Model abstention or referral based on uncertainty; this is a technical behaviour, not evidence of a validated clinical workflow.\\
\bottomrule
\end{longtable}

These definitions are intentionally narrower than ordinary-language usage. In
particular, ``external validation'' does not mean that a model has been proven safe
for deployment, and ``uncertainty'' does not mean that the network has measured a
patient's clinical risk.

% ===========================================================================
\section{Overview of the Research Design}
\label{sec:intro-design-overview}

The research is a retrospective, multi-dataset medical-imaging methods study. Its
logic is cumulative, but the conclusions remain component-specific.

\begin{enumerate}
  \item \textbf{Data and anatomy foundation.} Public data are inventoried,
        partitioned by patient, reconstructed in DICOM patient space, anatomically
        localised and converted into disease-specific 2.5D regions of interest.

  \item \textbf{Strong baseline establishment.} An independent ROI classifier and
        competitive multiview/context baselines establish what can be achieved
        without the proposed relational formulation.

  \item \textbf{Representation experiment.} Anatomically aligned cross-sequence
        self-supervision is evaluated against generic and medical pretraining under
        several label fractions and sequence-specific preservation probes.

  \item \textbf{Fusion experiment.} Disease-conditioned sequence routing is tested
        against fixed and attention-based fusion under complete, missing and
        corrupted sequence conditions.

  \item \textbf{Relational experiment.} Homogeneous and heterogeneous target graphs
        are compared with independent and ordered-transformer controls, including
        shuffled topology and isolated-lesion stress testing.

  \item \textbf{Output and uncertainty experiment.} Categorical, ordinal and
        cost-sensitive objectives are evaluated together with calibration and
        uncertainty, with both under-grading and over-grading reported.

  \item \textbf{External transport experiment.} The selected public-data model is
        frozen and evaluated zero-shot on the fixed Rizgary Hospital test set. A
        verified-ROI control and harmonised reader reference are used where feasible
        to interpret the source of degradation.

  \item \textbf{Adaptation experiment.} Controlled local annotation budgets are
        sampled from a separate adaptation pool to estimate how quickly performance
        recovers and where a limited trainable-parameter budget is most effective.
\end{enumerate}

This order deliberately prevents the final architecture from becoming an indivisible
engineering object. Each stage introduces one assumption that can be rejected
without invalidating the remainder of the thesis. The final model is therefore the
simplest configuration supported by the evidence, not necessarily the configuration
containing every proposed module.

% ===========================================================================
\section{Expected Outcomes and Criteria for Scientific Success}
\label{sec:intro-success}

The protocol does not define success as reaching a predetermined accuracy. Such a
threshold would be difficult to justify because performance depends on target,
severity distribution, reference-standard reliability and the final patient split.
Scientific success is instead defined by whether the research questions can be
answered under controlled conditions.

For RQ1, success means establishing whether typed anatomical relations add value
beyond capacity-matched controls and whether they avoid adjacent-level
contamination. A null result is informative if the comparison is sufficiently
powered and the topology controls are valid. For RQ2, success means establishing
whether anatomical cross-sequence pairing improves label efficiency, transfer or
neither, while documenting whether sequence-specific information is preserved. For
RQ3, success means quantifying robustness under missing and degraded modalities and
determining whether adaptive routing adds value over simpler fusion. For RQ4,
success means producing an interpretable zero-shot transfer estimate and a controlled
adaptation curve, even if the external drop is large. For RQ5, success means
characterising the trade-off between severe under-grading, severe over-grading and
probability reliability rather than selecting a loss solely because it raises one
aggregate metric.

At final submission, the thesis will distinguish three categories explicitly:
\emph{supported contributions}, \emph{negative findings}, and \emph{unresolved
questions}. This is preferable to rewriting the original proposal so that only
successful components appear to have been planned.

% Written after Chapter 4. High-level answers only; metrics live in Chapter 4.
\begin{description}
  \item[RQ1 --- Relational disease modelling.] \textbf{Partly supported.} A typed
  heterogeneous graph improves grading over independent heads and over a
  homogeneous graph, on every seed. The anatomical identity of the relations does
  not contribute: an anatomically correct topology does not outperform a
  degree-preserving random shuffle.

  \item[RQ2 --- Anatomically aligned self-supervision.] \textbf{Not supported.}
  The pretext task is learnable and was learned, but the resulting representation
  yields no detectable improvement in grading, in robustness to withheld
  sequences, or in where the model attends.

  \item[RQ3 --- Disease-conditioned modality routing.] \textbf{Mechanism
  supported, benefit not.} The router reliably learns the clinically expected
  sequence weights, but a router-free model shows the same sequence dependence
  under intervention, and routing improves neither accuracy nor missing-sequence
  robustness.

  \item[RQ4 --- Cross-institutional transfer.] \textbf{Not executed.} Cohort
  preparation is complete; reference-label adjudication and de-identification
  remain outstanding.

  \item[RQ5 --- Clinically asymmetric error and uncertainty.] \textbf{Supported
  for error asymmetry.} Ordinal and cost-sensitive objectives improve both
  aggregate agreement and recall on Severe targets, on every seed, and form the
  most reliable single contributor to final performance. Selective prediction
  under calibrated uncertainty was not evaluated confirmatorily.
\end{description}

% ===========================================================================
\section{Thesis Organisation}
\label{sec:intro-organisation}

At the current protocol stage, the thesis is organised around the progression from
problem definition to evidence.

\textbf{Chapter 1 --- Introduction} defines the clinical and methodological
motivation, research problem, gaps, aim, objectives, research questions, hypotheses,
proposed contributions, significance and scope of the thesis.

\textbf{Chapter 2 --- Literature Review} examines the clinical grading systems that
define the reference labels, MRI sequence--pathology correspondence, localisation
and segmentation, backbone and contextual architectures, class imbalance and
ordinal grading, public datasets, external validation, interpretability and the
specific gaps that motivate the present study.

\textbf{Chapter 3 --- Methodology} translates those gaps into a protocol-grade
experimental design. It defines the public and local cohorts, leakage controls,
DICOM reconstruction, localisation and ROI construction, self-supervised
pretraining, adaptive routing, heterogeneous graph reasoning, grading objectives,
calibration, ablations, external validation, domain adaptation, statistical analysis
and reproducibility safeguards.

\textbf{Chapter 4 --- Results and Analysis} will report the staged E0--E8
experiments, including component-specific ablations, robustness analyses,
calibration, zero-shot external transfer, few-shot adaptation and failure analysis.
% Written after Chapter 4. Retained experiment groups; research questions unchanged.
The retained experiment groups are the eight-rung ablation ladder E0--E7 with two
matched controls on the graph configuration, a controlled input ablation testing
sequence dependence by intervention rather than by learned allocation, and a
Grad-CAM attribution analysis against an architecture-matched untrained floor.
The cross-institutional transfer group was prepared but not executed and is
reported as such. The research questions are unchanged from those fixed in
Section~\ref{sec:intro-rqs}.

\textbf{Chapter 5 --- Discussion and Conclusions} will interpret the findings
against the literature, distinguish demonstrated contributions from negative
results, state limitations, discuss implications for transport to other institutions
and identify the next studies required before any clinical-impact claim.

\toconfirm{final institutional thesis structure if the university requires results
and discussion as separate chapters, publication-based chapters, or an additional
conclusion chapter}. The scientific ordering above should remain recognisable even
if the administrative chapter numbering changes.

% ===========================================================================
\section{Rules for Completing Prospective Placeholders}
\label{sec:intro-placeholders}

Because this chapter is written before completion of implementation, several fields
are intentionally prospective. The following rules govern their later completion.

\begin{enumerate}
  \item \textbf{[TO CONFIRM]} is used only for institutional or governance facts
        that are currently unknown, such as an approval reference or permitted
        compute environment. These fields are replaced only with documented facts.

  \item \textbf{[TO RECORD AFTER IMPLEMENTATION]} is used for executed technical or
        cohort quantities that cannot responsibly be specified before the pipeline
        runs, such as final quality-controlled sample counts.

  \item \textbf{[LATER INTEGRATION]} is used for a small number of places where the
        final thesis may state the principal findings or update the novelty boundary.
        These additions may report evidence; they must not alter the prospective
        research questions or make an unsuccessful hypothesis disappear.

  \item Any substantive protocol change made after results become visible is
        documented with date, reason and affected analysis. Confirmatory and
        exploratory analyses remain distinguishable.

  \item The introduction will not be retrofitted to imply that the final retained
        architecture was inevitable. The final thesis should preserve the fact that
        ACSSL, adaptive routing, heterogeneous graph reasoning and supporting losses
        were hypotheses subjected to testing.
\end{enumerate}

These rules turn placeholders into a transparency mechanism rather than an invitation
to hindsight bias.

% ===========================================================================
\section{Chapter Summary}
\label{sec:intro-summary}

This chapter has defined the thesis as a study of structured and transferable lumbar
MRI grading rather than as a competition between image backbones. The clinical task
is inherently multi-level, multi-condition, bilateral and multi-sequence, while its
reference labels are ordinal and imperfectly reproducible. The technical literature
has established that localisation should precede grading, that sophisticated
multiview systems can perform strongly on public data, and that external performance
frequently degrades. Those achievements narrow rather than remove the research gap.

The thesis therefore asks whether three explicit structural mechanisms add value:
anatomically aligned cross-sequence self-supervision, disease-conditioned and
quality-aware sequence routing, and a typed heterogeneous graph over
level--condition--laterality targets. These mechanisms are evaluated component by
component against strong controls rather than bundled into a single architecture and
assumed beneficial. Supporting ordinal, cost-sensitive and uncertainty methods are
used to test the clinical shape of error without being advertised as standalone
novelty.

The translational component separates model development from external evaluation.
LumbarDISC supplies the public development benchmark; SPIDER may support anatomical
localisation; and the independent Rizgary Hospital cohort is reserved for zero-shot
central-canal transfer and controlled local adaptation. The local reference process
is treated explicitly as a potential source of shift rather than assumed equivalent
to the public benchmark. Clinical impact remains outside the scope until a separate
reader-assistance or prospective study is conducted.

The aim, objectives, RQ1--RQ5 and H1--H6 stated here form the prospective frame for
Chapters~\ref{ch:litreview} and~\ref{ch:methodology}. The remainder of the thesis is
therefore not required to prove that every proposed mechanism succeeds. It is
required to determine, with reproducible evidence, which assumptions survive
controlled testing and institutional transfer.

% Written at final submission.
In summary, the complete system improves on a single-sequence baseline
reproducibly and by a margin that survives correction for multiple comparisons
while remaining far too small to alter clinical practice,
and the objective function designed around clinically asymmetric error is the
most reliable single contributor to that improvement. Of the three proposed
structural mechanisms, relational typing is supported, anatomical topology and
disease-conditioned routing are not, and anatomically aligned self-supervision is
rejected on every axis on which it was tested; attribution analysis indicates
that the convolutional encoder already performs the localisation these mechanisms
were intended to supply. The principal limitations are that all results are
internal to one benchmark, that cross-institutional transfer was not executed,
and that the study is underpowered for single-component effects of the magnitude
observed.

% ===========================================================================
\printbibliography[heading=bibintoc,title={References}]

\end{document}
 from the main thesis file.
% ===========================================================================

\documentclass[12pt,a4paper]{report}

\usepackage{lmodern}
\usepackage[T1]{fontenc}
\usepackage[utf8]{inputenc}
\usepackage[english]{babel}
\usepackage{csquotes}
\usepackage{microtype}
\usepackage[margin=2.5cm]{geometry}
\usepackage{setspace}
\onehalfspacing

\usepackage{amsmath,amssymb}
\usepackage{booktabs}
\usepackage{longtable}
\usepackage{array}
\usepackage{graphicx}
\usepackage{thesisstyle}   % shared TikZ figure styles; see thesisstyle.sty
\usepackage[hidelinks]{hyperref}

\usepackage[
  backend=biber,
  style=numeric-comp,
  sorting=none,
  maxbibnames=6,
  minbibnames=3,
  giveninits=true,
  doi=true,
  url=false,
  isbn=false
]{biblatex}

% Production bibliography cleaned for thesis compilation. Citation keys are
% unchanged from the working literature inventory.
\addbibresource{lumbar_spine_mri_ai_references_clean.bib}

\newcommand{\toconfirm}[1]{\textbf{[TO CONFIRM: #1]}}
\newcommand{\torecord}[1]{\textbf{[TO RECORD AFTER IMPLEMENTATION: #1]}}
\newcommand{\laterinsert}[1]{\textbf{[LATER INTEGRATION: #1]}}

\begin{document}

\chapter{Introduction}
\label{ch:introduction}

% ===========================================================================
\section{Introduction and Thesis Position}
\label{sec:intro-position}

Lumbar degenerative disease presents an unusually demanding problem for medical
image analysis. The anatomical region is compact, but the diagnostic task is not:
one examination contains several intervertebral levels, several pathological
compartments, left--right distinctions, multiple magnetic resonance imaging (MRI)
sequences acquired in different planes, and severity labels whose boundaries are
partly subjective. A clinically meaningful automated system must therefore do more
than recognise that degeneration is present. It must identify \emph{what} is
abnormal, \emph{where} it is abnormal, \emph{how severe} the abnormality is, and
\emph{which imaging evidence} supports that judgement. It must also remain
interpretable as a measurement of radiological grading rather than be mistaken for
an autonomous clinical diagnosis.

The need for such systems arises from both disease burden and reporting workload.
Low back pain remains one of the largest contributors to disability worldwide, and
lumbar degenerative change is a frequent structural correlate examined by MRI
\cite{compte2023are,wang2024machine}. Routine interpretation is repetitive because
the reader must inspect multiple levels across sagittal and axial acquisitions,
while the severity of central canal, foraminal and subarticular narrowing is judged
against morphological criteria that are not perfectly reproducible between human
readers \cite{lurie2008reliability,schizas2010qualitative,lee2011a,lee2010a}.
Automation is attractive in this setting not because a neural network is assumed to
possess a superior biological truth, but because a reproducible system may help
standardise a high-volume structured task whose human reference process itself has
known variability.

The technical field has advanced rapidly. Early systems demonstrated that lumbar
anatomy could be localised and individual discs graded with convolutional networks;
more recent pipelines combine explicit localisation, 2.5D regions of interest,
multiview fusion, attention and large public benchmarks
\cite{lu2018deepspine,pan2021automatically,batra2025mscan,richards2026the}. The
central problem has consequently shifted. The thesis no longer asks whether a
convolutional neural network or transformer can classify a lumbar MRI in isolation.
It asks whether the \emph{structure of the grading problem itself} has been
represented faithfully enough.

This thesis takes the position that lumbar MRI grading is a structured prediction
problem with four coupled sources of information. First, disease occurs at ordered
anatomical levels rather than at unrelated image locations. Second, several
conditions coexist within one motion segment and lateral compartments occur as
paired left--right structures. Third, the diagnostic value of an MRI sequence is
target-dependent: sagittal T1, sagittal T2/STIR and axial T2 provide complementary
rather than interchangeable evidence \cite{hallinan2021deep}. Fourth, a model that
performs well on one benchmark may degrade when scanners, protocols, population and
reference-standard conventions change across institutions
\cite{mcsweeney2023external,zhang2023deep,guan2022domain}.

The research therefore evaluates three methodological contributions and one
translational validation programme. The integrated framework evaluating these components
is designated \textbf{AMOG-Net} (\textbf{A}natomy-aware \textbf{M}ulti-sequence \textbf{O}rdinal \textbf{G}raph Network). The first contribution is anatomically aligned cross-sequence
self-supervision, in which DICOM patient-space correspondence defines which views
should agree during representation learning. The second is disease-conditioned
adaptive sequence routing, in which the model learns target- and patient-dependent
weights over available MRI sequences and is explicitly tested under both missing
and degraded inputs. The third is a heterogeneous disease--anatomy graph whose
nodes represent level--condition--laterality targets and whose typed edges encode
adjacent-level, same-level cross-condition and bilateral relations. These are then
evaluated under a deliberately separate transfer programme: development occurs on
a multinational public benchmark; the selected model is frozen; zero-shot
performance is measured on an independent Rizgary Hospital cohort; and controlled
few-shot adaptation is performed only afterwards.

Several related techniques are necessary but are not claimed as doctoral novelty.
Anatomical localisation, segmentation, 2.5D region construction, conventional CNN
or transformer encoders, class-imbalance handling, ordinal loss functions,
calibration and parameter-efficient fine-tuning are established tools. They are
used to construct a fair experiment around the thesis contributions rather than
presented as original merely because they appear in the same architecture. This
novelty boundary is important in a field where anatomy-guided contextual grading,
cross-sequence attention and graph-based disc modelling have already been reported
\cite{chai2026anatomy,nguyen2026crossspine,baur2024automated}. The question is not
whether any one of those mechanisms exists elsewhere; it is whether the particular
representation tested here---typed relationships among the benchmark's disease
outputs, anatomical cross-sequence supervision, target-conditioned evidence
allocation and explicit institutional transfer---adds measurable knowledge beyond
strong matched baselines.

% Written after the confirmatory experiments were locked. Reports the outcome
% only; the problem statement, research questions and hypotheses are unchanged.
The principal empirical outcome is that the complete system improves agreement
with the reference standard over a single-sequence baseline by $+0.0177$
quadratic weighted kappa, on every one of seven training seeds and after
correction for multiple comparisons, but that the anatomical mechanisms proposed
to achieve this do not account for the improvement: attribution analysis
indicates the convolutional encoder already localises the graded structure
without them.

% ===========================================================================
\section{Background and Motivation}
\label{sec:intro-background}

\subsection{Clinical Motivation: A Repetitive but Structured Imaging Task}
\label{sec:intro-clinical-motivation}

Lumbar degeneration is not a single lesion. It is a family of related structural
changes involving the intervertebral disc, facet joints, ligamentum flavum,
vertebral endplates and the neural spaces that these structures surround. Disc
dehydration and height loss alter load distribution; annular bulging or herniation
may encroach on the canal and foramina; posterior element hypertrophy may further
reduce the space available to neural tissue. These changes are assessed at several
levels, most commonly from L1--L2 through L5--S1, and their significance depends on
which anatomical compartment is narrowed.

Three stenotic compartments are especially relevant to the benchmark task used in
this thesis. Central canal stenosis concerns the thecal sac and cauda equina,
foraminal narrowing concerns the exiting nerve roots, and subarticular or lateral
recess stenosis concerns the traversing roots. Their visual criteria are related but
not identical, and neither are the acquisitions from which they are best judged.
Central canal and lateral recess morphology depend heavily on axial T2 information,
whereas neural foraminal assessment benefits particularly from sagittal or
parasagittal T1 because perineural fat provides the relevant contrast
\cite{hallinan2021deep,schizas2010qualitative,lee2010a}. A model that treats all
sequences as equivalent therefore ignores a clinical property of the task before
learning has even begun.

The labels themselves are another source of complexity. Morphological grading
systems convert continuous anatomical variation into ordinal categories. Human
agreement is strongest for some compartments and weaker for others: the reliability
study of Lurie et al. found substantially better inter-reader agreement for central
canal stenosis than for foraminal and subarticular findings
\cite{lurie2008reliability}. This matters methodologically. If the reference standard
contains boundary uncertainty, model evaluation must distinguish gross errors from
one-grade disagreements and must interpret performance relative to the reliability
of the grading instrument. Apparent model error and reference-standard disagreement
are not always separable.

This combination---high examination volume, repeated level-wise assessment,
complementary sequence evidence and imperfect categorical boundaries---makes lumbar
MRI a suitable domain for artificial intelligence, but only if the automation is
framed correctly. The objective is not to replace full clinical interpretation.
Symptoms, neurological findings, prior imaging and management context remain outside
the information available to an image-only grading system. The narrower objective is
to make a defined radiological grading task more consistent, measurable and
portable. Reader-assistance studies support this framing: deep-learning assistance
has been shown to reduce reporting time and improve or preserve inter-reader
agreement, with the greatest gains often occurring among less experienced readers
\cite{lim2022improved}. This thesis therefore treats the system as a potential
component of decision support, not as an autonomous treatment decision-maker.

\subsection{Technical Motivation: From Classification to Structured Prediction}
\label{sec:intro-technical-motivation}

The evolution of lumbar MRI automation has progressively narrowed the unit of
prediction. Early work often asked whether a lesion or abnormality was present in an
image. Subsequent systems localised discs, graded individual levels, and shared
encoders across several pathologies. Public benchmark design has now made the task
more explicit. The RSNA Lumbar Degenerative Imaging Spine Classification resource,
published as LumbarDISC, contains 2,697 patients and 8,593 MRI series collected
across eight institutions in six countries \cite{richards2026the}. The associated
challenge requires five condition/laterality targets at five lumbar levels, producing
25 ordinal outputs per study \cite{rsna2024rsna,richards2026the}. A correct model
therefore does not simply output ``mild'', ``moderate'' or ``severe'' for an image;
it attaches one of those grades to a particular condition at a particular level and,
for lateral compartments, a particular side.

This target structure exposes a limitation in the prevailing formulation. Most
systems share image features but ultimately emit independent output heads. Shared
features allow several tasks to benefit from a common representation, but a
prediction at L4--L5 does not explicitly inform an uncertain prediction at L3--L4;
a confident left foraminal prediction does not interact structurally with its right
counterpart; and the central, foraminal and subarticular targets at one motion
segment do not exchange typed messages simply because they coexist anatomically.
The representation is therefore multi-task but not fully relational.

Graph learning provides a natural language for testing whether those relationships
carry useful information. Graph neural networks are not novel by themselves, and a
lumbar application has already reconstructed an individual disc as a point cloud and
graded it with a graph network \cite{baur2024automated}. That precedent is important
because it prevents an unjustified claim that ``using a graph in lumbar MRI'' is new.
The unresolved question is narrower: whether the \emph{prediction targets} can be
represented as nodes connected by clinically motivated relation types, and whether
those relations improve grading beyond independent heads, an ordered level
transformer and a homogeneous graph. The distinction between a graph over points
within one disc and a graph over disease--anatomy targets across the spine is the
novelty boundary this thesis tests.

A second technical motivation concerns multimodal evidence. Recent methods have
already shown that cross-view and cross-sequence attention can improve lumbar
assessment \cite{batra2025mscan,nguyen2026crossspine}, while anatomy-guided
inter-level context has also become an active direction \cite{chai2026anatomy}.
These developments mean that novelty cannot rest on ``using multiple sequences'' or
``using attention''. The remaining question is whether evidence allocation should be
\emph{conditioned on the target and the patient}, and whether one trained model can
remain operational when an expected sequence is unavailable or present but of poor
quality. A fixed fusion rule assumes the same contribution of sagittal T1, sagittal
T2/STIR and axial T2 for every target and every study. Clinical sequence--pathology
correspondence makes that assumption doubtful.

A third motivation is representation learning. Generic ImageNet pretraining is
convenient but does not exploit the particular correspondence available in MRI: the
same anatomical level can be observed under different contrasts and planes. Earlier
spinal work demonstrated that self-supervision derived from patient-level or
longitudinal information can reduce dependence on manual labels
\cite{jamaludin2017selfsupervised}, and disc-centric contrastive work continues that
line \cite{acharya2026disccentric}. DICOM geometry makes a stronger, directly
anatomical pairing possible. If sagittal and axial observations can be matched in
patient space, then the correspondence itself can serve as supervision without
pretending that the images should look visually alike.

Finally, there is a translational motivation. Systematic reviews repeatedly find
that internally impressive medical-imaging models underperform when re-evaluated
outside their development distribution \cite{compte2023are,wang2024machine}.
Lumbar-specific external validations show the same pattern, with fine-grained
grading generally less portable than localisation or detection
\cite{mcsweeney2023external,zhang2023deep}. A model intended for eventual use in
Iraq or the wider region therefore cannot be justified solely by performance on a
multinational benchmark. The model must be transported to a genuinely independent
hospital cohort, its zero-shot degradation measured before local tuning, and the
amount of local annotation required to recover performance quantified explicitly.
That annotation--performance curve is more actionable than a single statement that
``domain adaptation improved accuracy''.

\subsection{Why the Problem Should Be Addressed Now}
\label{sec:intro-why-now}

Three developments make the present study timely. First, LumbarDISC supplies a
large, multi-institutional and anatomically structured benchmark whose target schema
is rich enough to test relational modelling rather than generic image
classification \cite{richards2026the}. Second, public anatomical resources such as
SPIDER reduce the need to claim localisation or segmentation as the central
contribution \cite{vandergraaf2024lumbar}. Third, the current methodological
literature is sufficiently mature to provide strong controls: multi-view
cross-attention, anatomy-guided contextual models, contemporary CNN and transformer
backbones, self-supervised pretraining, calibration and parameter-efficient transfer
all exist as comparators or supporting tools. The research can therefore ask a more
precise question than whether a new architecture attains a high score.

The resulting opportunity is methodological rather than merely computational. The
field now has enough data and enough baseline competence to test whether explicit
clinical structure provides value over capacity alone. If a typed target graph,
anatomical self-supervision or disease-conditioned router fails to outperform its
control, that negative result is still informative. Conversely, if a gain disappears
when topology is shuffled, when a simpler transformer is matched for capacity, or
when the model is transported to another hospital, the interpretation must change.
This insistence on falsifiable components is part of the motivation for the study,
not an afterthought added to evaluation.

% ===========================================================================
\section{Problem Statement}
\label{sec:intro-problem}

Despite rapid progress in automated lumbar MRI analysis, four connected
methodological problems remain.

\subsection{Problem 1: The Output Space Is Structured but Commonly Treated as Independent}
\label{sec:intro-problem-relational}

The benchmark output space contains ordered lumbar levels, multiple conditions at
each level and bilateral counterparts. These targets arise from one connected
anatomical system, yet the dominant design predicts them independently after a
shared encoder. This separation is convenient for implementation but does not test
whether adjacent-level context, within-level co-occurrence or bilateral relations
carry predictive information. Existing graph work in lumbar MRI does not resolve
this question because its graph is constructed within a disc rather than across the
disease targets \cite{baur2024automated}. CrossSpine introduces anatomical priors
and level-specific behaviour but still does not formulate the full 25 target outputs
as a typed communication graph \cite{nguyen2026crossspine}. The first problem is
therefore not the absence of contextual modelling in general; it is the absence of a
controlled test of explicit disease--anatomy relations among the prediction targets.

A relational model also carries a specific risk: useful context can become
``information contagion''. A severe lesion at one level should not cause the model
to elevate an adjacent normal level simply because degeneration often co-occurs.
Any graph proposed as a solution must therefore preserve local visual evidence,
include non-relational and homogeneous controls, and be stress-tested on naturally
isolated lesions. Without those controls, improved aggregate performance would not
establish that the anatomical topology is beneficial.

\subsection{Problem 2: Sequence Fusion Is Usually Fixed While Evidence Is Target- and Quality-Dependent}
\label{sec:intro-problem-routing}

Sagittal T1, sagittal T2/STIR and axial T2 are complementary acquisitions. The
foramina, central canal and lateral recess are not equally well represented in each
view, and routine data may be incomplete or technically degraded. Strong recent
systems address multi-view fusion through cross-attention
\cite{batra2025mscan,nguyen2026crossspine}, but a remaining deployment question is
whether one model can dynamically allocate evidence by condition, level, patient and
sequence quality. A model trained only on complete, high-quality benchmark studies
may fail when one sequence is missing, truncated or motion-corrupted. Missingness
and corruption are also different problems: an absent sequence can be masked,
whereas a degraded sequence remains present and can attract misleading attention.
The second problem is therefore to learn adaptive sequence allocation without
turning the router into an unexamined source of instability.

\subsection{Problem 3: Cross-Sequence Correspondence Is Available but Underused as Supervision}
\label{sec:intro-problem-ssl}

Self-supervised learning can reduce dependence on costly expert labels, but common
contrastive objectives construct positive pairs through augmentations of the same
image. Multi-sequence MRI contains a more domain-specific relationship: different
acquisitions can depict the same anatomical level even when their intensity pattern,
orientation and diagnostic role differ. DICOM patient-space geometry provides the
physical correspondence required to identify those pairs. Prior spinal
self-supervision demonstrates the value of metadata-derived or longitudinal
relationships \cite{jamaludin2017selfsupervised}, but anatomical cross-sequence
correspondence has not been established as the supervision signal for the target
formulation studied here.

The problem is not solved simply by forcing all sequence representations to become
identical. T1 and T2 contrasts contain sequence-specific diagnostic information, so
excessive invariance could improve embedding alignment while damaging grading. The
study must therefore separate the self-supervised projection space from downstream
features and test explicitly whether sequence-specific information survives
pretraining.

\subsection{Problem 4: Internal Accuracy Does Not Establish Transportability or Clinical Error Alignment}
\label{sec:intro-problem-transfer}

The strongest evidence against uncritical internal benchmarking is the repeated
performance decline under external validation
\cite{compte2023are,mcsweeney2023external,zhang2023deep}. Scanner manufacturer,
field strength, resolution, axial prescription, preprocessing and patient mix can
all change. More subtly, the reference standard may change: a public benchmark may
use structured expert annotation while a hospital's routine reports reflect local
reporting habits and may omit findings that were not considered clinically salient.
An external performance drop can therefore combine imaging-domain shift,
localisation failure and reference-standard shift.

A defensible transfer study must measure zero-shot performance before local tuning,
keep the local test set fixed, separate it from the adaptation pool, and harmonise
the local reference standard as far as resources allow. The present work therefore
uses the Rizgary Hospital cohort for external evaluation rather than for iterative
architecture development. The principal local comparison is restricted to central
canal stenosis, for which a defensible overlap with the public benchmark can be
created; local disc bulge, protrusion and extrusion remain a separate morphology
task rather than being relabelled as stenosis.

The error objective also requires care. Severity grades are ordinal, and distant
under-grading is clinically different from adjacent disagreement. Yet the literature
does not establish that ordinal losses automatically improve classification:
Niemeyer et al. found several ordinal formulations unable to outperform standard
cross-entropy for Pfirrmann grading \cite{niemeyer2021a}. A recent preprint has begun to evaluate distance-aware ordinal assessment on
LumbarDISC \cite{trinc2026beyond}; however, the combination of asymmetric clinical
cost with target-level relational modelling appears unaddressed in the literature
reviewed for this thesis.
At the same time, reducing
Severe$\rightarrow$Normal/Mild errors by simply shifting the model toward predicting
Severe would create an unacceptable false-positive burden. The fourth problem is
therefore to evaluate error direction, calibration and selective prediction without
converting radiological grading into an unsupported claim of treatment utility.

\subsection{Integrated Problem Statement}
\label{sec:intro-integrated-problem}

Taken together, the research problem can be stated as follows:

\begin{quote}
Current automated lumbar MRI grading systems can localise anatomy and predict
per-level disease severity with increasingly strong internal performance, but the
prevailing formulation does not fully exploit the relational structure of the
25-target output space, does not explicitly learn target- and quality-dependent
sequence allocation, underuses DICOM-defined cross-sequence anatomical
correspondence as a self-supervised signal, and provides limited evidence on how
these design choices behave under true institutional transfer. A controlled
methodology is therefore required to determine whether explicitly modelling these
properties improves grading, robustness, label efficiency and portability beyond
strong matched baselines, while preserving local evidence and respecting the
uncertainty of the clinical reference standard.
\end{quote}

\laterinsert{If implementation reveals a previously unrecognised but reproducible
problem that materially changes the framing---for example a systematic DICOM
geometry failure mode across a substantial cohort---add a short factual paragraph
here and document the protocol amendment. Do not add a new problem merely because a
proposed contribution underperformed.}

% ===========================================================================
\section{Research Gap and Novelty Boundary}
\label{sec:intro-gap}

Chapter~\ref{ch:litreview} develops the evidence in detail; the purpose here is to
state the gap concisely enough to define the thesis.

\subsection{Gap A: Explicit Relational Modelling of Disease--Anatomy Targets}
\label{sec:intro-gap-graph}

The reviewed literature contains multi-task systems, inter-level attention,
anatomical priors and one lumbar graph model, but no reviewed study constructs the
benchmark's level--condition--laterality targets as a heterogeneous graph with typed
relations for adjacent levels, within-level condition coexistence and bilateral
symmetry. Baur et al. use a graph within one reconstructed disc
\cite{baur2024automated}; Nguyen et al. introduce level-specific priors and
cross-sequence attention without target-to-target graph communication
\cite{nguyen2026crossspine}; Chai et al. demonstrate anatomy-guided contextual
multi-level modelling without the disease-target graph formulation proposed here
\cite{chai2026anatomy}. The intended contribution is therefore not ``graph neural
networks for lumbar MRI'' but a controlled test of an explicit typed graph over the
prediction targets themselves.

\subsection{Gap B: Anatomically Aligned Self-Supervision and Disease-Conditioned Evidence Routing}
\label{sec:intro-gap-sequence}

The literature establishes both the complementarity of MRI sequences and the value
of self-supervised pretraining, but these strands remain only partially connected in
lumbar grading. Anatomical correspondence across sequences is usually treated as a
preprocessing requirement for fusion rather than as supervision. At the same time,
multisequence systems generally learn a common fusion mechanism rather than a
condition- and quality-dependent allocation that is deliberately challenged under
missing and corrupted modalities. The thesis therefore tests two linked ideas: that
DICOM geometry can define cross-sequence positive relationships, and that a single
model can use the resulting features through a target-conditioned router whose
failure modes are directly measurable.

\subsection{Gap C: Quantified Transfer and Adaptation Under a Harmonised External Task}
\label{sec:intro-gap-transfer}

Existing external validations establish that transport is difficult, but they
usually report one transported model at one point in the adaptation process---often
with no adaptation at all. The operational question for a receiving hospital is
instead how performance changes before adaptation and how quickly it recovers as
local labels are supplied. The literature reviewed for this thesis contains no
Middle Eastern evaluation of the same scope and no quantified parameter-efficient
annotation curve for this lumbar target formulation. The gap is therefore not
simply ``lack of external validation''; it is lack of a staged zero-shot-to-few-shot
transfer experiment with fixed local testing, controlled annotation budgets and
explicit reference-standard harmonisation.

\subsection{Supporting Gap: Directional Error and Calibrated Uncertainty}
\label{sec:intro-gap-error}

Ordinal severity and asymmetric error are clinically meaningful, but ordinal
objectives have not consistently outperformed categorical baselines
\cite{niemeyer2021a}. Calibration is also rarely reported in lumbar MRI despite the
known overconfidence of modern neural networks \cite{guo2017calibration}. These are
therefore treated as supporting questions rather than headline novelty. A
cost-sensitive objective is retained only if it reduces clinically consequential
under-grading without creating indiscriminate Severe false positives, and selective
prediction is evaluated as a model behaviour rather than proof of a validated
clinical workflow.

The novelty claim throughout the thesis will remain appropriately bounded. The
phrasing ``to our knowledge'' or ``in the reviewed literature'' is used where
necessary, and the final claim will be updated against the literature available at
the time of submission. A new paper that adopts one element of the proposed system
does not invalidate the PhD if the controlled experimental question remains
distinct; it does, however, change how originality must be described.

\laterinsert{Immediately before thesis submission, update this novelty boundary
against publications appearing after the current literature-search cut-off. Record
which claims remain original and which should be reframed as independent validation,
extension or comparative evidence.}

% ===========================================================================
\section{Research Aim}
\label{sec:intro-aim}

The overall aim of this PhD is:

\begin{quote}
\emph{to determine whether explicitly modelling anatomical relationships among
lumbar disease targets, anatomically aligned evidence across MRI sequences, and
target-dependent sequence availability improves the accuracy, robustness, label
efficiency and cross-institutional portability of automated lumbar spine MRI
severity grading, relative to strong non-relational and fixed-fusion baselines.}
\end{quote}

The aim is intentionally phrased as a determination rather than a promise of
superiority. The thesis succeeds scientifically if it establishes under controlled
conditions where these mechanisms help, where they fail, and what their limits are.

% ===========================================================================
\section{Research Objectives}
\label{sec:intro-objectives}

The aim is operationalised through the following objectives.

\begin{enumerate}
  \item To construct a reproducible patient-level lumbar MRI analysis pipeline that
        preserves DICOM patient-space geometry, performs reliable level
        localisation, and supplies identical anatomically defined regions of
        interest to all grading baselines and proposed models.

  \item To establish competitive independent and contextual grading baselines on
        the public LumbarDISC development data, including fixed fusion,
        cross-attention and ordered inter-level context where reproducibility
        permits, so that improvements cannot be attributed to a weak comparator.

  \item To develop and evaluate anatomically aligned cross-sequence
        self-supervision in which positive relationships are defined by the same
        patient and lumbar level across MRI sequences using DICOM geometry.

  \item To quantify whether this anatomical self-supervision improves supervised
        grading under 10\%, 25\%, 50\% and 100\% label availability and whether it
        preserves sequence-specific diagnostic information.

  \item To develop and evaluate disease-conditioned adaptive sequence routing that
        learns patient- and target-dependent weights over available MRI evidence
        while explicitly representing sequence availability and acquisition quality.

  \item To test the router under both missing-sequence and corrupted-sequence
        conditions, and to compare its behaviour with concatenation, averaging and
        ordinary cross-attention under matched inputs and training budgets.

  \item To construct a heterogeneous disease--anatomy graph over the 25
        level--condition--laterality targets, with typed adjacent-level,
        same-level cross-condition and bilateral relations, and to compare it with
        independent heads, an ordered level transformer, a homogeneous graph,
        ungated message passing and shuffled topology.

  \item To test whether relational reasoning introduces adjacent-level false
        positives in naturally isolated lesions and to retain the graph only where
        any performance or robustness benefit is not purchased by clinically
        incoherent information propagation.

  \item To compare categorical, ordinal and clinically asymmetric cost-sensitive
        grading objectives, together with probability calibration and uncertainty
        analysis, while explicitly reporting both severe under-grading and severe
        over-grading.

  \item To quantify zero-shot transfer from the multinational public benchmark to
        the independent Rizgary Hospital cohort on the shared central-canal grading
        task, separating localisation, image-domain and reference-standard effects
        as far as the retrospective data allow.

  \item To estimate an annotation--performance recovery curve using controlled
        local adaptation sizes $N\in\{10,25,50,100\}$ and to compare where a limited
        trainable-parameter budget is most effectively placed.

  \item To report all principal comparisons with patient-level uncertainty,
        repeated training where feasible, contribution-specific ablations, negative
        results and explicit methodological limitations.
\end{enumerate}

These objectives define a staged experiment rather than a monolithic architecture.
A component is not retained because it appears in the proposal; it is retained only
if the corresponding experiment demonstrates a reproducible contribution or a
meaningful robustness, calibration or interpretability advantage.

% ===========================================================================
\section{Research Questions}
\label{sec:intro-rqs}

The research questions are fixed prospectively and correspond directly to the
methodological tests in Chapter~\ref{ch:methodology}.

\begin{description}
  \item[RQ1 --- Relational disease modelling.] Does representing lumbar disease
  targets as a typed heterogeneous graph improve grading over independent target
  heads, an ordered five-level transformer and a homogeneous graph, and are any
  gains concentrated where single-target evidence is weakest?

  \item[RQ2 --- Anatomically aligned self-supervision.] Does cross-sequence
  pretraining based on DICOM-defined anatomical correspondence improve grading,
  label efficiency and transfer compared with ImageNet initialisation, generic
  medical pretraining and ordinary augmentation-based contrastive learning?

  \item[RQ3 --- Disease-conditioned modality routing.] Can one model learn
  target- and patient-dependent sequence weights while remaining robust to both
  missing and quality-degraded MRI sequences, without requiring a separately
  trained network for every combination of sagittal T1, sagittal T2/STIR and axial
  T2?

  \item[RQ4 --- Cross-institutional transfer.] What is the zero-shot degradation
  when a benchmark-trained model is applied to a previously unseen Middle Eastern
  clinical cohort, how much of that degradation is attributable to localisation
  versus grading/reference-standard shift, and how does performance recover as the
  amount of local annotation available for adaptation increases?

  \item[RQ5 --- Clinically asymmetric error and uncertainty.] Do ordinal and
  cost-sensitive objectives reduce clinically consequential distant errors,
  particularly Severe$\rightarrow$Normal/Mild, and can calibrated uncertainty
  support selective prediction without overstating clinical safety?
\end{description}

These questions are intentionally answerable by negative as well as positive
results. If a typed graph performs no better than a homogeneous graph, RQ1 is still
answered. If anatomical self-supervision improves label efficiency but not external
transfer, RQ2 yields a more limited but still informative result. This design avoids
defining doctoral success as the requirement that every proposed component win.

% ===========================================================================
\section{Research Hypotheses}
\label{sec:intro-hypotheses}

The study uses directional hypotheses where the literature justifies a prediction,
but avoids arbitrary target accuracies that would have no defensible basis before
implementation.

\begin{description}
  \item[H1 --- Anatomical self-supervision.] Anatomical cross-sequence
  self-supervision will improve label-efficient grading and cross-institutional
  robustness relative to generic pretraining, particularly when only a fraction of
  benchmark labels is used.

  \item[H2 --- Adaptive routing.] Disease-conditioned routing with explicit
  availability masking, modality dropout and quality-aware gating will degrade less
  under missing-sequence and corruption experiments than fixed concatenation or
  ordinary cross-attention trained only on clean complete studies.

  \item[H3 --- Heterogeneous relational structure.] A typed heterogeneous graph
  will outperform independent target heads and a homogeneous graph if the proposed
  anatomical edge types carry useful information; performance under a shuffled-edge
  control should deteriorate if the topology itself matters.

  \item[H4 --- External degradation.] Zero-shot performance on the Rizgary Hospital
  cohort will be lower than internal benchmark performance because of scanner,
  protocol, population, localisation and reference-standard shift; the magnitude
  and pattern of that degradation are treated as findings rather than as failure of
  the research programme.

  \item[H5 --- Annotation-dependent recovery.] Local adaptation will improve
  performance as annotation increases, but no predetermined percentage of recovery
  is required. The shape of the recovery curve and the location in the network at
  which limited adaptation is most effective are empirical outcomes.

  \item[H6 --- Asymmetric grading error.] Cost-sensitive and ordinal objectives may
  reduce distant Severe$\rightarrow$Normal/Mild errors even if they do not improve
  aggregate macro-F1; any reduction must be evaluated against false-positive Severe
  predictions and against the cross-entropy baseline
  \cite{niemeyer2021a}.
\end{description}

The hypotheses will not be rewritten after inspection of the held-out public or
Rizgary test results. If an implementation decision changes the operational test of
a hypothesis, that change will be entered into the protocol changelog with its date
and rationale.

% ===========================================================================
\section{Conceptual Framework}
\label{sec:intro-conceptual}

The thesis can be understood as a progression from \emph{local evidence} to
\emph{structured evidence} and finally to \emph{transported evidence}.

At the first level, a model receives anatomically localised image evidence for a
single level and target. This is the baseline condition: if local evidence alone is
sufficient, more complex mechanisms should not be retained. At the second level,
the model is permitted to combine complementary MRI sequences, but the combination
must respect anatomical correspondence and sequence availability. Anatomical
self-supervision asks whether the same structure under different acquisitions can
produce a more label-efficient representation, while disease-conditioned routing
asks how strongly each sequence should contribute to a given target.

At the third level, a target is no longer treated as an isolated endpoint. Its
representation becomes a node in a typed graph whose relations correspond to three
known structures of the lumbar grading problem: the ordered chain of levels,
coexisting compartments at a level, and bilateral pairing. Message passing is soft
rather than deterministic; no rule states that disease at one level forces disease
at its neighbour. The graph is therefore a hypothesis about useful context, not a
hard-coded disease progression model.

At the fourth level, predictions are evaluated as ordinal and probabilistic outputs.
The model is judged not only by aggregate discrimination but by error distance,
Severe-class performance, calibration and the false-positive cost of any
cost-sensitive objective. This stage asks whether a statistically improved model is
also better aligned with the shape of clinically consequential grading error.

The final level is institutional transport. The representation, router, graph and
output head developed on public data are frozen before external evaluation. The
Rizgary Hospital cohort then asks a different question from internal validation:
what remains true when the acquisition and reporting environment changes? Only
after zero-shot performance has been measured is local adaptation permitted. This
ordering prevents the local test cohort from becoming an informal development set.

Conceptually, the thesis therefore tests the proposition

\begin{equation}
\text{robust grading}
\;=\;
f(\text{local anatomy},\;\text{cross-sequence evidence},\;
\text{target relations},\;\text{uncertainty},\;\text{domain}),
\label{eq:intro-conceptual}
\end{equation}

where the equation is explanatory rather than a literal computational model. Each
term corresponds to a controlled experiment, and no term is assumed necessary until
its ablation supports that conclusion.

% ===========================================================================
\section{Significance of the Study}
\label{sec:intro-significance}

\subsection{Scientific Significance}
\label{sec:intro-significance-scientific}

The principal scientific significance lies in changing the unit of reasoning from
an isolated image label to a structured set of disease targets. Many improvements in
medical imaging can be produced by additional parameters, larger backbones or more
aggressive tuning. Such gains are difficult to interpret mechanistically. By
contrast, this thesis attaches each proposed contribution to a specific hypothesis
and a destructive control: anatomical edges are compared with homogeneous and
shuffled topology; adaptive routing is compared with fixed fusion and challenged by
missing and corrupted modalities; anatomical self-supervision is compared with
generic pretraining and tested for over-invariance. The resulting evidence can show
not merely whether the full model scored higher, but \emph{why} it did or did not.

This matters because recent lumbar literature already demonstrates sophisticated
contextual models \cite{chai2026anatomy,batra2025mscan,nguyen2026crossspine}. The
research contribution cannot therefore be the generic use of attention, graphs or
multiple sequences. Its value lies in isolating the anatomical assumptions and
measuring them against alternatives under one controlled pipeline.

\subsection{Methodological Significance}
\label{sec:intro-significance-methodological}

The study places unusually strong emphasis on validity controls. Patient identity,
not slice or target, defines the data split; self-supervised learning is prohibited
from seeing held-out patients; localisation is evaluated separately from grading;
zero-shot transfer is measured before adaptation; the external reference standard is
kept distinct from routine report extraction; and every claimed model contribution
requires an ablation. These choices respond directly to recurring weaknesses in the
medical-imaging literature, where leakage, non-comparable splits and weak external
validation can produce optimistic performance estimates
\cite{compte2023are,wang2024machine}.

The transfer design is methodologically significant for a second reason. The local
hospital cohort is not used simply as another test set. The experiment attempts to
disentangle localisation failure, image-domain shift and reference-standard shift,
then measures recovery as local labels are added. This converts ``generalisation''
from a binary claim into a measurable process.

\subsection{Translational and Regional Significance}
\label{sec:intro-significance-regional}

The regional significance follows from the use of an independent Iraqi hospital
cohort that was not curated to mirror the public benchmark. Models developed on
large international datasets are often assumed to transfer into settings whose
scanner mix, acquisition practice, disease prevalence and reporting conventions were
not represented in development. Testing that assumption is particularly important
for institutions that do not have the resources to reproduce a large annotation
campaign locally. If useful performance can be recovered from tens rather than
hundreds of local annotations, that is a practical finding; if not, the negative
result is equally important because it quantifies the annotation burden required for
responsible transfer.

The study does not claim that one external hospital represents all of Iraq or the
Middle East. Rizgary Hospital provides evidence of transport to one genuinely
independent setting. Regional generalisation would require additional hospitals and
is therefore a future study rather than an inference drawn from one centre.

\subsection{Clinical Significance and its Boundary}
\label{sec:intro-significance-clinical}

The intended clinical significance is improved technical support for structured MRI
grading: consistency, reproducibility, calibrated confidence and explicit handling
of incomplete imaging. The thesis does not claim that this will improve treatment
outcomes, referral decisions or radiologist productivity. Those claims require a
prospective reader-assistance or impact study, as demonstrated by work that measures
reader performance directly \cite{lim2022improved}. The present study therefore
uses language such as ``decision-support potential'' and ``technical portability'',
not ``clinical benefit''.

% Written after Chapter 4. States only what the executed experiments support.
Of the benefits anticipated above, the executed experiments support
\emph{discrimination} and, within it, the reduction of clinically distant
errors. Quadratic weighted kappa improves by $+0.0177$ over a single-sequence
baseline, Severe$\rightarrow$Normal confusions fall from $5.7\%$ to $3.8\%$ and
recall on Severe targets rises from $61.1\%$ to $63.1\%$, the last two
attributable to the ordinal and cost-sensitive objective rather than to any
proposed structural mechanism.

Robustness to missing sequences improves, but the improvement is attributable to
modality dropout rather than to disease-conditioned routing, which changes
sequence reliance by an amount indistinguishable from zero. Label efficiency and
cross-institutional transfer were not evaluated and no claim is made for either.
Calibration is reported but was not the object of a confirmatory comparison.

No inference is drawn from these figures about patient benefit. The study
measures agreement with a reference standard on a retrospective benchmark; it
does not measure any clinical outcome.

% ===========================================================================
\section{Proposed Contributions}
\label{sec:intro-contributions}

At the protocol stage, the contributions below collectively form the proposed \textbf{AMOG-Net} framework and are \emph{proposed contributions}.
They become demonstrated thesis contributions only if the experiments support them.

\subsection{Contribution 1: Anatomically Aligned Cross-Sequence Self-Supervision}
\label{sec:intro-contrib-ssl}

The first proposed contribution uses DICOM-defined patient-space correspondence to
construct positive relationships between different MRI sequences depicting the same
patient and lumbar level. The objective is not to invent a new generic contrastive
loss; established contrastive or non-contrastive objectives may be used. The
contribution under test is the anatomical definition of correspondence and its value
for label efficiency, grading and transport. A dedicated preservation experiment
checks whether alignment erases sequence-specific diagnostic information.

\subsection{Contribution 2: Disease-Conditioned and Quality-Aware Sequence Routing}
\label{sec:intro-contrib-routing}

The second proposed contribution allows one study to be fused differently for
different disease targets. Routing depends on target identity, patient features,
sequence availability and quality evidence, and is evaluated under genuine or
simulated missingness and corruption. The contribution is not the existence of a
softmax gate or mixture-of-experts mechanism; it is the target-conditioned allocation
of clinically complementary evidence and the explicit robustness evaluation that
accompanies it.

\subsection{Contribution 3: Heterogeneous Disease--Anatomy Graph Reasoning}
\label{sec:intro-contrib-graph}

The third proposed contribution represents the 25 benchmark outputs as nodes in a
heterogeneous graph. Typed edges encode adjacent-level, within-level cross-condition
and bilateral relations. The graph includes a mechanism for preserving local
evidence when neighbour messages are unhelpful and is evaluated against independent,
ordered-transformer, homogeneous, ungated and shuffled-topology controls. A focal
isolated-lesion test determines whether relational reasoning improves uncertain
predictions without artificially spreading disease to adjacent normal levels.

\subsection{Contribution 4: A Structured Cross-Institutional Transfer Evaluation}
\label{sec:intro-contrib-transfer}

The translational contribution is an experimental design rather than a new network
layer. The selected public-data model is evaluated zero-shot on a fixed Rizgary
Hospital test set before any local adaptation. A separate local adaptation pool is
then sampled at controlled annotation sizes and used to compare constrained update
strategies. The analysis is strengthened by a verified-localisation control and,
where resources permit, harmonised independent reader grading of central canal
stenosis. This design estimates both the cost of transport and the cost of recovery.

\subsection{Supporting Contributions}
\label{sec:intro-contrib-supporting}

Ordinal and cost-sensitive objectives, calibration, uncertainty estimation and
parameter-efficient adaptation support the principal contributions but are not
claimed as original methods. Their value lies in testing whether the proposed
structured model behaves sensibly under clinically relevant error definitions and
limited local annotation. They are retained only when their contribution is
measurable.

% --------------------------------------------------------------------------
% Written at final submission, replacing the [LATER INTEGRATION] placeholder.
% Only components supported by the locked experiments are claimed. Components
% that failed are reported below as negative findings and remain in the
% research history and in the ablation ladder.
% --------------------------------------------------------------------------
\subsection{Demonstrated Contributions}
\label{sec:intro-contrib-demonstrated}

The experiments reported in Chapter~\ref{ch:results} support three
contributions. They are not the three proposed above, and the difference is
itself part of what the thesis contributes.

\paragraph{D1 --- A verified anatomy-aware multi-sequence grading system.}
AMOG-Net improves quadratic weighted kappa over a single-sequence baseline by
$+0.0177$ (95\% CI $[+0.0097, +0.0260]$), on every one of seven training seeds
and surviving false discovery rate correction across the pre-specified
comparison family. Two further comparisons also survive: typed heterogeneous
edges and the ordinal, cost-sensitive objective. The system, its frozen data
partition and its behavioural test suite are released.

\paragraph{D2 --- A controlled ablation protocol for anatomical priors.} Each
claimed mechanism is tested against a control matched in capacity rather than
against its own absence: a degree-preserving edge shuffle for graph topology, a
router-free fixed-fusion arm for sequence routing, an architecture-matched
untrained network as an attribution floor, and a label-shuffled model as a
pipeline-level negative control. Section~\ref{sec:results-threats} documents
three analyses that produced incorrect positive conclusions before the
corresponding control was introduced, in each case in the direction of the
hypothesis under test.

\paragraph{D3 --- Evidence that anatomical structural priors do not improve
grading at this data scale, with a mechanism.} Neither anatomically aligned
cross-sequence self-supervision nor disease-conditioned routing produces a
detectable improvement, and anatomically correct graph topology does not
outperform a degree-preserving random shuffle. These are bounds rather than
absences: the paired intervals constrain the three mechanisms to at most
$1.32\%$, $1.55\%$ and $1.05\%$ of baseline performance respectively. Gradient-weighted class activation
mapping shows why: the convolutional encoder already concentrates $2.4$ times the
chance share of its attention on the annotated target before any prior is added.
The localisation these priors were designed to supply is largely accomplished
upstream of them.

\subsection{Proposed Contributions Not Supported}
\label{sec:intro-contrib-negative}

Contribution~1 (anatomical cross-sequence self-supervision) is rejected on
accuracy, on cross-sequence robustness under controlled input ablation, and on
attention concentration. The pretext task is learnable and was learned ---
validation InfoNCE $1.0771$ against a chance value of $3.4657$ --- but the
representation confers no measurable downstream benefit.

Contribution~2 (disease-conditioned routing) divides. The router does learn
clinically sensible target-dependent sequence weights, reproducing the expected
allocation in all fifteen runs in which a gate is present. It does not improve
grading ($+0.0053$ QWK, five seeds of seven, bounded at $1.55\%$ of baseline)
and it does not improve robustness to missing sequences.
Its allocation tracks the annotation structure of the benchmark rather than
creating a dependence of its own.

Contribution~3 (heterogeneous disease--anatomy graph) is partly supported. Typed
relational structure improves grading over a homogeneous graph ($+0.0093$ QWK on
every one of seven seeds, surviving correction at $p_{\mathrm{FDR}} = 0.009$);
the anatomical identity of the relations does not. The resulting
claim is narrower than proposed and more specific than the original: relational
typing carries information, anatomical adjacency does not.

Contribution~4 (cross-institutional transfer evaluation) was not executed. Cohort
preparation is complete and is reported in Section~\ref{sec:results-rq4};
adjudication of conflicting reference labels and a correct de-identification
procedure remain outstanding. It is retained here as unexecuted rather than
withdrawn.

% ===========================================================================
\section{Scope and Delimitations}
\label{sec:intro-scope}

Clear boundaries are necessary because lumbar MRI supports a much broader clinical
and computational agenda than one PhD can test.

\subsection{Clinical Scope}
\label{sec:intro-scope-clinical}

The primary public task is degenerative lumbar disease severity grading at five
levels from L1--L2 through L5--S1. On LumbarDISC, the full target space contains
central canal stenosis, left and right neural foraminal narrowing, and left and
right subarticular stenosis, each graded as Normal/Mild, Moderate or Severe
\cite{richards2026the}. The study does not attempt to diagnose every cause of low
back pain, infer symptoms, determine surgical candidacy or predict long-term patient
outcomes.

The primary local transfer task is narrower. Rizgary Hospital reports support a
defensible central-canal comparison more readily than the complete benchmark schema.
The confirmatory external analysis is therefore restricted to central canal stenosis
at the five lumbar levels unless new expert grading expands the local reference
standard. Disc bulge, protrusion and extrusion are retained as a separate local
morphology task and are not heuristically converted into stenosis grades.

\subsection{Imaging Scope}
\label{sec:intro-scope-imaging}

The principal imaging inputs are sagittal T1, sagittal T2/STIR and axial T2 MRI.
Other sequences may appear in clinical studies but do not define the primary model
unless a justified later protocol amendment is made. DICOM geometry is preserved
because cross-plane correspondence is part of the scientific design. PNG or other
rendered formats may be used for visual inspection but do not replace the geometric
metadata required for alignment.

\subsection{Methodological Scope}
\label{sec:intro-scope-method}

Anatomical localisation and segmentation are enabling technologies. They are
measured and controlled because a poor crop can invalidate grading, but a new
localisation architecture is not the headline doctoral contribution. Likewise,
backbone selection is controlled rather than treated as a novelty competition. The
research may use ResNet-, EfficientNet-, Swin- or other contemporary encoders, but
the primary backbone is frozen before the main contribution-specific ablations.

The graph topology is defined anatomically rather than discovered freely from the
test data. Adjacent levels, same-level conditions and bilateral counterparts form the
pre-specified relation families. Learning the globally optimal topology is a
different research problem and is excluded from the primary scope.

\subsection{Data and Institutional Scope}
\label{sec:intro-scope-data}

Model development uses the public LumbarDISC benchmark and, where permitted, SPIDER
for anatomical localisation or segmentation support. The locally held labelled
competition subset contains approximately 1,975 studies at the protocol stage;
\torecord{final public development, validation and held-out counts after quality
control}. The local project inventory currently identifies 299 narrative reports
and 294 candidate matched multi-sequence cases, while a wider raw DICOM inventory
contains additional studies. These numbers are not presented as the final cohort
until one-to-one reconciliation is complete:
\torecord{final Rizgary case-flow counts and reasons for exclusion}.

Before analytical use of the local cohort, the Rizgary data-preparation phase will
follow the approved institutional governance protocol. A controlled research copy will
be de-identified before access by the model-development team, and the permitted storage
and compute environment will be verified before local processing begins. These safeguards
precede, rather than form part of, model development.

Rizgary Hospital is a single external institution. The study can therefore support a
claim of transport to that independent setting, not universal Middle Eastern
performance. Additional regional hospitals would be required for a multicentre
regional generalisation claim.

\subsection{Ethical and Governance Scope}
\label{sec:intro-scope-ethics}

The local component is retrospective and does not alter patient care. Formal use of
the local data remains conditional on the relevant institutional permissions:
\toconfirm{hospital/ethics approval reference, date and responsible institution};
\toconfirm{retrospective consent or waiver basis}; and
\toconfirm{permitted compute environment, including whether external cloud/GPU
processing is authorised}. Identifiable raw DICOM is not supplied to students or
external compute; a controlled de-identified research copy is required before local
analysis.

\subsection{Evaluation Scope}
\label{sec:intro-scope-evaluation}

The thesis evaluates technical grading performance, robustness, calibration,
computational cost and cross-institutional portability. The unit of statistical
independence is the patient. It does not claim prospective clinical impact,
radiologist replacement or treatment benefit. Reader-assistance evaluation is a
logical future study if the technical system proves sufficiently reliable.

% ===========================================================================
\section{Assumptions and Operational Definitions}
\label{sec:intro-definitions}

The following terms are used consistently throughout the thesis.

\begin{longtable}{p{3.6cm}p{10.0cm}}
\toprule
\textbf{Term} & \textbf{Operational meaning in this thesis}\\
\midrule
\endfirsthead
\toprule
\textbf{Term} & \textbf{Operational meaning in this thesis}\\
\midrule
\endhead
Lumbar level & One of L1--L2, L2--L3, L3--L4, L4--L5 or L5--S1.\\
Target & A level--condition--laterality output on LumbarDISC; 25 targets per study at full benchmark scope.\\
Severity grade & Normal/Mild, Moderate or Severe, encoded as an ordered three-class outcome.\\
Local evidence & Image-derived representation for one target before graph message passing.\\
Anatomical correspondence & Physical correspondence between acquisitions defined through DICOM patient-space geometry and lumbar localisation, not through image appearance alone.\\
Sequence routing & Learned allocation of weight over available MRI sequence features for a specific patient and target.\\
Missing modality & An expected MRI sequence is absent and explicitly unavailable to the model.\\
Degraded modality & A sequence is present but technically compromised, for example by motion, truncation, noise, slice loss or poor coverage.\\
Heterogeneous graph & A graph whose nodes are disease--anatomy targets and whose edges belong to clinically defined relation types.\\
Zero-shot transfer & Evaluation of the frozen public-data model on the fixed local test set before local fine-tuning, threshold selection or local calibration.\\
Few-shot adaptation & Model adaptation using only the separate local adaptation pool under pre-specified annotation budgets.\\
Reference-standard shift & Difference between the process used to generate labels in the benchmark and the process used to establish the local external reference.\\
Selective prediction & Model abstention or referral based on uncertainty; this is a technical behaviour, not evidence of a validated clinical workflow.\\
\bottomrule
\end{longtable}

These definitions are intentionally narrower than ordinary-language usage. In
particular, ``external validation'' does not mean that a model has been proven safe
for deployment, and ``uncertainty'' does not mean that the network has measured a
patient's clinical risk.

% ===========================================================================
\section{Overview of the Research Design}
\label{sec:intro-design-overview}

The research is a retrospective, multi-dataset medical-imaging methods study. Its
logic is cumulative, but the conclusions remain component-specific.

\begin{enumerate}
  \item \textbf{Data and anatomy foundation.} Public data are inventoried,
        partitioned by patient, reconstructed in DICOM patient space, anatomically
        localised and converted into disease-specific 2.5D regions of interest.

  \item \textbf{Strong baseline establishment.} An independent ROI classifier and
        competitive multiview/context baselines establish what can be achieved
        without the proposed relational formulation.

  \item \textbf{Representation experiment.} Anatomically aligned cross-sequence
        self-supervision is evaluated against generic and medical pretraining under
        several label fractions and sequence-specific preservation probes.

  \item \textbf{Fusion experiment.} Disease-conditioned sequence routing is tested
        against fixed and attention-based fusion under complete, missing and
        corrupted sequence conditions.

  \item \textbf{Relational experiment.} Homogeneous and heterogeneous target graphs
        are compared with independent and ordered-transformer controls, including
        shuffled topology and isolated-lesion stress testing.

  \item \textbf{Output and uncertainty experiment.} Categorical, ordinal and
        cost-sensitive objectives are evaluated together with calibration and
        uncertainty, with both under-grading and over-grading reported.

  \item \textbf{External transport experiment.} The selected public-data model is
        frozen and evaluated zero-shot on the fixed Rizgary Hospital test set. A
        verified-ROI control and harmonised reader reference are used where feasible
        to interpret the source of degradation.

  \item \textbf{Adaptation experiment.} Controlled local annotation budgets are
        sampled from a separate adaptation pool to estimate how quickly performance
        recovers and where a limited trainable-parameter budget is most effective.
\end{enumerate}

This order deliberately prevents the final architecture from becoming an indivisible
engineering object. Each stage introduces one assumption that can be rejected
without invalidating the remainder of the thesis. The final model is therefore the
simplest configuration supported by the evidence, not necessarily the configuration
containing every proposed module.

% ===========================================================================
\section{Expected Outcomes and Criteria for Scientific Success}
\label{sec:intro-success}

The protocol does not define success as reaching a predetermined accuracy. Such a
threshold would be difficult to justify because performance depends on target,
severity distribution, reference-standard reliability and the final patient split.
Scientific success is instead defined by whether the research questions can be
answered under controlled conditions.

For RQ1, success means establishing whether typed anatomical relations add value
beyond capacity-matched controls and whether they avoid adjacent-level
contamination. A null result is informative if the comparison is sufficiently
powered and the topology controls are valid. For RQ2, success means establishing
whether anatomical cross-sequence pairing improves label efficiency, transfer or
neither, while documenting whether sequence-specific information is preserved. For
RQ3, success means quantifying robustness under missing and degraded modalities and
determining whether adaptive routing adds value over simpler fusion. For RQ4,
success means producing an interpretable zero-shot transfer estimate and a controlled
adaptation curve, even if the external drop is large. For RQ5, success means
characterising the trade-off between severe under-grading, severe over-grading and
probability reliability rather than selecting a loss solely because it raises one
aggregate metric.

At final submission, the thesis will distinguish three categories explicitly:
\emph{supported contributions}, \emph{negative findings}, and \emph{unresolved
questions}. This is preferable to rewriting the original proposal so that only
successful components appear to have been planned.

% Written after Chapter 4. High-level answers only; metrics live in Chapter 4.
\begin{description}
  \item[RQ1 --- Relational disease modelling.] \textbf{Partly supported.} A typed
  heterogeneous graph improves grading over independent heads and over a
  homogeneous graph, on every seed. The anatomical identity of the relations does
  not contribute: an anatomically correct topology does not outperform a
  degree-preserving random shuffle.

  \item[RQ2 --- Anatomically aligned self-supervision.] \textbf{Not supported.}
  The pretext task is learnable and was learned, but the resulting representation
  yields no detectable improvement in grading, in robustness to withheld
  sequences, or in where the model attends.

  \item[RQ3 --- Disease-conditioned modality routing.] \textbf{Mechanism
  supported, benefit not.} The router reliably learns the clinically expected
  sequence weights, but a router-free model shows the same sequence dependence
  under intervention, and routing improves neither accuracy nor missing-sequence
  robustness.

  \item[RQ4 --- Cross-institutional transfer.] \textbf{Not executed.} Cohort
  preparation is complete; reference-label adjudication and de-identification
  remain outstanding.

  \item[RQ5 --- Clinically asymmetric error and uncertainty.] \textbf{Supported
  for error asymmetry.} Ordinal and cost-sensitive objectives improve both
  aggregate agreement and recall on Severe targets, on every seed, and form the
  most reliable single contributor to final performance. Selective prediction
  under calibrated uncertainty was not evaluated confirmatorily.
\end{description}

% ===========================================================================
\section{Thesis Organisation}
\label{sec:intro-organisation}

At the current protocol stage, the thesis is organised around the progression from
problem definition to evidence.

\textbf{Chapter 1 --- Introduction} defines the clinical and methodological
motivation, research problem, gaps, aim, objectives, research questions, hypotheses,
proposed contributions, significance and scope of the thesis.

\textbf{Chapter 2 --- Literature Review} examines the clinical grading systems that
define the reference labels, MRI sequence--pathology correspondence, localisation
and segmentation, backbone and contextual architectures, class imbalance and
ordinal grading, public datasets, external validation, interpretability and the
specific gaps that motivate the present study.

\textbf{Chapter 3 --- Methodology} translates those gaps into a protocol-grade
experimental design. It defines the public and local cohorts, leakage controls,
DICOM reconstruction, localisation and ROI construction, self-supervised
pretraining, adaptive routing, heterogeneous graph reasoning, grading objectives,
calibration, ablations, external validation, domain adaptation, statistical analysis
and reproducibility safeguards.

\textbf{Chapter 4 --- Results and Analysis} will report the staged E0--E8
experiments, including component-specific ablations, robustness analyses,
calibration, zero-shot external transfer, few-shot adaptation and failure analysis.
% Written after Chapter 4. Retained experiment groups; research questions unchanged.
The retained experiment groups are the eight-rung ablation ladder E0--E7 with two
matched controls on the graph configuration, a controlled input ablation testing
sequence dependence by intervention rather than by learned allocation, and a
Grad-CAM attribution analysis against an architecture-matched untrained floor.
The cross-institutional transfer group was prepared but not executed and is
reported as such. The research questions are unchanged from those fixed in
Section~\ref{sec:intro-rqs}.

\textbf{Chapter 5 --- Discussion and Conclusions} will interpret the findings
against the literature, distinguish demonstrated contributions from negative
results, state limitations, discuss implications for transport to other institutions
and identify the next studies required before any clinical-impact claim.

\toconfirm{final institutional thesis structure if the university requires results
and discussion as separate chapters, publication-based chapters, or an additional
conclusion chapter}. The scientific ordering above should remain recognisable even
if the administrative chapter numbering changes.

% ===========================================================================
\section{Rules for Completing Prospective Placeholders}
\label{sec:intro-placeholders}

Because this chapter is written before completion of implementation, several fields
are intentionally prospective. The following rules govern their later completion.

\begin{enumerate}
  \item \textbf{[TO CONFIRM]} is used only for institutional or governance facts
        that are currently unknown, such as an approval reference or permitted
        compute environment. These fields are replaced only with documented facts.

  \item \textbf{[TO RECORD AFTER IMPLEMENTATION]} is used for executed technical or
        cohort quantities that cannot responsibly be specified before the pipeline
        runs, such as final quality-controlled sample counts.

  \item \textbf{[LATER INTEGRATION]} is used for a small number of places where the
        final thesis may state the principal findings or update the novelty boundary.
        These additions may report evidence; they must not alter the prospective
        research questions or make an unsuccessful hypothesis disappear.

  \item Any substantive protocol change made after results become visible is
        documented with date, reason and affected analysis. Confirmatory and
        exploratory analyses remain distinguishable.

  \item The introduction will not be retrofitted to imply that the final retained
        architecture was inevitable. The final thesis should preserve the fact that
        ACSSL, adaptive routing, heterogeneous graph reasoning and supporting losses
        were hypotheses subjected to testing.
\end{enumerate}

These rules turn placeholders into a transparency mechanism rather than an invitation
to hindsight bias.

% ===========================================================================
\section{Chapter Summary}
\label{sec:intro-summary}

This chapter has defined the thesis as a study of structured and transferable lumbar
MRI grading rather than as a competition between image backbones. The clinical task
is inherently multi-level, multi-condition, bilateral and multi-sequence, while its
reference labels are ordinal and imperfectly reproducible. The technical literature
has established that localisation should precede grading, that sophisticated
multiview systems can perform strongly on public data, and that external performance
frequently degrades. Those achievements narrow rather than remove the research gap.

The thesis therefore asks whether three explicit structural mechanisms add value:
anatomically aligned cross-sequence self-supervision, disease-conditioned and
quality-aware sequence routing, and a typed heterogeneous graph over
level--condition--laterality targets. These mechanisms are evaluated component by
component against strong controls rather than bundled into a single architecture and
assumed beneficial. Supporting ordinal, cost-sensitive and uncertainty methods are
used to test the clinical shape of error without being advertised as standalone
novelty.

The translational component separates model development from external evaluation.
LumbarDISC supplies the public development benchmark; SPIDER may support anatomical
localisation; and the independent Rizgary Hospital cohort is reserved for zero-shot
central-canal transfer and controlled local adaptation. The local reference process
is treated explicitly as a potential source of shift rather than assumed equivalent
to the public benchmark. Clinical impact remains outside the scope until a separate
reader-assistance or prospective study is conducted.

The aim, objectives, RQ1--RQ5 and H1--H6 stated here form the prospective frame for
Chapters~\ref{ch:litreview} and~\ref{ch:methodology}. The remainder of the thesis is
therefore not required to prove that every proposed mechanism succeeds. It is
required to determine, with reproducible evidence, which assumptions survive
controlled testing and institutional transfer.

% Written at final submission.
In summary, the complete system improves on a single-sequence baseline
reproducibly and by a margin that survives correction for multiple comparisons
while remaining far too small to alter clinical practice,
and the objective function designed around clinically asymmetric error is the
most reliable single contributor to that improvement. Of the three proposed
structural mechanisms, relational typing is supported, anatomical topology and
disease-conditioned routing are not, and anatomically aligned self-supervision is
rejected on every axis on which it was tested; attribution analysis indicates
that the convolutional encoder already performs the localisation these mechanisms
were intended to supply. The principal limitations are that all results are
internal to one benchmark, that cross-institutional transfer was not executed,
and that the study is underpowered for single-component effects of the magnitude
observed.

% ===========================================================================
\printbibliography[heading=bibintoc,title={References}]

\end{document}
 from the main thesis file.
% ===========================================================================

\documentclass[12pt,a4paper]{report}

\usepackage{lmodern}
\usepackage[T1]{fontenc}
\usepackage[utf8]{inputenc}
\usepackage[english]{babel}
\usepackage{csquotes}
\usepackage{microtype}
\usepackage[margin=2.5cm]{geometry}
\usepackage{setspace}
\onehalfspacing

\usepackage{amsmath,amssymb}
\usepackage{booktabs}
\usepackage{longtable}
\usepackage{array}
\usepackage{graphicx}
\usepackage{thesisstyle}   % shared TikZ figure styles; see thesisstyle.sty
\usepackage[hidelinks]{hyperref}

\usepackage[
  backend=biber,
  style=numeric-comp,
  sorting=none,
  maxbibnames=6,
  minbibnames=3,
  giveninits=true,
  doi=true,
  url=false,
  isbn=false
]{biblatex}

% Production bibliography cleaned for thesis compilation. Citation keys are
% unchanged from the working literature inventory.
\addbibresource{lumbar_spine_mri_ai_references_clean.bib}

\newcommand{\toconfirm}[1]{\textbf{[TO CONFIRM: #1]}}
\newcommand{\torecord}[1]{\textbf{[TO RECORD AFTER IMPLEMENTATION: #1]}}
\newcommand{\laterinsert}[1]{\textbf{[LATER INTEGRATION: #1]}}

\begin{document}

\chapter{Introduction}
\label{ch:introduction}

% ===========================================================================
\section{Introduction and Thesis Position}
\label{sec:intro-position}

Lumbar degenerative disease presents an unusually demanding problem for medical
image analysis. The anatomical region is compact, but the diagnostic task is not:
one examination contains several intervertebral levels, several pathological
compartments, left--right distinctions, multiple magnetic resonance imaging (MRI)
sequences acquired in different planes, and severity labels whose boundaries are
partly subjective. A clinically meaningful automated system must therefore do more
than recognise that degeneration is present. It must identify \emph{what} is
abnormal, \emph{where} it is abnormal, \emph{how severe} the abnormality is, and
\emph{which imaging evidence} supports that judgement. It must also remain
interpretable as a measurement of radiological grading rather than be mistaken for
an autonomous clinical diagnosis.

The need for such systems arises from both disease burden and reporting workload.
Low back pain remains one of the largest contributors to disability worldwide, and
lumbar degenerative change is a frequent structural correlate examined by MRI
\cite{compte2023are,wang2024machine}. Routine interpretation is repetitive because
the reader must inspect multiple levels across sagittal and axial acquisitions,
while the severity of central canal, foraminal and subarticular narrowing is judged
against morphological criteria that are not perfectly reproducible between human
readers \cite{lurie2008reliability,schizas2010qualitative,lee2011a,lee2010a}.
Automation is attractive in this setting not because a neural network is assumed to
possess a superior biological truth, but because a reproducible system may help
standardise a high-volume structured task whose human reference process itself has
known variability.

The technical field has advanced rapidly. Early systems demonstrated that lumbar
anatomy could be localised and individual discs graded with convolutional networks;
more recent pipelines combine explicit localisation, 2.5D regions of interest,
multiview fusion, attention and large public benchmarks
\cite{lu2018deepspine,pan2021automatically,batra2025mscan,richards2026the}. The
central problem has consequently shifted. The thesis no longer asks whether a
convolutional neural network or transformer can classify a lumbar MRI in isolation.
It asks whether the \emph{structure of the grading problem itself} has been
represented faithfully enough.

This thesis takes the position that lumbar MRI grading is a structured prediction
problem with four coupled sources of information. First, disease occurs at ordered
anatomical levels rather than at unrelated image locations. Second, several
conditions coexist within one motion segment and lateral compartments occur as
paired left--right structures. Third, the diagnostic value of an MRI sequence is
target-dependent: sagittal T1, sagittal T2/STIR and axial T2 provide complementary
rather than interchangeable evidence \cite{hallinan2021deep}. Fourth, a model that
performs well on one benchmark may degrade when scanners, protocols, population and
reference-standard conventions change across institutions
\cite{mcsweeney2023external,zhang2023deep,guan2022domain}.

The research therefore evaluates three methodological contributions and one
translational validation programme. The integrated framework evaluating these components
is designated \textbf{AMOG-Net} (\textbf{A}natomy-aware \textbf{M}ulti-sequence \textbf{O}rdinal \textbf{G}raph Network). The first contribution is anatomically aligned cross-sequence
self-supervision, in which DICOM patient-space correspondence defines which views
should agree during representation learning. The second is disease-conditioned
adaptive sequence routing, in which the model learns target- and patient-dependent
weights over available MRI sequences and is explicitly tested under both missing
and degraded inputs. The third is a heterogeneous disease--anatomy graph whose
nodes represent level--condition--laterality targets and whose typed edges encode
adjacent-level, same-level cross-condition and bilateral relations. These are then
evaluated under a deliberately separate transfer programme: development occurs on
a multinational public benchmark; the selected model is frozen; zero-shot
performance is measured on an independent Rizgary Hospital cohort; and controlled
few-shot adaptation is performed only afterwards.

Several related techniques are necessary but are not claimed as doctoral novelty.
Anatomical localisation, segmentation, 2.5D region construction, conventional CNN
or transformer encoders, class-imbalance handling, ordinal loss functions,
calibration and parameter-efficient fine-tuning are established tools. They are
used to construct a fair experiment around the thesis contributions rather than
presented as original merely because they appear in the same architecture. This
novelty boundary is important in a field where anatomy-guided contextual grading,
cross-sequence attention and graph-based disc modelling have already been reported
\cite{chai2026anatomy,nguyen2026crossspine,baur2024automated}. The question is not
whether any one of those mechanisms exists elsewhere; it is whether the particular
representation tested here---typed relationships among the benchmark's disease
outputs, anatomical cross-sequence supervision, target-conditioned evidence
allocation and explicit institutional transfer---adds measurable knowledge beyond
strong matched baselines.

% Written after the confirmatory experiments were locked. Reports the outcome
% only; the problem statement, research questions and hypotheses are unchanged.
The principal empirical outcome is that the complete system improves agreement
with the reference standard over a single-sequence baseline by $+0.0177$
quadratic weighted kappa, on every one of seven training seeds and after
correction for multiple comparisons, but that the anatomical mechanisms proposed
to achieve this do not account for the improvement: attribution analysis
indicates the convolutional encoder already localises the graded structure
without them.

% ===========================================================================
\section{Background and Motivation}
\label{sec:intro-background}

\subsection{Clinical Motivation: A Repetitive but Structured Imaging Task}
\label{sec:intro-clinical-motivation}

Lumbar degeneration is not a single lesion. It is a family of related structural
changes involving the intervertebral disc, facet joints, ligamentum flavum,
vertebral endplates and the neural spaces that these structures surround. Disc
dehydration and height loss alter load distribution; annular bulging or herniation
may encroach on the canal and foramina; posterior element hypertrophy may further
reduce the space available to neural tissue. These changes are assessed at several
levels, most commonly from L1--L2 through L5--S1, and their significance depends on
which anatomical compartment is narrowed.

Three stenotic compartments are especially relevant to the benchmark task used in
this thesis. Central canal stenosis concerns the thecal sac and cauda equina,
foraminal narrowing concerns the exiting nerve roots, and subarticular or lateral
recess stenosis concerns the traversing roots. Their visual criteria are related but
not identical, and neither are the acquisitions from which they are best judged.
Central canal and lateral recess morphology depend heavily on axial T2 information,
whereas neural foraminal assessment benefits particularly from sagittal or
parasagittal T1 because perineural fat provides the relevant contrast
\cite{hallinan2021deep,schizas2010qualitative,lee2010a}. A model that treats all
sequences as equivalent therefore ignores a clinical property of the task before
learning has even begun.

The labels themselves are another source of complexity. Morphological grading
systems convert continuous anatomical variation into ordinal categories. Human
agreement is strongest for some compartments and weaker for others: the reliability
study of Lurie et al. found substantially better inter-reader agreement for central
canal stenosis than for foraminal and subarticular findings
\cite{lurie2008reliability}. This matters methodologically. If the reference standard
contains boundary uncertainty, model evaluation must distinguish gross errors from
one-grade disagreements and must interpret performance relative to the reliability
of the grading instrument. Apparent model error and reference-standard disagreement
are not always separable.

This combination---high examination volume, repeated level-wise assessment,
complementary sequence evidence and imperfect categorical boundaries---makes lumbar
MRI a suitable domain for artificial intelligence, but only if the automation is
framed correctly. The objective is not to replace full clinical interpretation.
Symptoms, neurological findings, prior imaging and management context remain outside
the information available to an image-only grading system. The narrower objective is
to make a defined radiological grading task more consistent, measurable and
portable. Reader-assistance studies support this framing: deep-learning assistance
has been shown to reduce reporting time and improve or preserve inter-reader
agreement, with the greatest gains often occurring among less experienced readers
\cite{lim2022improved}. This thesis therefore treats the system as a potential
component of decision support, not as an autonomous treatment decision-maker.

\subsection{Technical Motivation: From Classification to Structured Prediction}
\label{sec:intro-technical-motivation}

The evolution of lumbar MRI automation has progressively narrowed the unit of
prediction. Early work often asked whether a lesion or abnormality was present in an
image. Subsequent systems localised discs, graded individual levels, and shared
encoders across several pathologies. Public benchmark design has now made the task
more explicit. The RSNA Lumbar Degenerative Imaging Spine Classification resource,
published as LumbarDISC, contains 2,697 patients and 8,593 MRI series collected
across eight institutions in six countries \cite{richards2026the}. The associated
challenge requires five condition/laterality targets at five lumbar levels, producing
25 ordinal outputs per study \cite{rsna2024rsna,richards2026the}. A correct model
therefore does not simply output ``mild'', ``moderate'' or ``severe'' for an image;
it attaches one of those grades to a particular condition at a particular level and,
for lateral compartments, a particular side.

This target structure exposes a limitation in the prevailing formulation. Most
systems share image features but ultimately emit independent output heads. Shared
features allow several tasks to benefit from a common representation, but a
prediction at L4--L5 does not explicitly inform an uncertain prediction at L3--L4;
a confident left foraminal prediction does not interact structurally with its right
counterpart; and the central, foraminal and subarticular targets at one motion
segment do not exchange typed messages simply because they coexist anatomically.
The representation is therefore multi-task but not fully relational.

Graph learning provides a natural language for testing whether those relationships
carry useful information. Graph neural networks are not novel by themselves, and a
lumbar application has already reconstructed an individual disc as a point cloud and
graded it with a graph network \cite{baur2024automated}. That precedent is important
because it prevents an unjustified claim that ``using a graph in lumbar MRI'' is new.
The unresolved question is narrower: whether the \emph{prediction targets} can be
represented as nodes connected by clinically motivated relation types, and whether
those relations improve grading beyond independent heads, an ordered level
transformer and a homogeneous graph. The distinction between a graph over points
within one disc and a graph over disease--anatomy targets across the spine is the
novelty boundary this thesis tests.

A second technical motivation concerns multimodal evidence. Recent methods have
already shown that cross-view and cross-sequence attention can improve lumbar
assessment \cite{batra2025mscan,nguyen2026crossspine}, while anatomy-guided
inter-level context has also become an active direction \cite{chai2026anatomy}.
These developments mean that novelty cannot rest on ``using multiple sequences'' or
``using attention''. The remaining question is whether evidence allocation should be
\emph{conditioned on the target and the patient}, and whether one trained model can
remain operational when an expected sequence is unavailable or present but of poor
quality. A fixed fusion rule assumes the same contribution of sagittal T1, sagittal
T2/STIR and axial T2 for every target and every study. Clinical sequence--pathology
correspondence makes that assumption doubtful.

A third motivation is representation learning. Generic ImageNet pretraining is
convenient but does not exploit the particular correspondence available in MRI: the
same anatomical level can be observed under different contrasts and planes. Earlier
spinal work demonstrated that self-supervision derived from patient-level or
longitudinal information can reduce dependence on manual labels
\cite{jamaludin2017selfsupervised}, and disc-centric contrastive work continues that
line \cite{acharya2026disccentric}. DICOM geometry makes a stronger, directly
anatomical pairing possible. If sagittal and axial observations can be matched in
patient space, then the correspondence itself can serve as supervision without
pretending that the images should look visually alike.

Finally, there is a translational motivation. Systematic reviews repeatedly find
that internally impressive medical-imaging models underperform when re-evaluated
outside their development distribution \cite{compte2023are,wang2024machine}.
Lumbar-specific external validations show the same pattern, with fine-grained
grading generally less portable than localisation or detection
\cite{mcsweeney2023external,zhang2023deep}. A model intended for eventual use in
Iraq or the wider region therefore cannot be justified solely by performance on a
multinational benchmark. The model must be transported to a genuinely independent
hospital cohort, its zero-shot degradation measured before local tuning, and the
amount of local annotation required to recover performance quantified explicitly.
That annotation--performance curve is more actionable than a single statement that
``domain adaptation improved accuracy''.

\subsection{Why the Problem Should Be Addressed Now}
\label{sec:intro-why-now}

Three developments make the present study timely. First, LumbarDISC supplies a
large, multi-institutional and anatomically structured benchmark whose target schema
is rich enough to test relational modelling rather than generic image
classification \cite{richards2026the}. Second, public anatomical resources such as
SPIDER reduce the need to claim localisation or segmentation as the central
contribution \cite{vandergraaf2024lumbar}. Third, the current methodological
literature is sufficiently mature to provide strong controls: multi-view
cross-attention, anatomy-guided contextual models, contemporary CNN and transformer
backbones, self-supervised pretraining, calibration and parameter-efficient transfer
all exist as comparators or supporting tools. The research can therefore ask a more
precise question than whether a new architecture attains a high score.

The resulting opportunity is methodological rather than merely computational. The
field now has enough data and enough baseline competence to test whether explicit
clinical structure provides value over capacity alone. If a typed target graph,
anatomical self-supervision or disease-conditioned router fails to outperform its
control, that negative result is still informative. Conversely, if a gain disappears
when topology is shuffled, when a simpler transformer is matched for capacity, or
when the model is transported to another hospital, the interpretation must change.
This insistence on falsifiable components is part of the motivation for the study,
not an afterthought added to evaluation.

% ===========================================================================
\section{Problem Statement}
\label{sec:intro-problem}

Despite rapid progress in automated lumbar MRI analysis, four connected
methodological problems remain.

\subsection{Problem 1: The Output Space Is Structured but Commonly Treated as Independent}
\label{sec:intro-problem-relational}

The benchmark output space contains ordered lumbar levels, multiple conditions at
each level and bilateral counterparts. These targets arise from one connected
anatomical system, yet the dominant design predicts them independently after a
shared encoder. This separation is convenient for implementation but does not test
whether adjacent-level context, within-level co-occurrence or bilateral relations
carry predictive information. Existing graph work in lumbar MRI does not resolve
this question because its graph is constructed within a disc rather than across the
disease targets \cite{baur2024automated}. CrossSpine introduces anatomical priors
and level-specific behaviour but still does not formulate the full 25 target outputs
as a typed communication graph \cite{nguyen2026crossspine}. The first problem is
therefore not the absence of contextual modelling in general; it is the absence of a
controlled test of explicit disease--anatomy relations among the prediction targets.

A relational model also carries a specific risk: useful context can become
``information contagion''. A severe lesion at one level should not cause the model
to elevate an adjacent normal level simply because degeneration often co-occurs.
Any graph proposed as a solution must therefore preserve local visual evidence,
include non-relational and homogeneous controls, and be stress-tested on naturally
isolated lesions. Without those controls, improved aggregate performance would not
establish that the anatomical topology is beneficial.

\subsection{Problem 2: Sequence Fusion Is Usually Fixed While Evidence Is Target- and Quality-Dependent}
\label{sec:intro-problem-routing}

Sagittal T1, sagittal T2/STIR and axial T2 are complementary acquisitions. The
foramina, central canal and lateral recess are not equally well represented in each
view, and routine data may be incomplete or technically degraded. Strong recent
systems address multi-view fusion through cross-attention
\cite{batra2025mscan,nguyen2026crossspine}, but a remaining deployment question is
whether one model can dynamically allocate evidence by condition, level, patient and
sequence quality. A model trained only on complete, high-quality benchmark studies
may fail when one sequence is missing, truncated or motion-corrupted. Missingness
and corruption are also different problems: an absent sequence can be masked,
whereas a degraded sequence remains present and can attract misleading attention.
The second problem is therefore to learn adaptive sequence allocation without
turning the router into an unexamined source of instability.

\subsection{Problem 3: Cross-Sequence Correspondence Is Available but Underused as Supervision}
\label{sec:intro-problem-ssl}

Self-supervised learning can reduce dependence on costly expert labels, but common
contrastive objectives construct positive pairs through augmentations of the same
image. Multi-sequence MRI contains a more domain-specific relationship: different
acquisitions can depict the same anatomical level even when their intensity pattern,
orientation and diagnostic role differ. DICOM patient-space geometry provides the
physical correspondence required to identify those pairs. Prior spinal
self-supervision demonstrates the value of metadata-derived or longitudinal
relationships \cite{jamaludin2017selfsupervised}, but anatomical cross-sequence
correspondence has not been established as the supervision signal for the target
formulation studied here.

The problem is not solved simply by forcing all sequence representations to become
identical. T1 and T2 contrasts contain sequence-specific diagnostic information, so
excessive invariance could improve embedding alignment while damaging grading. The
study must therefore separate the self-supervised projection space from downstream
features and test explicitly whether sequence-specific information survives
pretraining.

\subsection{Problem 4: Internal Accuracy Does Not Establish Transportability or Clinical Error Alignment}
\label{sec:intro-problem-transfer}

The strongest evidence against uncritical internal benchmarking is the repeated
performance decline under external validation
\cite{compte2023are,mcsweeney2023external,zhang2023deep}. Scanner manufacturer,
field strength, resolution, axial prescription, preprocessing and patient mix can
all change. More subtly, the reference standard may change: a public benchmark may
use structured expert annotation while a hospital's routine reports reflect local
reporting habits and may omit findings that were not considered clinically salient.
An external performance drop can therefore combine imaging-domain shift,
localisation failure and reference-standard shift.

A defensible transfer study must measure zero-shot performance before local tuning,
keep the local test set fixed, separate it from the adaptation pool, and harmonise
the local reference standard as far as resources allow. The present work therefore
uses the Rizgary Hospital cohort for external evaluation rather than for iterative
architecture development. The principal local comparison is restricted to central
canal stenosis, for which a defensible overlap with the public benchmark can be
created; local disc bulge, protrusion and extrusion remain a separate morphology
task rather than being relabelled as stenosis.

The error objective also requires care. Severity grades are ordinal, and distant
under-grading is clinically different from adjacent disagreement. Yet the literature
does not establish that ordinal losses automatically improve classification:
Niemeyer et al. found several ordinal formulations unable to outperform standard
cross-entropy for Pfirrmann grading \cite{niemeyer2021a}. A recent preprint has begun to evaluate distance-aware ordinal assessment on
LumbarDISC \cite{trinc2026beyond}; however, the combination of asymmetric clinical
cost with target-level relational modelling appears unaddressed in the literature
reviewed for this thesis.
At the same time, reducing
Severe$\rightarrow$Normal/Mild errors by simply shifting the model toward predicting
Severe would create an unacceptable false-positive burden. The fourth problem is
therefore to evaluate error direction, calibration and selective prediction without
converting radiological grading into an unsupported claim of treatment utility.

\subsection{Integrated Problem Statement}
\label{sec:intro-integrated-problem}

Taken together, the research problem can be stated as follows:

\begin{quote}
Current automated lumbar MRI grading systems can localise anatomy and predict
per-level disease severity with increasingly strong internal performance, but the
prevailing formulation does not fully exploit the relational structure of the
25-target output space, does not explicitly learn target- and quality-dependent
sequence allocation, underuses DICOM-defined cross-sequence anatomical
correspondence as a self-supervised signal, and provides limited evidence on how
these design choices behave under true institutional transfer. A controlled
methodology is therefore required to determine whether explicitly modelling these
properties improves grading, robustness, label efficiency and portability beyond
strong matched baselines, while preserving local evidence and respecting the
uncertainty of the clinical reference standard.
\end{quote}

\laterinsert{If implementation reveals a previously unrecognised but reproducible
problem that materially changes the framing---for example a systematic DICOM
geometry failure mode across a substantial cohort---add a short factual paragraph
here and document the protocol amendment. Do not add a new problem merely because a
proposed contribution underperformed.}

% ===========================================================================
\section{Research Gap and Novelty Boundary}
\label{sec:intro-gap}

Chapter~\ref{ch:litreview} develops the evidence in detail; the purpose here is to
state the gap concisely enough to define the thesis.

\subsection{Gap A: Explicit Relational Modelling of Disease--Anatomy Targets}
\label{sec:intro-gap-graph}

The reviewed literature contains multi-task systems, inter-level attention,
anatomical priors and one lumbar graph model, but no reviewed study constructs the
benchmark's level--condition--laterality targets as a heterogeneous graph with typed
relations for adjacent levels, within-level condition coexistence and bilateral
symmetry. Baur et al. use a graph within one reconstructed disc
\cite{baur2024automated}; Nguyen et al. introduce level-specific priors and
cross-sequence attention without target-to-target graph communication
\cite{nguyen2026crossspine}; Chai et al. demonstrate anatomy-guided contextual
multi-level modelling without the disease-target graph formulation proposed here
\cite{chai2026anatomy}. The intended contribution is therefore not ``graph neural
networks for lumbar MRI'' but a controlled test of an explicit typed graph over the
prediction targets themselves.

\subsection{Gap B: Anatomically Aligned Self-Supervision and Disease-Conditioned Evidence Routing}
\label{sec:intro-gap-sequence}

The literature establishes both the complementarity of MRI sequences and the value
of self-supervised pretraining, but these strands remain only partially connected in
lumbar grading. Anatomical correspondence across sequences is usually treated as a
preprocessing requirement for fusion rather than as supervision. At the same time,
multisequence systems generally learn a common fusion mechanism rather than a
condition- and quality-dependent allocation that is deliberately challenged under
missing and corrupted modalities. The thesis therefore tests two linked ideas: that
DICOM geometry can define cross-sequence positive relationships, and that a single
model can use the resulting features through a target-conditioned router whose
failure modes are directly measurable.

\subsection{Gap C: Quantified Transfer and Adaptation Under a Harmonised External Task}
\label{sec:intro-gap-transfer}

Existing external validations establish that transport is difficult, but they
usually report one transported model at one point in the adaptation process---often
with no adaptation at all. The operational question for a receiving hospital is
instead how performance changes before adaptation and how quickly it recovers as
local labels are supplied. The literature reviewed for this thesis contains no
Middle Eastern evaluation of the same scope and no quantified parameter-efficient
annotation curve for this lumbar target formulation. The gap is therefore not
simply ``lack of external validation''; it is lack of a staged zero-shot-to-few-shot
transfer experiment with fixed local testing, controlled annotation budgets and
explicit reference-standard harmonisation.

\subsection{Supporting Gap: Directional Error and Calibrated Uncertainty}
\label{sec:intro-gap-error}

Ordinal severity and asymmetric error are clinically meaningful, but ordinal
objectives have not consistently outperformed categorical baselines
\cite{niemeyer2021a}. Calibration is also rarely reported in lumbar MRI despite the
known overconfidence of modern neural networks \cite{guo2017calibration}. These are
therefore treated as supporting questions rather than headline novelty. A
cost-sensitive objective is retained only if it reduces clinically consequential
under-grading without creating indiscriminate Severe false positives, and selective
prediction is evaluated as a model behaviour rather than proof of a validated
clinical workflow.

The novelty claim throughout the thesis will remain appropriately bounded. The
phrasing ``to our knowledge'' or ``in the reviewed literature'' is used where
necessary, and the final claim will be updated against the literature available at
the time of submission. A new paper that adopts one element of the proposed system
does not invalidate the PhD if the controlled experimental question remains
distinct; it does, however, change how originality must be described.

\laterinsert{Immediately before thesis submission, update this novelty boundary
against publications appearing after the current literature-search cut-off. Record
which claims remain original and which should be reframed as independent validation,
extension or comparative evidence.}

% ===========================================================================
\section{Research Aim}
\label{sec:intro-aim}

The overall aim of this PhD is:

\begin{quote}
\emph{to determine whether explicitly modelling anatomical relationships among
lumbar disease targets, anatomically aligned evidence across MRI sequences, and
target-dependent sequence availability improves the accuracy, robustness, label
efficiency and cross-institutional portability of automated lumbar spine MRI
severity grading, relative to strong non-relational and fixed-fusion baselines.}
\end{quote}

The aim is intentionally phrased as a determination rather than a promise of
superiority. The thesis succeeds scientifically if it establishes under controlled
conditions where these mechanisms help, where they fail, and what their limits are.

% ===========================================================================
\section{Research Objectives}
\label{sec:intro-objectives}

The aim is operationalised through the following objectives.

\begin{enumerate}
  \item To construct a reproducible patient-level lumbar MRI analysis pipeline that
        preserves DICOM patient-space geometry, performs reliable level
        localisation, and supplies identical anatomically defined regions of
        interest to all grading baselines and proposed models.

  \item To establish competitive independent and contextual grading baselines on
        the public LumbarDISC development data, including fixed fusion,
        cross-attention and ordered inter-level context where reproducibility
        permits, so that improvements cannot be attributed to a weak comparator.

  \item To develop and evaluate anatomically aligned cross-sequence
        self-supervision in which positive relationships are defined by the same
        patient and lumbar level across MRI sequences using DICOM geometry.

  \item To quantify whether this anatomical self-supervision improves supervised
        grading under 10\%, 25\%, 50\% and 100\% label availability and whether it
        preserves sequence-specific diagnostic information.

  \item To develop and evaluate disease-conditioned adaptive sequence routing that
        learns patient- and target-dependent weights over available MRI evidence
        while explicitly representing sequence availability and acquisition quality.

  \item To test the router under both missing-sequence and corrupted-sequence
        conditions, and to compare its behaviour with concatenation, averaging and
        ordinary cross-attention under matched inputs and training budgets.

  \item To construct a heterogeneous disease--anatomy graph over the 25
        level--condition--laterality targets, with typed adjacent-level,
        same-level cross-condition and bilateral relations, and to compare it with
        independent heads, an ordered level transformer, a homogeneous graph,
        ungated message passing and shuffled topology.

  \item To test whether relational reasoning introduces adjacent-level false
        positives in naturally isolated lesions and to retain the graph only where
        any performance or robustness benefit is not purchased by clinically
        incoherent information propagation.

  \item To compare categorical, ordinal and clinically asymmetric cost-sensitive
        grading objectives, together with probability calibration and uncertainty
        analysis, while explicitly reporting both severe under-grading and severe
        over-grading.

  \item To quantify zero-shot transfer from the multinational public benchmark to
        the independent Rizgary Hospital cohort on the shared central-canal grading
        task, separating localisation, image-domain and reference-standard effects
        as far as the retrospective data allow.

  \item To estimate an annotation--performance recovery curve using controlled
        local adaptation sizes $N\in\{10,25,50,100\}$ and to compare where a limited
        trainable-parameter budget is most effectively placed.

  \item To report all principal comparisons with patient-level uncertainty,
        repeated training where feasible, contribution-specific ablations, negative
        results and explicit methodological limitations.
\end{enumerate}

These objectives define a staged experiment rather than a monolithic architecture.
A component is not retained because it appears in the proposal; it is retained only
if the corresponding experiment demonstrates a reproducible contribution or a
meaningful robustness, calibration or interpretability advantage.

% ===========================================================================
\section{Research Questions}
\label{sec:intro-rqs}

The research questions are fixed prospectively and correspond directly to the
methodological tests in Chapter~\ref{ch:methodology}.

\begin{description}
  \item[RQ1 --- Relational disease modelling.] Does representing lumbar disease
  targets as a typed heterogeneous graph improve grading over independent target
  heads, an ordered five-level transformer and a homogeneous graph, and are any
  gains concentrated where single-target evidence is weakest?

  \item[RQ2 --- Anatomically aligned self-supervision.] Does cross-sequence
  pretraining based on DICOM-defined anatomical correspondence improve grading,
  label efficiency and transfer compared with ImageNet initialisation, generic
  medical pretraining and ordinary augmentation-based contrastive learning?

  \item[RQ3 --- Disease-conditioned modality routing.] Can one model learn
  target- and patient-dependent sequence weights while remaining robust to both
  missing and quality-degraded MRI sequences, without requiring a separately
  trained network for every combination of sagittal T1, sagittal T2/STIR and axial
  T2?

  \item[RQ4 --- Cross-institutional transfer.] What is the zero-shot degradation
  when a benchmark-trained model is applied to a previously unseen Middle Eastern
  clinical cohort, how much of that degradation is attributable to localisation
  versus grading/reference-standard shift, and how does performance recover as the
  amount of local annotation available for adaptation increases?

  \item[RQ5 --- Clinically asymmetric error and uncertainty.] Do ordinal and
  cost-sensitive objectives reduce clinically consequential distant errors,
  particularly Severe$\rightarrow$Normal/Mild, and can calibrated uncertainty
  support selective prediction without overstating clinical safety?
\end{description}

These questions are intentionally answerable by negative as well as positive
results. If a typed graph performs no better than a homogeneous graph, RQ1 is still
answered. If anatomical self-supervision improves label efficiency but not external
transfer, RQ2 yields a more limited but still informative result. This design avoids
defining doctoral success as the requirement that every proposed component win.

% ===========================================================================
\section{Research Hypotheses}
\label{sec:intro-hypotheses}

The study uses directional hypotheses where the literature justifies a prediction,
but avoids arbitrary target accuracies that would have no defensible basis before
implementation.

\begin{description}
  \item[H1 --- Anatomical self-supervision.] Anatomical cross-sequence
  self-supervision will improve label-efficient grading and cross-institutional
  robustness relative to generic pretraining, particularly when only a fraction of
  benchmark labels is used.

  \item[H2 --- Adaptive routing.] Disease-conditioned routing with explicit
  availability masking, modality dropout and quality-aware gating will degrade less
  under missing-sequence and corruption experiments than fixed concatenation or
  ordinary cross-attention trained only on clean complete studies.

  \item[H3 --- Heterogeneous relational structure.] A typed heterogeneous graph
  will outperform independent target heads and a homogeneous graph if the proposed
  anatomical edge types carry useful information; performance under a shuffled-edge
  control should deteriorate if the topology itself matters.

  \item[H4 --- External degradation.] Zero-shot performance on the Rizgary Hospital
  cohort will be lower than internal benchmark performance because of scanner,
  protocol, population, localisation and reference-standard shift; the magnitude
  and pattern of that degradation are treated as findings rather than as failure of
  the research programme.

  \item[H5 --- Annotation-dependent recovery.] Local adaptation will improve
  performance as annotation increases, but no predetermined percentage of recovery
  is required. The shape of the recovery curve and the location in the network at
  which limited adaptation is most effective are empirical outcomes.

  \item[H6 --- Asymmetric grading error.] Cost-sensitive and ordinal objectives may
  reduce distant Severe$\rightarrow$Normal/Mild errors even if they do not improve
  aggregate macro-F1; any reduction must be evaluated against false-positive Severe
  predictions and against the cross-entropy baseline
  \cite{niemeyer2021a}.
\end{description}

The hypotheses will not be rewritten after inspection of the held-out public or
Rizgary test results. If an implementation decision changes the operational test of
a hypothesis, that change will be entered into the protocol changelog with its date
and rationale.

% ===========================================================================
\section{Conceptual Framework}
\label{sec:intro-conceptual}

The thesis can be understood as a progression from \emph{local evidence} to
\emph{structured evidence} and finally to \emph{transported evidence}.

At the first level, a model receives anatomically localised image evidence for a
single level and target. This is the baseline condition: if local evidence alone is
sufficient, more complex mechanisms should not be retained. At the second level,
the model is permitted to combine complementary MRI sequences, but the combination
must respect anatomical correspondence and sequence availability. Anatomical
self-supervision asks whether the same structure under different acquisitions can
produce a more label-efficient representation, while disease-conditioned routing
asks how strongly each sequence should contribute to a given target.

At the third level, a target is no longer treated as an isolated endpoint. Its
representation becomes a node in a typed graph whose relations correspond to three
known structures of the lumbar grading problem: the ordered chain of levels,
coexisting compartments at a level, and bilateral pairing. Message passing is soft
rather than deterministic; no rule states that disease at one level forces disease
at its neighbour. The graph is therefore a hypothesis about useful context, not a
hard-coded disease progression model.

At the fourth level, predictions are evaluated as ordinal and probabilistic outputs.
The model is judged not only by aggregate discrimination but by error distance,
Severe-class performance, calibration and the false-positive cost of any
cost-sensitive objective. This stage asks whether a statistically improved model is
also better aligned with the shape of clinically consequential grading error.

The final level is institutional transport. The representation, router, graph and
output head developed on public data are frozen before external evaluation. The
Rizgary Hospital cohort then asks a different question from internal validation:
what remains true when the acquisition and reporting environment changes? Only
after zero-shot performance has been measured is local adaptation permitted. This
ordering prevents the local test cohort from becoming an informal development set.

Conceptually, the thesis therefore tests the proposition

\begin{equation}
\text{robust grading}
\;=\;
f(\text{local anatomy},\;\text{cross-sequence evidence},\;
\text{target relations},\;\text{uncertainty},\;\text{domain}),
\label{eq:intro-conceptual}
\end{equation}

where the equation is explanatory rather than a literal computational model. Each
term corresponds to a controlled experiment, and no term is assumed necessary until
its ablation supports that conclusion.

% ===========================================================================
\section{Significance of the Study}
\label{sec:intro-significance}

\subsection{Scientific Significance}
\label{sec:intro-significance-scientific}

The principal scientific significance lies in changing the unit of reasoning from
an isolated image label to a structured set of disease targets. Many improvements in
medical imaging can be produced by additional parameters, larger backbones or more
aggressive tuning. Such gains are difficult to interpret mechanistically. By
contrast, this thesis attaches each proposed contribution to a specific hypothesis
and a destructive control: anatomical edges are compared with homogeneous and
shuffled topology; adaptive routing is compared with fixed fusion and challenged by
missing and corrupted modalities; anatomical self-supervision is compared with
generic pretraining and tested for over-invariance. The resulting evidence can show
not merely whether the full model scored higher, but \emph{why} it did or did not.

This matters because recent lumbar literature already demonstrates sophisticated
contextual models \cite{chai2026anatomy,batra2025mscan,nguyen2026crossspine}. The
research contribution cannot therefore be the generic use of attention, graphs or
multiple sequences. Its value lies in isolating the anatomical assumptions and
measuring them against alternatives under one controlled pipeline.

\subsection{Methodological Significance}
\label{sec:intro-significance-methodological}

The study places unusually strong emphasis on validity controls. Patient identity,
not slice or target, defines the data split; self-supervised learning is prohibited
from seeing held-out patients; localisation is evaluated separately from grading;
zero-shot transfer is measured before adaptation; the external reference standard is
kept distinct from routine report extraction; and every claimed model contribution
requires an ablation. These choices respond directly to recurring weaknesses in the
medical-imaging literature, where leakage, non-comparable splits and weak external
validation can produce optimistic performance estimates
\cite{compte2023are,wang2024machine}.

The transfer design is methodologically significant for a second reason. The local
hospital cohort is not used simply as another test set. The experiment attempts to
disentangle localisation failure, image-domain shift and reference-standard shift,
then measures recovery as local labels are added. This converts ``generalisation''
from a binary claim into a measurable process.

\subsection{Translational and Regional Significance}
\label{sec:intro-significance-regional}

The regional significance follows from the use of an independent Iraqi hospital
cohort that was not curated to mirror the public benchmark. Models developed on
large international datasets are often assumed to transfer into settings whose
scanner mix, acquisition practice, disease prevalence and reporting conventions were
not represented in development. Testing that assumption is particularly important
for institutions that do not have the resources to reproduce a large annotation
campaign locally. If useful performance can be recovered from tens rather than
hundreds of local annotations, that is a practical finding; if not, the negative
result is equally important because it quantifies the annotation burden required for
responsible transfer.

The study does not claim that one external hospital represents all of Iraq or the
Middle East. Rizgary Hospital provides evidence of transport to one genuinely
independent setting. Regional generalisation would require additional hospitals and
is therefore a future study rather than an inference drawn from one centre.

\subsection{Clinical Significance and its Boundary}
\label{sec:intro-significance-clinical}

The intended clinical significance is improved technical support for structured MRI
grading: consistency, reproducibility, calibrated confidence and explicit handling
of incomplete imaging. The thesis does not claim that this will improve treatment
outcomes, referral decisions or radiologist productivity. Those claims require a
prospective reader-assistance or impact study, as demonstrated by work that measures
reader performance directly \cite{lim2022improved}. The present study therefore
uses language such as ``decision-support potential'' and ``technical portability'',
not ``clinical benefit''.

% Written after Chapter 4. States only what the executed experiments support.
Of the benefits anticipated above, the executed experiments support
\emph{discrimination} and, within it, the reduction of clinically distant
errors. Quadratic weighted kappa improves by $+0.0177$ over a single-sequence
baseline, Severe$\rightarrow$Normal confusions fall from $5.7\%$ to $3.8\%$ and
recall on Severe targets rises from $61.1\%$ to $63.1\%$, the last two
attributable to the ordinal and cost-sensitive objective rather than to any
proposed structural mechanism.

Robustness to missing sequences improves, but the improvement is attributable to
modality dropout rather than to disease-conditioned routing, which changes
sequence reliance by an amount indistinguishable from zero. Label efficiency and
cross-institutional transfer were not evaluated and no claim is made for either.
Calibration is reported but was not the object of a confirmatory comparison.

No inference is drawn from these figures about patient benefit. The study
measures agreement with a reference standard on a retrospective benchmark; it
does not measure any clinical outcome.

% ===========================================================================
\section{Proposed Contributions}
\label{sec:intro-contributions}

At the protocol stage, the contributions below collectively form the proposed \textbf{AMOG-Net} framework and are \emph{proposed contributions}.
They become demonstrated thesis contributions only if the experiments support them.

\subsection{Contribution 1: Anatomically Aligned Cross-Sequence Self-Supervision}
\label{sec:intro-contrib-ssl}

The first proposed contribution uses DICOM-defined patient-space correspondence to
construct positive relationships between different MRI sequences depicting the same
patient and lumbar level. The objective is not to invent a new generic contrastive
loss; established contrastive or non-contrastive objectives may be used. The
contribution under test is the anatomical definition of correspondence and its value
for label efficiency, grading and transport. A dedicated preservation experiment
checks whether alignment erases sequence-specific diagnostic information.

\subsection{Contribution 2: Disease-Conditioned and Quality-Aware Sequence Routing}
\label{sec:intro-contrib-routing}

The second proposed contribution allows one study to be fused differently for
different disease targets. Routing depends on target identity, patient features,
sequence availability and quality evidence, and is evaluated under genuine or
simulated missingness and corruption. The contribution is not the existence of a
softmax gate or mixture-of-experts mechanism; it is the target-conditioned allocation
of clinically complementary evidence and the explicit robustness evaluation that
accompanies it.

\subsection{Contribution 3: Heterogeneous Disease--Anatomy Graph Reasoning}
\label{sec:intro-contrib-graph}

The third proposed contribution represents the 25 benchmark outputs as nodes in a
heterogeneous graph. Typed edges encode adjacent-level, within-level cross-condition
and bilateral relations. The graph includes a mechanism for preserving local
evidence when neighbour messages are unhelpful and is evaluated against independent,
ordered-transformer, homogeneous, ungated and shuffled-topology controls. A focal
isolated-lesion test determines whether relational reasoning improves uncertain
predictions without artificially spreading disease to adjacent normal levels.

\subsection{Contribution 4: A Structured Cross-Institutional Transfer Evaluation}
\label{sec:intro-contrib-transfer}

The translational contribution is an experimental design rather than a new network
layer. The selected public-data model is evaluated zero-shot on a fixed Rizgary
Hospital test set before any local adaptation. A separate local adaptation pool is
then sampled at controlled annotation sizes and used to compare constrained update
strategies. The analysis is strengthened by a verified-localisation control and,
where resources permit, harmonised independent reader grading of central canal
stenosis. This design estimates both the cost of transport and the cost of recovery.

\subsection{Supporting Contributions}
\label{sec:intro-contrib-supporting}

Ordinal and cost-sensitive objectives, calibration, uncertainty estimation and
parameter-efficient adaptation support the principal contributions but are not
claimed as original methods. Their value lies in testing whether the proposed
structured model behaves sensibly under clinically relevant error definitions and
limited local annotation. They are retained only when their contribution is
measurable.

% --------------------------------------------------------------------------
% Written at final submission, replacing the [LATER INTEGRATION] placeholder.
% Only components supported by the locked experiments are claimed. Components
% that failed are reported below as negative findings and remain in the
% research history and in the ablation ladder.
% --------------------------------------------------------------------------
\subsection{Demonstrated Contributions}
\label{sec:intro-contrib-demonstrated}

The experiments reported in Chapter~\ref{ch:results} support three
contributions. They are not the three proposed above, and the difference is
itself part of what the thesis contributes.

\paragraph{D1 --- A verified anatomy-aware multi-sequence grading system.}
AMOG-Net improves quadratic weighted kappa over a single-sequence baseline by
$+0.0177$ (95\% CI $[+0.0097, +0.0260]$), on every one of seven training seeds
and surviving false discovery rate correction across the pre-specified
comparison family. Two further comparisons also survive: typed heterogeneous
edges and the ordinal, cost-sensitive objective. The system, its frozen data
partition and its behavioural test suite are released.

\paragraph{D2 --- A controlled ablation protocol for anatomical priors.} Each
claimed mechanism is tested against a control matched in capacity rather than
against its own absence: a degree-preserving edge shuffle for graph topology, a
router-free fixed-fusion arm for sequence routing, an architecture-matched
untrained network as an attribution floor, and a label-shuffled model as a
pipeline-level negative control. Section~\ref{sec:results-threats} documents
three analyses that produced incorrect positive conclusions before the
corresponding control was introduced, in each case in the direction of the
hypothesis under test.

\paragraph{D3 --- Evidence that anatomical structural priors do not improve
grading at this data scale, with a mechanism.} Neither anatomically aligned
cross-sequence self-supervision nor disease-conditioned routing produces a
detectable improvement, and anatomically correct graph topology does not
outperform a degree-preserving random shuffle. These are bounds rather than
absences: the paired intervals constrain the three mechanisms to at most
$1.32\%$, $1.55\%$ and $1.05\%$ of baseline performance respectively. Gradient-weighted class activation
mapping shows why: the convolutional encoder already concentrates $2.4$ times the
chance share of its attention on the annotated target before any prior is added.
The localisation these priors were designed to supply is largely accomplished
upstream of them.

\subsection{Proposed Contributions Not Supported}
\label{sec:intro-contrib-negative}

Contribution~1 (anatomical cross-sequence self-supervision) is rejected on
accuracy, on cross-sequence robustness under controlled input ablation, and on
attention concentration. The pretext task is learnable and was learned ---
validation InfoNCE $1.0771$ against a chance value of $3.4657$ --- but the
representation confers no measurable downstream benefit.

Contribution~2 (disease-conditioned routing) divides. The router does learn
clinically sensible target-dependent sequence weights, reproducing the expected
allocation in all fifteen runs in which a gate is present. It does not improve
grading ($+0.0053$ QWK, five seeds of seven, bounded at $1.55\%$ of baseline)
and it does not improve robustness to missing sequences.
Its allocation tracks the annotation structure of the benchmark rather than
creating a dependence of its own.

Contribution~3 (heterogeneous disease--anatomy graph) is partly supported. Typed
relational structure improves grading over a homogeneous graph ($+0.0093$ QWK on
every one of seven seeds, surviving correction at $p_{\mathrm{FDR}} = 0.009$);
the anatomical identity of the relations does not. The resulting
claim is narrower than proposed and more specific than the original: relational
typing carries information, anatomical adjacency does not.

Contribution~4 (cross-institutional transfer evaluation) was not executed. Cohort
preparation is complete and is reported in Section~\ref{sec:results-rq4};
adjudication of conflicting reference labels and a correct de-identification
procedure remain outstanding. It is retained here as unexecuted rather than
withdrawn.

% ===========================================================================
\section{Scope and Delimitations}
\label{sec:intro-scope}

Clear boundaries are necessary because lumbar MRI supports a much broader clinical
and computational agenda than one PhD can test.

\subsection{Clinical Scope}
\label{sec:intro-scope-clinical}

The primary public task is degenerative lumbar disease severity grading at five
levels from L1--L2 through L5--S1. On LumbarDISC, the full target space contains
central canal stenosis, left and right neural foraminal narrowing, and left and
right subarticular stenosis, each graded as Normal/Mild, Moderate or Severe
\cite{richards2026the}. The study does not attempt to diagnose every cause of low
back pain, infer symptoms, determine surgical candidacy or predict long-term patient
outcomes.

The primary local transfer task is narrower. Rizgary Hospital reports support a
defensible central-canal comparison more readily than the complete benchmark schema.
The confirmatory external analysis is therefore restricted to central canal stenosis
at the five lumbar levels unless new expert grading expands the local reference
standard. Disc bulge, protrusion and extrusion are retained as a separate local
morphology task and are not heuristically converted into stenosis grades.

\subsection{Imaging Scope}
\label{sec:intro-scope-imaging}

The principal imaging inputs are sagittal T1, sagittal T2/STIR and axial T2 MRI.
Other sequences may appear in clinical studies but do not define the primary model
unless a justified later protocol amendment is made. DICOM geometry is preserved
because cross-plane correspondence is part of the scientific design. PNG or other
rendered formats may be used for visual inspection but do not replace the geometric
metadata required for alignment.

\subsection{Methodological Scope}
\label{sec:intro-scope-method}

Anatomical localisation and segmentation are enabling technologies. They are
measured and controlled because a poor crop can invalidate grading, but a new
localisation architecture is not the headline doctoral contribution. Likewise,
backbone selection is controlled rather than treated as a novelty competition. The
research may use ResNet-, EfficientNet-, Swin- or other contemporary encoders, but
the primary backbone is frozen before the main contribution-specific ablations.

The graph topology is defined anatomically rather than discovered freely from the
test data. Adjacent levels, same-level conditions and bilateral counterparts form the
pre-specified relation families. Learning the globally optimal topology is a
different research problem and is excluded from the primary scope.

\subsection{Data and Institutional Scope}
\label{sec:intro-scope-data}

Model development uses the public LumbarDISC benchmark and, where permitted, SPIDER
for anatomical localisation or segmentation support. The locally held labelled
competition subset contains approximately 1,975 studies at the protocol stage;
\torecord{final public development, validation and held-out counts after quality
control}. The local project inventory currently identifies 299 narrative reports
and 294 candidate matched multi-sequence cases, while a wider raw DICOM inventory
contains additional studies. These numbers are not presented as the final cohort
until one-to-one reconciliation is complete:
\torecord{final Rizgary case-flow counts and reasons for exclusion}.

Before analytical use of the local cohort, the Rizgary data-preparation phase will
follow the approved institutional governance protocol. A controlled research copy will
be de-identified before access by the model-development team, and the permitted storage
and compute environment will be verified before local processing begins. These safeguards
precede, rather than form part of, model development.

Rizgary Hospital is a single external institution. The study can therefore support a
claim of transport to that independent setting, not universal Middle Eastern
performance. Additional regional hospitals would be required for a multicentre
regional generalisation claim.

\subsection{Ethical and Governance Scope}
\label{sec:intro-scope-ethics}

The local component is retrospective and does not alter patient care. Formal use of
the local data remains conditional on the relevant institutional permissions:
\toconfirm{hospital/ethics approval reference, date and responsible institution};
\toconfirm{retrospective consent or waiver basis}; and
\toconfirm{permitted compute environment, including whether external cloud/GPU
processing is authorised}. Identifiable raw DICOM is not supplied to students or
external compute; a controlled de-identified research copy is required before local
analysis.

\subsection{Evaluation Scope}
\label{sec:intro-scope-evaluation}

The thesis evaluates technical grading performance, robustness, calibration,
computational cost and cross-institutional portability. The unit of statistical
independence is the patient. It does not claim prospective clinical impact,
radiologist replacement or treatment benefit. Reader-assistance evaluation is a
logical future study if the technical system proves sufficiently reliable.

% ===========================================================================
\section{Assumptions and Operational Definitions}
\label{sec:intro-definitions}

The following terms are used consistently throughout the thesis.

\begin{longtable}{p{3.6cm}p{10.0cm}}
\toprule
\textbf{Term} & \textbf{Operational meaning in this thesis}\\
\midrule
\endfirsthead
\toprule
\textbf{Term} & \textbf{Operational meaning in this thesis}\\
\midrule
\endhead
Lumbar level & One of L1--L2, L2--L3, L3--L4, L4--L5 or L5--S1.\\
Target & A level--condition--laterality output on LumbarDISC; 25 targets per study at full benchmark scope.\\
Severity grade & Normal/Mild, Moderate or Severe, encoded as an ordered three-class outcome.\\
Local evidence & Image-derived representation for one target before graph message passing.\\
Anatomical correspondence & Physical correspondence between acquisitions defined through DICOM patient-space geometry and lumbar localisation, not through image appearance alone.\\
Sequence routing & Learned allocation of weight over available MRI sequence features for a specific patient and target.\\
Missing modality & An expected MRI sequence is absent and explicitly unavailable to the model.\\
Degraded modality & A sequence is present but technically compromised, for example by motion, truncation, noise, slice loss or poor coverage.\\
Heterogeneous graph & A graph whose nodes are disease--anatomy targets and whose edges belong to clinically defined relation types.\\
Zero-shot transfer & Evaluation of the frozen public-data model on the fixed local test set before local fine-tuning, threshold selection or local calibration.\\
Few-shot adaptation & Model adaptation using only the separate local adaptation pool under pre-specified annotation budgets.\\
Reference-standard shift & Difference between the process used to generate labels in the benchmark and the process used to establish the local external reference.\\
Selective prediction & Model abstention or referral based on uncertainty; this is a technical behaviour, not evidence of a validated clinical workflow.\\
\bottomrule
\end{longtable}

These definitions are intentionally narrower than ordinary-language usage. In
particular, ``external validation'' does not mean that a model has been proven safe
for deployment, and ``uncertainty'' does not mean that the network has measured a
patient's clinical risk.

% ===========================================================================
\section{Overview of the Research Design}
\label{sec:intro-design-overview}

The research is a retrospective, multi-dataset medical-imaging methods study. Its
logic is cumulative, but the conclusions remain component-specific.

\begin{enumerate}
  \item \textbf{Data and anatomy foundation.} Public data are inventoried,
        partitioned by patient, reconstructed in DICOM patient space, anatomically
        localised and converted into disease-specific 2.5D regions of interest.

  \item \textbf{Strong baseline establishment.} An independent ROI classifier and
        competitive multiview/context baselines establish what can be achieved
        without the proposed relational formulation.

  \item \textbf{Representation experiment.} Anatomically aligned cross-sequence
        self-supervision is evaluated against generic and medical pretraining under
        several label fractions and sequence-specific preservation probes.

  \item \textbf{Fusion experiment.} Disease-conditioned sequence routing is tested
        against fixed and attention-based fusion under complete, missing and
        corrupted sequence conditions.

  \item \textbf{Relational experiment.} Homogeneous and heterogeneous target graphs
        are compared with independent and ordered-transformer controls, including
        shuffled topology and isolated-lesion stress testing.

  \item \textbf{Output and uncertainty experiment.} Categorical, ordinal and
        cost-sensitive objectives are evaluated together with calibration and
        uncertainty, with both under-grading and over-grading reported.

  \item \textbf{External transport experiment.} The selected public-data model is
        frozen and evaluated zero-shot on the fixed Rizgary Hospital test set. A
        verified-ROI control and harmonised reader reference are used where feasible
        to interpret the source of degradation.

  \item \textbf{Adaptation experiment.} Controlled local annotation budgets are
        sampled from a separate adaptation pool to estimate how quickly performance
        recovers and where a limited trainable-parameter budget is most effective.
\end{enumerate}

This order deliberately prevents the final architecture from becoming an indivisible
engineering object. Each stage introduces one assumption that can be rejected
without invalidating the remainder of the thesis. The final model is therefore the
simplest configuration supported by the evidence, not necessarily the configuration
containing every proposed module.

% ===========================================================================
\section{Expected Outcomes and Criteria for Scientific Success}
\label{sec:intro-success}

The protocol does not define success as reaching a predetermined accuracy. Such a
threshold would be difficult to justify because performance depends on target,
severity distribution, reference-standard reliability and the final patient split.
Scientific success is instead defined by whether the research questions can be
answered under controlled conditions.

For RQ1, success means establishing whether typed anatomical relations add value
beyond capacity-matched controls and whether they avoid adjacent-level
contamination. A null result is informative if the comparison is sufficiently
powered and the topology controls are valid. For RQ2, success means establishing
whether anatomical cross-sequence pairing improves label efficiency, transfer or
neither, while documenting whether sequence-specific information is preserved. For
RQ3, success means quantifying robustness under missing and degraded modalities and
determining whether adaptive routing adds value over simpler fusion. For RQ4,
success means producing an interpretable zero-shot transfer estimate and a controlled
adaptation curve, even if the external drop is large. For RQ5, success means
characterising the trade-off between severe under-grading, severe over-grading and
probability reliability rather than selecting a loss solely because it raises one
aggregate metric.

At final submission, the thesis will distinguish three categories explicitly:
\emph{supported contributions}, \emph{negative findings}, and \emph{unresolved
questions}. This is preferable to rewriting the original proposal so that only
successful components appear to have been planned.

% Written after Chapter 4. High-level answers only; metrics live in Chapter 4.
\begin{description}
  \item[RQ1 --- Relational disease modelling.] \textbf{Partly supported.} A typed
  heterogeneous graph improves grading over independent heads and over a
  homogeneous graph, on every seed. The anatomical identity of the relations does
  not contribute: an anatomically correct topology does not outperform a
  degree-preserving random shuffle.

  \item[RQ2 --- Anatomically aligned self-supervision.] \textbf{Not supported.}
  The pretext task is learnable and was learned, but the resulting representation
  yields no detectable improvement in grading, in robustness to withheld
  sequences, or in where the model attends.

  \item[RQ3 --- Disease-conditioned modality routing.] \textbf{Mechanism
  supported, benefit not.} The router reliably learns the clinically expected
  sequence weights, but a router-free model shows the same sequence dependence
  under intervention, and routing improves neither accuracy nor missing-sequence
  robustness.

  \item[RQ4 --- Cross-institutional transfer.] \textbf{Not executed.} Cohort
  preparation is complete; reference-label adjudication and de-identification
  remain outstanding.

  \item[RQ5 --- Clinically asymmetric error and uncertainty.] \textbf{Supported
  for error asymmetry.} Ordinal and cost-sensitive objectives improve both
  aggregate agreement and recall on Severe targets, on every seed, and form the
  most reliable single contributor to final performance. Selective prediction
  under calibrated uncertainty was not evaluated confirmatorily.
\end{description}

% ===========================================================================
\section{Thesis Organisation}
\label{sec:intro-organisation}

At the current protocol stage, the thesis is organised around the progression from
problem definition to evidence.

\textbf{Chapter 1 --- Introduction} defines the clinical and methodological
motivation, research problem, gaps, aim, objectives, research questions, hypotheses,
proposed contributions, significance and scope of the thesis.

\textbf{Chapter 2 --- Literature Review} examines the clinical grading systems that
define the reference labels, MRI sequence--pathology correspondence, localisation
and segmentation, backbone and contextual architectures, class imbalance and
ordinal grading, public datasets, external validation, interpretability and the
specific gaps that motivate the present study.

\textbf{Chapter 3 --- Methodology} translates those gaps into a protocol-grade
experimental design. It defines the public and local cohorts, leakage controls,
DICOM reconstruction, localisation and ROI construction, self-supervised
pretraining, adaptive routing, heterogeneous graph reasoning, grading objectives,
calibration, ablations, external validation, domain adaptation, statistical analysis
and reproducibility safeguards.

\textbf{Chapter 4 --- Results and Analysis} will report the staged E0--E8
experiments, including component-specific ablations, robustness analyses,
calibration, zero-shot external transfer, few-shot adaptation and failure analysis.
% Written after Chapter 4. Retained experiment groups; research questions unchanged.
The retained experiment groups are the eight-rung ablation ladder E0--E7 with two
matched controls on the graph configuration, a controlled input ablation testing
sequence dependence by intervention rather than by learned allocation, and a
Grad-CAM attribution analysis against an architecture-matched untrained floor.
The cross-institutional transfer group was prepared but not executed and is
reported as such. The research questions are unchanged from those fixed in
Section~\ref{sec:intro-rqs}.

\textbf{Chapter 5 --- Discussion and Conclusions} will interpret the findings
against the literature, distinguish demonstrated contributions from negative
results, state limitations, discuss implications for transport to other institutions
and identify the next studies required before any clinical-impact claim.

\toconfirm{final institutional thesis structure if the university requires results
and discussion as separate chapters, publication-based chapters, or an additional
conclusion chapter}. The scientific ordering above should remain recognisable even
if the administrative chapter numbering changes.

% ===========================================================================
\section{Rules for Completing Prospective Placeholders}
\label{sec:intro-placeholders}

Because this chapter is written before completion of implementation, several fields
are intentionally prospective. The following rules govern their later completion.

\begin{enumerate}
  \item \textbf{[TO CONFIRM]} is used only for institutional or governance facts
        that are currently unknown, such as an approval reference or permitted
        compute environment. These fields are replaced only with documented facts.

  \item \textbf{[TO RECORD AFTER IMPLEMENTATION]} is used for executed technical or
        cohort quantities that cannot responsibly be specified before the pipeline
        runs, such as final quality-controlled sample counts.

  \item \textbf{[LATER INTEGRATION]} is used for a small number of places where the
        final thesis may state the principal findings or update the novelty boundary.
        These additions may report evidence; they must not alter the prospective
        research questions or make an unsuccessful hypothesis disappear.

  \item Any substantive protocol change made after results become visible is
        documented with date, reason and affected analysis. Confirmatory and
        exploratory analyses remain distinguishable.

  \item The introduction will not be retrofitted to imply that the final retained
        architecture was inevitable. The final thesis should preserve the fact that
        ACSSL, adaptive routing, heterogeneous graph reasoning and supporting losses
        were hypotheses subjected to testing.
\end{enumerate}

These rules turn placeholders into a transparency mechanism rather than an invitation
to hindsight bias.

% ===========================================================================
\section{Chapter Summary}
\label{sec:intro-summary}

This chapter has defined the thesis as a study of structured and transferable lumbar
MRI grading rather than as a competition between image backbones. The clinical task
is inherently multi-level, multi-condition, bilateral and multi-sequence, while its
reference labels are ordinal and imperfectly reproducible. The technical literature
has established that localisation should precede grading, that sophisticated
multiview systems can perform strongly on public data, and that external performance
frequently degrades. Those achievements narrow rather than remove the research gap.

The thesis therefore asks whether three explicit structural mechanisms add value:
anatomically aligned cross-sequence self-supervision, disease-conditioned and
quality-aware sequence routing, and a typed heterogeneous graph over
level--condition--laterality targets. These mechanisms are evaluated component by
component against strong controls rather than bundled into a single architecture and
assumed beneficial. Supporting ordinal, cost-sensitive and uncertainty methods are
used to test the clinical shape of error without being advertised as standalone
novelty.

The translational component separates model development from external evaluation.
LumbarDISC supplies the public development benchmark; SPIDER may support anatomical
localisation; and the independent Rizgary Hospital cohort is reserved for zero-shot
central-canal transfer and controlled local adaptation. The local reference process
is treated explicitly as a potential source of shift rather than assumed equivalent
to the public benchmark. Clinical impact remains outside the scope until a separate
reader-assistance or prospective study is conducted.

The aim, objectives, RQ1--RQ5 and H1--H6 stated here form the prospective frame for
Chapters~\ref{ch:litreview} and~\ref{ch:methodology}. The remainder of the thesis is
therefore not required to prove that every proposed mechanism succeeds. It is
required to determine, with reproducible evidence, which assumptions survive
controlled testing and institutional transfer.

% Written at final submission.
In summary, the complete system improves on a single-sequence baseline
reproducibly and by a margin that survives correction for multiple comparisons
while remaining far too small to alter clinical practice,
and the objective function designed around clinically asymmetric error is the
most reliable single contributor to that improvement. Of the three proposed
structural mechanisms, relational typing is supported, anatomical topology and
disease-conditioned routing are not, and anatomically aligned self-supervision is
rejected on every axis on which it was tested; attribution analysis indicates
that the convolutional encoder already performs the localisation these mechanisms
were intended to supply. The principal limitations are that all results are
internal to one benchmark, that cross-institutional transfer was not executed,
and that the study is underpowered for single-component effects of the magnitude
observed.

% ===========================================================================
\printbibliography[heading=bibintoc,title={References}]

\end{document}
